\section*{Introduction}

Operads are mathematical structures that were first introduced as a way to formalize the 
notion of type of algebra: given a set of multilinear maps and relations between them,
the associated operad is composed of all the multilinear maps obtained by composing those in the
initial set.
These form in fact the generators of the operad. As detailed in \cite{MSJ02}, operad theory was first 
used in algebraic topology in the 1960s but had a 'renaissance' in the 1990s when it began to be used 
in many other fields, see for example \cite{LV12} for a very general algebraic approach to the theory. 
In particular, in combinatorics, operads provide the right framework to define the embedding of a 
combinatorial object in another \cite{CSLC, Gir18}. In this context, operads are defined by directly 
describing their elements and how to embed them. This is in contrast with the original way to define
them by generators and relations, and passing from one definition to the other is often a difficult
question and provides a lot of insight on the structure of the operad. 

A particularly successful example of such a transition is the pre-Lie operad which was first defined by 
being generated by the pre-Lie operator, i.e. an antisymmetric bilinear map $x\vartriangleleft y$ such that 
$(x\vartriangleleft y)\vartriangleleft z -x \vartriangleleft(y\vartriangleleft z) = (x\vartriangleleft z)\vartriangleleft y - x\vartriangleleft(z\vartriangleleft y)$.
Chapoton and Livernet later proved in \cite{CL01} that the elements of the pre-Lie operad could be seen 
as rooted trees and described the composition of two elements in a purely combinatorial way.

The pre-Lie operad along with the non-associative permutative operad~\cite{Liv06} are two examples
of interesting operads for which the combinatorial interpretation is given by trees. Even though 
interesting operads on graphs exist in the literature \cite{Kreimer:2000ja,Kon99,Wil05,MV19}, none 
of them is studied in a purely combinatorial framework. We propose a way to do so in this paper by 
extending the pre-Lie operad from trees to graphs.
Within our approach, we find that it is more natural to work on multigraphs and define an operad 
structure $\MG$ on the species of multigraphs, and an operad structure $\point{\MG}_{or}$ on the 
species of pointed oriented multigraphs. The operad $\MG$ in particular restricts to the 
Kontsevich-Willwacher operad~\cite{MV19} on the species of graphs $\G$. We show that these operads 
and the pre-Lie operad $\PLie$ relate to each other by the following commutative diagram:
\begin{equation*}
    \begin{tikzcd}
        \T \arrow[r, "\sim"] \arrow[d, hook]
        & \PLie\cap\K\Operad/\mathcal{I} \arrow[r, equal] \arrow[d, hook]
        & \PLie\cap\Operad \arrow[d,hook] \arrow[r, hook]
        & \PLie \arrow[d,hook]\\
        \G_c \arrow[r, "\sim"] \arrow[d, hook]
        & \Operad\cap\point{\G}_{orc}/\mathcal{I}\cap\point{\G_{orc}} \arrow[d, hook]
        & \point{\G}_{orc}\cap\Operad \arrow[l, two heads] \arrow[r, hook] \arrow[d, hook]
        & \point{\G}_{orc}\cap\ST \arrow[d, hook] \\
        \MG_c \arrow[r, "\sim"]
        & \Operad/\mathcal{I}
        & \Operad \arrow[l, two heads] \arrow[r, hook]
        & \MG\times\PLie
    \end{tikzcd},
\end{equation*}
where $\T$, $\G_c$ and $\MG_c$ are respectively the sub-operads of trees, connected graphs and 
connected multigraphs of $\MG$, $\point{\G_{orc}}$ is the sub-operad of pointed connected oriented 
graphs of $\point{\MG}_{or}$, and $\ST$ is an operad on spanning trees with $\Operad$ and 
$\mathcal{I}$ sub-species of $\ST$ generated by ad hoc sums of spanning trees. 
We show that while, contrarily to $\PLie$, the operad $\MG$ does not admit a simple presentation by 
generators and relations, its sub-operads generated by combinations of the empty graph over two vertices 
and the segment graph, have interesting links with the commutative operad $\Com$ and the magmatic 
operad $\ComMag$. In particular, one of them is Koszul, which allows us to exhibit its Koszul dual.
We end our study by providing new constructions on species and operads. These constructions enable us 
to define a general way of constructing operads on graphs, which we use to justify the operad structure 
of $\MG$ and $\point{\MG}_{or}$. More specifically, we give a general way to construct operads on 
multigraphs where the partial composition of two elements $g_1\comp g_2$ can be informally described 
by the following steps:
\begin{enumerate}
    \item take the union of $g_1$ and $g_2$,
    \item remove the vertex $\ast$ from the union, we now have loose ends,
    \item \underline{independently} connect each loose end to vertices of $g_2$,
\end{enumerate}
where by independently we mean that the way connect a loose end does not depend on the 
way we connect the others. A formal definition of what we mean loose end and connecting them
is given in section~\ref{intro_graphs}.

This paper is organized as follows. In Section~\ref{sec_intro} we give some general definitions of
species theory and of operad theory, as well as some general definitions on graphs and multigraphs. 
In Section~\ref{sec_extending_plie} we define the aforementioned two operads on multigraphs and 
pointed oriented multigraphs and we exhibit the correspondence between these operads and the pre-Lie 
and Kontsevich-Willwacher operads. We then proceed to study the operad on multigraphs and some of its 
sub-operads. Finally, in Section~\ref{sec_graph_operads} we exhibit our proposal of constructing species 
and operads and we apply them in order to obtain our general construction of operads on multigraphs.

This paper is an extended version of the extended abstract for FPSAC 2020~\cite{AGKT20} with proofs and 
additional results.





\section{Context}\label{sec_intro}
In this entire paper, unless otherwise stated, $V$ denotes a finite set and $V_1$ and $V_2$ two disjoint 
sets such that $V=V_1\cup V_2$. The letter $n$ always denotes a non negative integer and we denote by $[n]$ 
the set $\set{1,\dots, n}$. All vector spaces appearing in this paper are defined over a field of 
characteristic 0 denoted by $\K$. Finally, for any set $E$, $\K E$ denotes the free vector space over $E$.





\subsection{Graphs and multigraphs}\label{intro_graphs}

\textbf{Multisets.}
Because multiset notations can be cumbersome, it is customary to consider them in the same way as sets and we do so whenever it
is possible. However, multiset are necessary in order to formally define multigraphs and we sometimes need notations proper to them.

A {\em multiset} over $V$ is subset of $V\times\N^*$ such that $\forall (v,i)\in V\times\N\setminus\set{0,1},\, (v,i)\in m \Rightarrow (v,i-1)\in m$. Let $m$ be a multiset 
over $V$. We say that an element $v\in V$ {\em is in} $m$ and denote by $v\in m$ if $(v,1)\in m$. In this case, we call {\em multiplicity of $v$ in $m$} 
and denote by $m(v)$ the greatest integer $i\in\N$ such that $(v,i)\in m$. If $v\not\in m$ we set $m(v)=0$. Let $m'$ be a multiset over a set $V'$. The
{\em deletion} of $m'$ from $m$ is the multiset $m\setminus m' = \set{(v,i)\in m\,|\, i\leq m(v)-m'(v)}$ and the {\em disjoint union} of $m$ and
$m'$ is the multiset $m\sqcup m'= m\cup\set{(v,i+m(v))\,|\,(v,i)\in m'}$.

The {\em size} of a multiset is its cardinality as a set and we call {\em unordered multipair} a multiset of size 2. 
We denote by $\m(V)$ the set of multisets over $V$. 
Every set can be seen as multiset by the bijection $W\cong W\times\set{1}$.

\noindent\textbf{Graphs and multigraphs.}
A {\em graph} or {\em simple graph} over a non empty set $V$ is a set of unordered pairs of distinct elements of $V$ and a 
{\em multigraph} is a multiset of unordered multipairs over $V$. In this context, the elements of a 
(multi)graph $g$ are called {\em edges} of $g$, the elements of $V$ are called {\em vertices} or {\em nodes} of $g$ and the pairs
$(e,v)$ with $e\in g$ and $v\in e$ are called {\em ends} of $g$. For $(e,v)$ an end of $g$, we say that $(e,v)$ is both
an {\em end of the edge} $e$ and an {\em end of the vertex} $v$. For $v$ a vertex, an edge of the form $\set{v,v}$ 
is called a {\em loop over $v$} and we call {\em neighbour of $v$} any vertex $u\not = v$ such that $\set{u,v}\in g$ 
($v$ can not be its own neighbour). For $W\subseteq V$, we call the {\em restriction of $g$ to $W$} and denote by $g|_W$,
the multigraph over $W$ with edges exactly the edges of $g$ contained in $W$. 
Graphs can be seen as multigraphs with at most one edge between two vertices and no loops. 

Since we describe different partial composition on multigraphs in term of removing vertices and reconnecting loose ends, we must
formally introduce the notion of a multigraph with loose ends. However, do note that the elements of study of this paper are still multigraphs
as defined above. Indeed, loose ends will only appears as a step in partial compositions, in order to provide a combinatorial understanding,
and never outside it. 
A {\em multigraph with loose ends} over $V$ is a multigraph over $V\cup\set{\emptyset}$ and we call {\em loose ends} the ends of 
$\emptyset$. We identify multigraphs with multigraphs with loose ends with no loose ends. Let $g$ be a
multigraph with loose ends over $V$. We define the operations of removing a vertex and connecting a loose end as follow.
\begin{itemize}
    \item For $v\in V$, the multigraph with loose ends obtained by {\em removing} $v$ of $g$ is the multigraph with
    loose ends over $V\setminus\set{v}$ obtained by replacing every occurrence of $v$ by $\emptyset$ in $g$:
    $\set{\sigma(e)\,|\, e\in g}$ where $\sigma(\set{v_1,v_2}) = \set{\sigma(v_1),\sigma(v_2)}$ and 
    $\sigma:V\cup\set{\emptyset}\to (V\setminus \set{v})\cup\set{\emptyset}$ sends $v$ to $\emptyset$ and is the identity elsewhere. 
    \item For $\set{u,\emptyset}$ an edge with a loose end of $g$ and $v_1,\dots, v_n \in V$, the multigraph with loose ends obtained 
    by {\em connecting $(\set{u,\emptyset},\emptyset)$ to $(v_1,\dots, v_n)$} is the multigraph with loose ends obtained by replacing the edge 
    $\set{u,\emptyset}$ with the edges $\set{u,v_1},\dots,\set{u,v_n}$: $\left(g\setminus\set{\set{u,\emptyset}}\right)\sqcup\set{\set{u,v_i}}_{1\leq i\leq n}$. If $n=1$ we write connecting to $v_1$ instead of connecting to $(v_1)$.
\end{itemize}

\begin{rema}
    An edge $\set{\emptyset, \emptyset}$ have two loose ends $(\set{\emptyset, \emptyset}, \emptyset)$. We could first connect one of them
    to $(u_1,\dots, u_n)$ to obtain the new edges $\set{\emptyset, u_1}, \dots, \set{\emptyset, u_n}$, and then respectively connect their associated
    loose ends to $(v_{1,1}, \dots, v_{1,k_1}), \dots, (v_{n,1}, \dots, v_{n,k_n})$ to obtain the edges 
    $\set{u_1, v_{1,1}}, \dots, \set{u_1, v_{1,k_1}}, \dots, \set{u_n, v_{n,k_n}}$. In practice we will always choose $k_i = k$ for all $1\leq i \leq n$ 
    and $v_{i,j} = v_j$ for all $1\leq i \leq n$, $1\leq j \leq k$ where $k$ is an integer and $v_1,\dots, v_n$ are vertices. 
    In these circumstances, connecting one of the loose ends of an edge $\set{\emptyset, \emptyset}$ to $(u_1,\dots, u_n)$ then connecting the 
    loose ends of all the new edges to $(v_1,\dots, v_k)$ returns the same edges than to first connect one of the loose ends of the edge to 
    $(v_1,\dots, v_k)$ then connecting the loose ends of all the new edges to $(u_1,\dots, u_n)$. To simplify, we hence call these two operation with 
    same results as {\em connecting the ends of $\set{\emptyset, \emptyset}$ to $(u_1,\dots, u_n)$ and $(v_1,\dots, v_k)$}.
\end{rema}

\begin{exam}\label{ex_graph}
    Let $g$ be the multigraph over $\set{a,b,c,d,e}$ with edges $\set{a,a}$, $\set{a,b}$, $\set{a,d}$, $\set{b,e}$, 
    $\set{c,e}$, $\set{c,d}$ and $\set{c,d}$. The multigraph with loose ends obtained by removing $a$ has edges 
    $\set{\emptyset, \emptyset}$, $\set{\emptyset,b}$, $\set{\emptyset,d}$, $\set{b,e}$, $\set{c,e}$, $\set{c,d}$ and $\set{c,d}$.
    Connecting the two loose ends of $\set{\emptyset, \emptyset}$ to $b$ and $(b,c)$, and the loose ends of 
    $\set{\emptyset,b}$, $\set{\emptyset,d}$ respectively to $d$ and $e$ gives us the multigraph with edges
    $\set{b,b}$, $\set{b,c}$, $\set{b,d}$, $\set{d,e}$, $\set{b,e}$, $\set{c,e}$, $\set{c,d}$, and $\set{c,d}$.

    We usually represent multigraphs with blue circles for vertices and full blue lines for edges. For example, we represent 
    $g$, the multigraph with loose end obtained by removing $a$ and the multigraph obtained by reconnecting the loose ends
    as above:
    \begin{equation}
        \begin{tikzpicture}[Centering,scale=1]
            \tikzset{every loop/.style={}}
            \node[NodeGraph, minimum size = 5mm](d)at(0,0){$d$};
            \node[NodeGraph, minimum size = 5mm](f)at(1,0){$c$};
            \node[NodeGraph, minimum size = 5mm](e)at(1,1.5){$e$};
            \node[NodeGraph, minimum size = 5mm](b)at(0,1.5){$b$};
            \node[NodeGraph, minimum size = 5mm](a)at(-1,0.75){$a$};
            \draw[EdgeGraph](b)--(a);
            \draw[EdgeGraph](d)--(a);
            \draw[EdgeGraph](d)edge[bend right=40](f);
            \draw[EdgeGraph](d)edge[bend left=40](f);
            \draw[EdgeGraph](f)--(e);
            \draw[EdgeGraph](b)--(e);
            \draw[EdgeGraph](a)edge[loop above](a);
        \end{tikzpicture}
        \quad\quad\quad\quad\quad
        \begin{tikzpicture}[Centering,scale=1]
            \tikzset{every loop/.style={}}
            \node[NodeGraph, minimum size = 5mm](d)at(0,0){$d$};
            \node[NodeGraph, minimum size = 5mm](f)at(1,0){$c$};
            \node[NodeGraph, minimum size = 5mm](e)at(1,1.5){$e$};
            \node[NodeGraph, minimum size = 5mm](b)at(0,1.5){$b$};
            \node[minimum size = 5mm](a)at(-1,0.75){};
            \draw[EdgeGraph](b)--(a);
            \draw[EdgeGraph](d)--(a);
            \draw[EdgeGraph](d)edge[bend right=40](f);
            \draw[EdgeGraph](d)edge[bend left=40](f);
            \draw[EdgeGraph](f)--(e);
            \draw[EdgeGraph](b)--(e);
            \draw[EdgeGraph](a)edge[loop above](a);
        \end{tikzpicture}
        \quad\quad\quad\quad\quad
        \begin{tikzpicture}[Centering,scale=1]
            \tikzset{every loop/.style={}}
            \node[NodeGraph, minimum size = 5mm](d)at(0,0){$d$};
            \node[NodeGraph, minimum size = 5mm](f)at(1,0){$c$};
            \node[NodeGraph, minimum size = 5mm](e)at(1,1.5){$e$};
            \node[NodeGraph, minimum size = 5mm](b)at(0,1.5){$b$};
            \draw[EdgeGraph](b)--(d);
            \draw[EdgeGraph](b)--(f);
            \draw[EdgeGraph](d)--(e);
            \draw[EdgeGraph](d)edge[bend right=40](f);
            \draw[EdgeGraph](d)edge[bend left=40](f);
            \draw[EdgeGraph](f)--(e);
            \draw[EdgeGraph](b)--(e);
            \draw[EdgeGraph](b)edge[loop above](b);
        \end{tikzpicture}.
    \end{equation}
\end{exam}

\noindent\textbf{Trees, spanning trees and rooted trees.}
A {\em path} in a multigraph $g$ is a sequence $v_1,\dots, v_n$ of vertices such that $n>1$ and $\set{v_i,v_{i+1}}$ is an edge of $g$
for $1\leq i < n$. A multigraph $g$ is {\em connected} if for any
two vertices $u$ and $v$ of $g$, the exists a path $v_1,\dots,v_n$ such that $v_1=u$ and $v_n=v$. In a connected graph,
the {\em distance} between two vertices is the length of the shortest path between these vertices. 
A {\em cycle} is a path such that $v_1=v_n$ and a {\em tree} or {\em abstract tree} is a graph without cycle. In a tree, a vertex is
said to be {\em internal} if it has at least two neighbours and said to be a {\em leaf} else. For $g$ a connected multigraph
over $V$, a {\em spanning tree} of $g$ is a tree over $V$ with edges contained in $g$.

A {\em rooted tree} is a pair $(t,r)$ with $t$ a tree and $r$ a vertex of $t$ which we call {\em root} or {\em root vertex} of $t$.
For $(t,r)$ a rooted tree and $v$ a vertex of $t$, we call {\em children of $v$} the neighbours of $v$ farther 
from $r$ than $v$. If $v\not = r$ we also call {\em parent of $v$} the neighbour of $v$ closer to $r$ than 
$v$. For $p$ the parent of $v$ and $c$ a child of $v$, the end $(\set{p,v},v)$ is the {\em parent end} of $v$ and the end
$(\set{v,c},v)$ is a child end of $v$. A {\em corolla} is a rooted tree where only the root has children.


\begin{exam}\ 
    \begin{itemize}
        \item The multigraph $g$ of Example~\ref{ex_graph} is connected but $g|_{\set{a,b,c}}$ is not.
        \item We represent spanning trees by using red lines for the edges that are part of the tree. Here is a representation
        of $g$ with the spanning tree $\set{\set{a,b},\set{a,d},\set{b,e},\set{d,c}}$:
        \begin{equation}
            \begin{tikzpicture}[Centering,scale=1]
                \tikzset{every loop/.style={}}
                \node[NodeGraph, minimum size = 5mm](d)at(0,0){$d$};
                \node[NodeGraph, minimum size = 5mm](f)at(1,0){$c$};
                \node[NodeGraph, minimum size = 5mm](e)at(1,1.5){$e$};
                \node[NodeGraph, minimum size = 5mm](b)at(0,1.5){$b$};
                \node[NodeGraph, minimum size = 5mm](a)at(-1,0.75){$a$};
                \draw[EdgeGraph, DarkRed](b)--(a);
                \draw[EdgeGraph, DarkRed](d)--(a);
                \draw[EdgeGraph, DarkRed](d)edge[bend right=40](f);
                \draw[EdgeGraph](d)edge[bend left=40](f);
                \draw[EdgeGraph](f)--(e);
                \draw[EdgeGraph, DarkRed](b)--(e);
                \draw[EdgeGraph](a)edge[loop above](a);
            \end{tikzpicture}.
        \end{equation}
        \item We generally represent rooted trees as we do multigraphs except with a square for the root vertex. 
        For example, the rooted tree $\left(\set{\set{a,c},\set{c,b}},c\right)$ over $\set{a,b,c}$ is represented as follows:
        \begin{equation}
            \begin{tikzpicture}[Centering,scale=.7]
                \node[RootGraph](s)at(0,0){$r$};
                \node[NodeGraph](a)at(1,1){$a$};
                \node[NodeGraph](b)at(1,-1){$b$};
                \draw[EdgeGraph](a)--(s);
                \draw[EdgeGraph](s)--(b);
            \end{tikzpicture}
            \enspace = \enspace\left(
            \begin{tikzpicture}[Centering,scale=.7]
                \node[NodeGraph](s)at(0,0){$r$};
                \node[NodeGraph](a)at(1,1){$a$};
                \node[NodeGraph](b)at(1,-1){$b$};
                \draw[EdgeGraph](a)--(s);
                \draw[EdgeGraph](s)--(b);
            \end{tikzpicture}
            \enspace , \enspace r\right).
        \end{equation}
        Sometimes -in particular in the case of Schröder trees- rooted trees may stand for other objects. In these cases, while the objects for which they
        stand for may be labelled with $V$, the rooted trees will only have their leaves labelled with $V$, e.g. Example~\ref{ex_stree}. Hence, in order 
        to avoid some confusion, we use a different representation for rooted trees in these cases, with the root being the bottom-most vertex and the
        vertices and edges not sharing the previously set color code.
        Furthermore, as we explain in Section~\ref{sec_operads}, while there exists different kind
        of operad structures on rooted trees, Schröder trees have an implicit operad structure given by the grafting of trees; hence we found
        important to clearly differentiate these Schröder trees, in particular since we can consider Schröder trees enriched with the species
        of rooted trees.
    \end{itemize}
\end{exam}

\noindent\textbf{Orientations.}
Let $g$ be a multigraph over $V$. An {\em orientation} of $g$ is a map $\ori$ from $g$ to $\m(V)\times\m(V)$ such that for any edge $e$ of $g$,
$e= \pi_1(\ori(e))\sqcup \pi_2(\ori(e))$, with $\pi_i$ the projection on the $i$-th coordinate. We respectively call $\pi_1(\ori(e))$ and $\pi_2(\ori(e))$ 
the {\em sources} and {\em targets} of $e$. For $(e,v)$ an end of $g$, we say that it is a {\em source end} if $v$ is a source of $e$ and a {\em target end}
else. 
An {\em oriented multigraph} is a pair of a multigraph and an orientation of this multigraph. 
The union of two oriented multigraphs over disjoint sets $(g_1,\ori_1)$ and $(g_2,\ori_2)$ is the oriented multigraph $(g_1\sqcup g_2,\ori)$, where
$\ori(e) = \ori_1(e)$ if $e\in g_1$ and $\ori(e)=\ori_2(e)$ else.

The notion of orientation directly generalizes to oriented multigraphs by considering them as multigraphs over $V\cup\set{\emptyset}$ and an 
{\em oriented multigraph with loose ends} is then a pair of a multigraph with loose ends and an orientation of it. Given an oriented multigraph
with loose ends $(g,\ori)$, {\em removing} a vertex $v$ from $g$ induces the orientation $\ori'$ such that 
$\ori'\circ\sigma(e)= (\sigma\times\sigma)\circ\ori(e)$ with $\sigma$ defined as before. {\em Connecting} the end $(\set{u,\emptyset},\emptyset)$ to
$v_1,\dots, v_n$ induces the orientation $\ori'$ where the sources and targets of $\set{u,v_i}$ are obtained by replacing $\emptyset$ by $v_i$ in
$\ori(\set{u,\emptyset})$.

\begin{rema}
    We do not use the standard definition of an orientation: since the source and target maps have their image in $\m(V)$, an edge can have two target 
    ends and no source ends and vice versa.
\end{rema}

\begin{exam}
    Let $g$ be the multigraph of Example~\ref{ex_graph}. We define the following orientation of $g$: $\ori(\set{a,a})=(\emptyset,\set{a,a})$, 
    $\ori(\set{a,b})=(\set{b},\set{a})$, $\ori(\set{a,d})=(\set{a},\set{d})$, $\ori(\set{b,e})=(\emptyset,\set{b,e})$, 
    $\ori(\set{c,e})=(\set{c,e},\emptyset)$, and $\ori(\set{c,d})=(\set{c},\set{d})$ and $\ori(\set{c,d})=(\set{d},\set{c})$. We represent
    oriented multigraphs by adding arrow heads to the targets:
    \begin{equation}
        (g,\ori)=\begin{tikzpicture}[Centering,scale=1]
            \tikzset{every loop/.style={}}
            \node[NodeGraph, minimum size = 5mm](d)at(0,0){$d$};
            \node[NodeGraph, minimum size = 5mm](f)at(1,0){$c$};
            \node[NodeGraph, minimum size = 5mm](e)at(1,1.5){$e$};
            \node[NodeGraph, minimum size = 5mm](b)at(0,1.5){$b$};
            \node[NodeGraph, minimum size = 5mm](a)at(-1,0.75){$a$};
            \draw[EdgeGraph, -Stealth](b)--(a);
            \draw[EdgeGraph, Stealth-](d)--(a);
            \draw[EdgeGraph, -Stealth](d)edge[bend right=40](f);
            \draw[EdgeGraph, Stealth-](d)edge[bend left=40](f);
            \draw[EdgeGraph](f)--(e);
            \draw[EdgeGraph, Stealth-Stealth](b)--(e);
            \draw[EdgeGraph, -Stealth](a)edge[loop above](a);
        \end{tikzpicture}.
    \end{equation}
\end{exam}



\subsection{Definitions and background on species} \label{sec_species}
We recall here basic definitions on species. We refer the interested reader to~\cite{BLL98} for a detailed
presentation of combinatorial species.

\begin{defi}
    A {\em linear species} $\spe{S}$ consists of the following data:
    \begin{itemize}
       \item For each finite set $V$, a vector space $\spe{S}[V]$ of finite dimension,
       \item For each bijection of finite sets $\sigma: V\to V'$, a linear map 
       $\spe{S}[\sigma]:\spe{S}[V]\to \spe{S}[V']$.
       These maps should be such that 
       $\spe{S}[\sigma_1\circ\sigma_2] = \spe{S}[\sigma_1]\circ \spe{S}[\sigma_2]$ 
       and $\spe{S}[\id] = \id$.
    \end{itemize}
    Furthermore if $\spe{S}[\emptyset]  = \set{0}$, then $\spe{S}$ is said to be {\em positive}.
    
    The elements of the vector spaces $\spe{S}[V]$ are called {\em elements of $\spe{S}$}.
\end{defi}

We will use the term {\em species} to refer to linear species.

When defining a species, the maps $\spe{S}[\sigma]$ are often clear from the context and we do not mention them.
Let $\spe{S}$ be a species. A {\em sub-species} of $\spe{S}$ is a species $\spe{R}$ such that $\spe{R}[V]$ is 
a sub-space of $\spe{S}[V]$ for every finite set $V$ and $\spe{R}[\sigma] = \spe{S}[\sigma]$ for every
bijection of finite sets $\sigma$. 
For $E$ a set of elements of $\spe{S}$, the {\em species generated by $E$} is the smallest sub-species of
$\spe{S}$ containing $E$.
We denote by $\spe{S}_{n}$ the sub-species of $\spe{S}$ defined by $\spe{S}_{n}[V]=\spe{S}[V]$ if 
$\card{V} = n$ and $\spe{S}_n[V]=\set{0}$ otherwise, and by $\spe{S}_{n+}$ the sub-species of $\spe{S}$
defined by $\spe{S}_{n+}[V]=\spe{S}[V]$ if $\card{V}\geq n$ and $\spe{S}_{n+}[V]=\set{0}$ otherwise.
We also denote by $\spe{S}_+$ the species $\spe{S}_{1^+}$. The {\em Hilbert series} of $\spe{S}$ is the 
formal power series defined by 
$\hilbert{\spe{S}}(x)=\sum_{n\geq 0}\frac{\dim\spe{S}[[n]]}{n!}x^n$.

A {\em morphism of species} from a species $\spe{R}$ to $\spe{S}$ is a collection of linear maps 
$f_V : \spe{R}[V] \to \spe{S}[V]$ such that for each bijection $\sigma: V\to V'$, we have 
$f_{V'}\circ \spe{R}[\sigma] = \spe{S}[\sigma]\circ f_V$. For easier reading, we will often forget 
the index $V$. 



\begin{exam}\label{ex_species}\
    \begin{itemize}
        \item The {\em Identity}, {\em exponential} and {\em singleton species} are respectively defined 
        by $\spe{\id}[V] = \K V$, $\Sym[V] = \K\set{V}$ and $\spe{X} = \Sym_1$.

        \item For $V$ a finite set, we denote by $\spe{Pol}[V]$ the \textit{set} (not the vector space) of 
        polynomials with coefficients in $\N$, variables in $V$ and null constant coefficient. To consider 
        the species $\K\spe{Pol}$, we must take into consideration the fact that we need to differentiate 
        the plus of polynomials and the addition of vectors. We will thus denote by $\oplus$ the former and
        keep $+$ for the latter and we will denote by $0_V\in\spe{Pol}[V]$ the polynomial constant to 0 and
        keep the notation $0$ for the null vector. For example, $ab\oplus c$ is an element of $\spe{Pol}
        [\{a,b,c\}]$, but $a\oplus b+c$ is a vector in $\K \spe{Pol}[\{a,b,c\}]$.

        \item We have a natural morphism from $\spe{\id}$ to $\K\spe{Pol}$ given by $v\mapsto v$ and three
        natural morphisms from $\Sym$ to $\K\spe{Pol}$ respectively given by $V\mapsto \prod_{v\in V}v$,
        $V\mapsto\bigoplus_{v\in V}v$ and $V\mapsto\sum_{v\in V}v$. In particular, the image of $\spe{\id}$ 
        by this morphism is the species $\K\spe{Pol}^1$ of homogeneous polynomials of degree 1.

        \item We respectively denote by $\T$, $\G$ and $\MG$ the species of trees, graphs and multigraphs.
        We also respectively denote by $\G_c$ and $\MG_c$ the species of connected graphs and connected
        multigraphs. These five species are related by the following diagram:
        \begin{equation}\label{graph_diagram}
            \begin{tikzcd}[row sep=scriptsize,column sep=scriptsize]
            \T \arrow[r, hook] & \G_c \arrow[r, hook] \arrow[d, hook] & \MG_c \arrow[d, hook] \\
            & \G \arrow[r, hook] & \MG & 
            \end{tikzcd}.
        \end{equation}
        We add the index $or$ to these five species to designate their oriented counterpart, e.g.
        $\G_{orc}$ is the species of oriented connected graphs. Notice that since we only defined
        multigraphs over non empty finite sets, these species are positive, i.e. we set $\MG[\emptyset] = \set{0}$.


        \item There is an monomorphism $\Phi$ from the species $\MG$ to $\K\spe{Pol}$ defined as follows:
        \begin{itemize}\label{graphpol}
            \item the empty graph $\emptyset_V\in \MG[V]$ is sent on the null polynomial $\Phi(\emptyset_V) = 0_V$;
            \item an edge $e=\set{v_1,v_2}$ is sent on the monomial $\Phi(e) = v_1v_2$;
            \item an element $g\in \MG[V]$ is sent on the polynomial $\Phi(g) = \bigoplus_{e\in g} \Phi(e)$.
        \end{itemize}
        This enables us to see the species of multigraphs as the sub-species $\K\spe{Pol}^2$ of homogeneous 
        polynomials of degree 2. This will be useful in the following sections to do computations on multigraphs since 
        it is easier to formally write operations, and in particular composition, on polynomials than on multigraphs.
        With this identification, the multigraph of Example~\ref{ex_graph} writes as the polynomial 
        $\Phi(g) = a^2\oplus ab\oplus ad\oplus be\oplus ec\oplus 2dc$.
    \end{itemize}
\end{exam}

One strong point of species is the existence of  different operations on them which enable to construct
new species from existing ones. Let $\spe{R}$ and $\spe{S}$ be two species. We can then construct new 
species which are defined as follows: {\em sum} $(\spe{R}+\spe{S})[V] = \spe{R}[V]\bigoplus \spe{S}[V]$, 
{\em product} $\spe{R}\cdot \spe{S}[V] = \bigoplus_{V_1\cup V_2 = V} \spe{R}[V_1]\otimes \spe{S}[V_2]$,
\linebreak[9] {\em Hadamard product} $(\spe{R}\times \spe{S})[V] = \spe{R}[V]\otimes \spe{S}[V]$, 
{\em derivative} $\spe{S}'[V] = \spe{S}[V\cup\set{\ast}]$ (where $\ast\not\in V$), 
{\em second derivative} $\spe{S}''[V] = \spe{S}[V\cup\set{\ast_1,\ast_2}]$ ($\ast_1,\ast_2\not\in V$) 
and {\em pointing} $\point{\spe{S}}[V] = \spe{S}[V]\otimes \K V$. Furthermore if $\spe{S}$ is positive we 
can also define the {\em composition} of $\spe{R}$ and $\spe{S}$ by
$$\spe{R}(\spe{S})[V] = \bigoplus_{\text{$P$ partition of $V$}}\spe{R}[P]\bigotimes_{P_i\in P} \spe{S}[P_i],$$
where $\bigotimes_{P_i\in P} \spe{S}[P_i]=\left(\spe{S}[P_1]\otimes\dots\otimes\spe{S}(P_k)\right)_{\symgrp_k}$ 
should be seen as an unordered tensor product. We call {\em assemblies} of $\spe{S}$ the elements of
this unordered tensor product.

\begin{rema}\label{remark_operations}\
    \begin{itemize}
        \item The product of species is associative: $(\spe{S}_1\cdot\spe{S}_2)\cdot\spe{S}_3 = \spe{S}_1\cdot(\spe{S}_2\cdot\spe{S}_3)$.
        We will forget the parenthesis when considering the product of three or more species and directly denote $\spe{S}_1\cdot\spe{S}_2\cdot\spe{S}_3$.
        \item These operations are compatible with Hilbert series: the Hilbert series of a sum of species is the sum of their respective
        Hilbert series, etc.
        \item When considering the product of two derivatives $\spe{R}'\cdot\spe{S}'$, to avoid confusion we will use 
        the notations $\spe{R}'[V] = \spe{R}[V\cup\set{\ast_1}]$ and $\spe{S}'[V] = \spe{S}[V\cup\set{\ast_2}]$.
    \end{itemize}
\end{rema}

\begin{exam}\label{example_derivation}\
    \begin{itemize}
        \item The species $\point{\T}$ is the species of rooted trees. In order to stay coherent with our representation
        of rooted trees, we will also represent the pointed vertex of the elements $\point{\MG}$ and its sub-species
        with a square.
        \item Since $\Sym[V]$ is always reduced to a space of dimension 1, the species $\Sym(\spe{S})$ can 
        be interpreted as the species of assemblies of $\spe{S}$.
    \end{itemize}
\end{exam}

We can construct a particularly interesting species by using the composition of positive species. For 
$\spe{S}$ a positive species such that $\spe{S} = \spe{X} + \spe{S}_{2^+}$, the species $\stree{\spe{S}}$ 
of {\em Schröder trees enriched with $\spe{S}$} is defined as the species satisfying the equation
\begin{math}
    \stree{\spe{S}} = \spe{X} + \spe{S}_{2^+}(\stree{\spe{S}}).
\end{math}


The vector space $\stree{\spe{S}}[V]$ is then generated by abstract rooted trees with internal nodes labelled
by elements of $\spe{S}$ and set of leaves $V$. In particular, remark that the internal nodes have a linear
behaviour: if $W$ is a subset of $V$, $x$ and $y$ are two elements of $\spe{S}[W]$ and $k$ is an element of
$\K$, a tree $t$ with a node labelled by $x+ky$ is equal to the sum $t_1+kt_2$ where $t_1$ and $t_2$ are
identical to $t$ except at the node which was labelled by $x+ky$ in $t$ which is now respectively labelled by 
$x$ and $y$. In this context, the elements of $\spe{S}$ are identified with the corollas of $\stree{\spe{S}}$.

\begin{rema}
    Contrarily to the usual case, when working in the context of species the Schröder trees we consider are not 
    planar. This is because we can already differentiate leaves through their label and hence do not need an
    order on them.
\end{rema}

\begin{exam}\label{ex_stree}
    Here is an example of an element in $\stree{\spe{S}}[[4]]$ where $\spe{S}[V] = \K{x}$ if $\card{V} = 3$,
    $\spe{S}[V] = \K{y}$ if $\card{V}=2$ and $\spe{S}[V] = \set{0}$ otherwise.
    \begin{equation}
        \begin{tikzpicture}[Centering, scale=0.6]
            \node[NodeFree, minimum size=8mm](r)at(0,0){$x_{12,3,4}$};
            \node[NodeFree, minimum size=8mm](n)at(-1.5,2.75){$y_{1,2}$};
            \node[NodeFree, minimum size=8mm](s)at(0.25,2.5){$4$};
            \node[NodeFree, minimum size=8mm](3)at(1.75,1.75){$3$};
            \node[NodeFree, minimum size=8mm](1)at(-3.5,5){$1$};
            \node[NodeFree, minimum size=8mm](2)at(-1.5,5.5){$2$};
            \draw[EdgeFree](r)--(n);
            \draw[EdgeFree](r)--(s);
            \draw[EdgeFree](r)--(3);
            \draw[EdgeFree](n)--(1);
            \draw[EdgeFree](n)--(2);
        \end{tikzpicture}
    \end{equation}
\end{exam}


\subsection{Definitions and background on operads}\label{sec_operads}
We give here basic definitions and results of the theory as well as some classical examples. We use the definition
by partial composition of an operad (see e.g.~\cite{F14}).
We refer the reader to~\cite{Men15} and \cite{LV12} for a more general approach to the theory of operads.

A {\em partial composition} of a species $\spe{S}$ is a collection of linear maps 
$\comp[v]: \spe{S}[V_1]\otimes\spe{S}[V_2]\to \spe{S}[V_1+V_2-\set{v}]$, for $V_1$ and $V_2$ disjoint finite
sets and $v\in V_1$. These maps must be natural: for $\sigma_1$ and $\sigma_2$ bijections with respective 
domain $V_1$ and $V_2$, we have 
$\comp[\sigma_1(v)]\circ\spe{S}[\sigma_1]\otimes\spe{S}[\sigma_2] = \spe{S}[\sigma_1\circ\sigma_2]\circ\comp[v]$. 
By naturality, specifying a partial composition amounts to specifying a morphism 
$\comp:\spe{S}'\cdot\spe{S}\to\spe{S}$ from which we can recover all the maps of the collection.

\begin{defi}
    A {\em symmetric linear operad} is a positive linear species $\Operad$ equipped with an {\em unity} 
    $e: \spe{X}\to \Operad$ and a  {\em partial composition} $\comp:\Operad'\cdot\Operad\to \Operad$ such
    that the following diagrams commute
    \begin{center}\label{def_op}
        \begin{tikzcd}
            \Operad''\cdot\Operad\cdot\Operad \arrow[r, "{\comp[\ast_1]}"] \arrow[d, "{\comp[\ast_2]}\circ\id\cdot\twist"] & \Operad'\cdot\Operad \arrow[d, "{\comp[\ast_2]}"] & &
            \Operad'\cdot \Operad'\cdot\Operad \arrow[r, "{\comp[\ast_1]}\cdot\id"] \arrow[d, "\id\cdot{\comp[\ast_2]}"] & \Operad'\cdot\Operad \arrow[d, "{\comp[\ast_2]}"] \\
            \Operad'\cdot \Operad \arrow[r, "{\comp[\ast_1]}"] & \Operad & & 
            \Operad'\cdot \Operad \arrow[r, "{\comp[\ast_1]}"] & \Operad
        \end{tikzcd}
    \end{center}
    \begin{center}
        \begin{tikzcd}
            \Operad'\cdot \K X \arrow[r, "\Operad'\cdot e"] \arrow[rd, "p"] & \Operad'\cdot\Operad \arrow[d, "\comp"] & \K X'\cdot \Operad \arrow[l, "e'\cdot \Operad" swap]\arrow[dl, "\cong" swap]\\
            & \Operad 
        \end{tikzcd}
    \end{center}
    where  $\tau:x\otimes y\mapsto y\otimes x$ and $p_V: x \mapsto \Operad[\sigma](x)$ with $\sigma$ the
    bijection that sends $\ast$ on $v$ and is the identity on $V\setminus\{v\}$. Also recall that $\ast_1$ and 
    $\ast_2$ are elements either introduced by a second derivative or the product of two derivatives (as introduced 
    in remark~\ref{remark_operations}).
\end{defi}


A {\em sub-operad} of an operad $\Operad$ is a sub-species of $\Operad$ containing the image of $e$ and 
which is stable under partial composition. For $\Operad$ an operad, the sub-operad of $\Operad$ 
{\em generated by} a sub-species $\spe{S}$ is the smallest sub-operad of $\Operad$ containing $\spe{S}$ and 
the sub-operad generated by a set $E$ of elements of $\Operad$ is the sub-operad generated by the species 
generated by $E$. A {\em morphism of operads} $f:\Operad_1\to\Operad_2$ is a morphism of species stable
under the structure maps: $f\circ e = e$ and $f(x\comp y) =f(x)\comp f(y)$.

In practice the map $e$ is often trivial and we do not mention it. Let us now give a series of examples of
operads.

\begin{description}\label{ex_operad}
    \item[Identity.] The identity species has a natural operad structure given by $v_1\comp v_2 = v_1$ if
    $v_1=\ast$ and $v_1\comp v_2=v_2$ otherwise.
    \item[Com.] The species $\Sym_+$ has a natural operad structure given by 
    $(V_1\cup\set{\ast})\comp V_2 = V_1\sqcup V_2$. This operad is called the {\em commutative operad} and denoted 
    by $\Com$. 
    In this context we denote by $\mu_V=V$ the basis element of $\Com[V]$.
    \item[NAP.] Let be $(t_1,r_1)$ and $(t_2,r_2)$ be two rooted trees and denote by $c(\ast)$ the child ends of $\ast$. 
    Let $(t_1,r_1)\comp (t_2,r_2)$ be the rooted tree obtained by the following construction.
    \begin{enumerate}
        \item Take the union of $t_1$ and $t_2$,

        \item remove the vertex $\ast$,

        \item connect the loose ends to $r_2$,

        \item if $r_1\not = \ast$, we choose $r_1$ as the root of the obtained tree, else we choose $r_2$.
    \end{enumerate}
    This construction is a partial composition and turns the species $\point{\T}$ of rooted trees into an 
    operad. This operad is called the {\em non-associative permutative operad}~\cite{Liv06} and is denoted 
    by $\NAP$. For instance we have:
    \begin{equation}
        \begin{tikzpicture}[Centering,scale=.7]
            \node[NodeGraph](s)at(0,0){$\ast$};
            \node[RootGraph](a)at(1,1){$a$};
            \node[NodeGraph](b)at(1,-1){$b$};
            \draw[EdgeGraph](a)--(s);
            \draw[EdgeGraph](s)--(b);
        \end{tikzpicture}
        \enspace \comp \enspace
        \begin{tikzpicture}[Centering,scale=.7]
            \node[RootGraph](c)at(0,0){$c$};
            \node[NodeGraph](d)at(1,0){$d$};
            \draw[EdgeGraph](c)--(d);
        \end{tikzpicture}
        \enspace = \enspace
        \begin{tikzpicture}[Centering,scale=.7]
            \node[RootGraph](a)at(1,1){$a$};
            \node[NodeGraph](b)at(1,-1){$b$};
            \node[NodeGraph](c)at(0,0){$c$};
            \node[NodeGraph](d)at(2,0){$d$};
            \draw[EdgeGraph](a)--(c);
            \draw[EdgeGraph](c)--(b);
            \draw[EdgeGraph](c)--(d);
        \end{tikzpicture}.
    \end{equation}
    \item[PreLie.]\label{ex_plie} Let be $(t_1,r_1)\in\T'[V_1]$ and $(t_2,r_2)\in\T[V_2]$ be two rooted trees and let $(t_1,r_1)\comp (t_2,r_2)$ be the 
    sum of rooted trees obtained by the following construction and indexed by the maps $f:\texttt{c}(\ast)\to V_2$, with $\texttt{c}(\ast)$ the set of 
    child ends of $\ast$.
    \begin{enumerate}
        \item Take the union of $t_1$ and $t_2$.

        \item remove the vertex $\ast$,

        \item connect the loose parent end to $r_2$,

        \item connect the loose child ends to their images by $f$,

        \item if $r_1\not = \ast$, we choose $r_1$ as the root of the obtained tree, else we choose $r_2$.
    \end{enumerate}
    This operation is a partial composition and turns the species $\point{\T}$ into an operad. This operad 
    is called {\em pre-Lie operad}~\cite{CL01} and is denoted by $\PLie$. Remark that the partial composition
    of $t_1$ and $t_2$ as elements of $\NAP$ is always in the support of the partial composition of $t_1$ and
    $t_2$ as elements of $\PLie$. For instance, we have:
    \begin{equation}
        \begin{tikzpicture}[Centering,scale=.7]
            \node[NodeGraph](s)at(0,0){$\ast$};
            \node[RootGraph](a)at(1,1){$a$};
            \node[NodeGraph](b)at(1,-1){$b$};
            \draw[EdgeGraph](a)--(s);
            \draw[EdgeGraph](s)--(b);
        \end{tikzpicture}
        \enspace \comp \enspace
        \begin{tikzpicture}[Centering,scale=.7]
            \node[RootGraph](c)at(0,0){$c$};
            \node[NodeGraph](d)at(1,0){$d$};
            \draw[EdgeGraph](c)--(d);
        \end{tikzpicture}
        \enspace = \enspace
        \begin{tikzpicture}[Centering,scale=.7]
            \node[RootGraph](a)at(1,1){$a$};
            \node[NodeGraph](b)at(1,-1){$b$};
            \node[NodeGraph](c)at(0,0){$c$};
            \node[NodeGraph](d)at(2,0){$d$};
            \draw[EdgeGraph](a)--(c);
            \draw[EdgeGraph](c)--(b);
            \draw[EdgeGraph](c)--(d);
        \end{tikzpicture}
        \enspace + \enspace
        \begin{tikzpicture}[Centering,scale=.7]
            \node[RootGraph](a)at(1,1){$a$};
            \node[NodeGraph](b)at(1,-1){$b$};
            \node[NodeGraph](c)at(0,0){$c$};
            \node[NodeGraph](d)at(2,0){$d$};
            \draw[EdgeGraph](a)--(c);
            \draw[EdgeGraph](d)--(b);
            \draw[EdgeGraph](c)--(d);
        \end{tikzpicture}.
    \end{equation}
    \item[Polynomials.] The species $\K\Poly$ has a natural partial composition given by the composition of
    polynomials: for $p_1(v_1,\dots,v_k,\ast)$ and $p_2(v_1',\dots,v_l')$ two polynomials over disjoint sets 
    of variables,
    \begin{equation}\label{pol_composition}
        (p_1\comp p_2)(v_1,\dots,v_k,v_1',\dots,v_l') = p_1|_{\ast \take p_2}
        = p_1(v_1,\dots, v_{k},p_2(v_1',\dots,v_l')).
    \end{equation}
    One can directly check that this partial composition satisfies the commutative diagrams of
    Definition~\ref{def_op}. This turns $\K\Poly$ into an operad where the units are the singleton 
    polynomials $v\in\Poly[\set{v}]$. Remark that since $\comp$ is a linear map, the product and
    addition of polynomials act as bilinear maps. Restricting the four species morphisms from 
    Example~\ref{ex_species} to $\Sym_+$ and $\K\Poly$ when necessary makes them operad morphisms.
    
    \item[Hadamard product.] If $\Operad_1$ and $\Operad_2$ are two operads, then the species 
    $\Operad_1\times\Operad_2$ is also an operad with partial composition: 
    $(x_1\otimes x_2)\comp (y_1\otimes y_2) =(x_1\comp y_1)\otimes (x_2\comp y_2)$. 
    
    \item[Assemblies.] For $\Operad$ an operad, the species $\Com(\Operad)$ of assemblies of $\Operad$ has 
    a natural operad structure. Let $V_1,\dots V_n$ and $W_1,\dots, W_k$ be $n+k$ disjoint sets such that 
    $\ast\in V_1$. Let be $x_i\in \Operad[V_i]$ for $1\leq i\leq n$ and $y_j\in\Operad[W_j]$ for 
    $1\leq j\leq k$.
    Then the partial composition of the assemblies $x_1\otimes\dots\otimes x_n$ and 
    $y_1\otimes\dots\otimes y_k$ is defined by
    $$(x_1\otimes\dots\otimes x_n) \comp (y_1\otimes\dots\otimes y_k) 
    = \sum_{i=1}^k y_1\otimes\dots\otimes y_{i-1}\otimes x_1\comp y_i \otimes y_{i+1}\otimes\dots\otimes y_k \otimes x_2\otimes\dots\otimes x_n.$$
\end{description}

In all the above examples, we defined operads by explicitly describing the vector spaces $\Operad[V]$ and 
the partial compositions. There is another way to present an operad which is as a quotient of a free operad.

Let $\spe{S}$ be a positive linear species such that $\spe{S}=\spe{X}+\spe{S}_{2^+}$. Recall that 
$\stree{\spe{S}}[V]$ is generated by trees with internal vertices decorated with elements of $\spe{S}$ 
and set of leaves $V$.
The species $\stree{\spe{S}}$ has then a natural operad structure given by the grafting of trees. For 
$t_1\in\stree{\spe{S}}'[V_1]$ and $t_2\in\stree{\spe{S}}[V_2]$, the partial composition $t_1\comp t_2$ is 
the tree obtained by grafting $t_2$ on the leaf $\ast$ of $t_1$ and relabeling the nodes of $t_1$ accordingly. 
This operad is called the {\em free operad over $\spe{S}$} and we denote it by $\free{\spe{S}}$.

\begin{exam}
    We give here an example of partial composition in a free operad over the same species as in Example~\ref{ex_stree}.
    \begin{equation}
        \begin{tikzpicture}[Centering, scale=0.6]
            \node[NodeFree, minimum size=8mm](r)at(0,0){$x_{12,\ast,3}$};
            \node[NodeFree, minimum size=8mm](n)at(-1.5,2.75){$y_{1,2}$};
            \node[NodeFree, minimum size=8mm](s)at(0.25,2.5){$\ast$};
            \node[NodeFree, minimum size=8mm](3)at(1.75,1.75){$3$};
            \node[NodeFree, minimum size=8mm](1)at(-3.5,5){$1$};
            \node[NodeFree, minimum size=8mm](2)at(-1.5,5.5){$2$};
            \draw[EdgeFree](r)--(n);
            \draw[EdgeFree](r)--(s);
            \draw[EdgeFree](r)--(3);
            \draw[EdgeFree](n)--(1);
            \draw[EdgeFree](n)--(2);
        \end{tikzpicture}
        \enspace \comp \enspace
        \begin{tikzpicture}[Centering, scale=0.6]
            \node[NodeFree, minimum size=8mm](r)at(0,0){$y_{4,5}$};
            \node[NodeFree, minimum size=8mm](4)at(-1.5,2.75){$4$};
            \node[NodeFree, minimum size=8mm](5)at(1.75,1.75){$5$};
            \draw[EdgeFree](r)--(4);
            \draw[EdgeFree](r)--(5);
        \end{tikzpicture}
        \enspace = \enspace
        \begin{tikzpicture}[Centering, scale=0.6]
            \node[NodeFree, minimum size=8mm](r)at(0,0){$x_{12,45,3}$};
            \node[NodeFree, minimum size=8mm](n)at(-1.5,2.75){$y_{1,2}$};
            \node[NodeFree, minimum size=8mm](s)at(0.25,2.5){$y_{4,5}$};
            \node[NodeFree, minimum size=8mm](4)at(0,5.25){$4$};
            \node[NodeFree, minimum size=8mm](5)at(2,4){$5$};
            \node[NodeFree, minimum size=8mm](3)at(1.75,1.75){$3$};
            \node[NodeFree, minimum size=8mm](1)at(-3.5,5){$1$};
            \node[NodeFree, minimum size=8mm](2)at(-1.5,5.5){$2$};
            \draw[EdgeFree](r)--(n);
            \draw[EdgeFree](r)--(s);
            \draw[EdgeFree](r)--(3);
            \draw[EdgeFree](n)--(1);
            \draw[EdgeFree](n)--(2);
            \draw[EdgeFree](s)--(4);
            \draw[EdgeFree](s)--(5);
        \end{tikzpicture}
    \end{equation}
\end{exam}

In the sequel, we consider free operads over species which are sub-species of an operad $\Operad$. When 
this happens, we denote by $\comp^\xi$ the partial composition in the free operad in order to not confuse 
it with the partial composition in $\Operad$.

\begin{exam}[$\ComMag$]
    The free operad over one symmetric generator $\free{\Sym_2}$ is the operad of abstract binary trees
    with partial composition the grafting of trees. This operad is called the {\em commutative magmatic 
    operad}~\cite{BL11} and is denoted by $\ComMag$. In this context, for $V$ a set of size $2$, we denote 
    by $s_V$ the generating element of $\Sym_2[V]$: $\Sym_2[V]=\free{\Sym_2}[V] = \K s_V$.
\end{exam}

An {\em ideal} of an operad $\Operad$ is a sub-species $\mathcal{I}$ such that the image of the products
$\Operad'\cdot \mathcal{I}$ and $\mathcal{I}'\cdot\Operad$ by the partial composition maps are in 
$\mathcal{I}$. The {\em quotient species} $\Operad/ \mathcal{I}$ defined by 
$(\Operad/\mathcal{I})[V] = \Operad[V]/\mathcal{I}[V]$ is then an operad with the natural partial 
composition and unit : $\iso{x}\comp\iso{y}=\iso{x\comp y}$ where $\iso{x}$ is the equivalence class of 
$x$. For $\mathcal{G}$ a species, if $\mathcal{R}$ is a sub-species of $\free{\mathcal{G}}$, we denote 
by $(\mathcal{R})$ the smallest ideal of $\free{\mathcal{G}}$ containing $\mathcal{R}$ and write that 
$(\mathcal{R})$ is {\em generated by $\mathcal{R}$}.

Denote by $\free{\mathcal{G}}^{(2)}$ the sub-species of $\free{\mathcal{G}}$ of trees with two internal nodes.

\begin{defi}
    Let $\mathcal{G}$ be a species and $\mathcal{R}$ be a sub-species of $\free{\mathcal{G}}$. We denote by 
    $\Ope(\mathcal{G},\mathcal{R})=\free{\mathcal{G}}/(\mathcal{R})$ the operad {\em generated by 
    $\mathcal{G}$ and with relations $\mathcal{R}$}. 
    The operad $\Ope(\mathcal{G},\mathcal{R})$ is {\em binary} if the species $\mathcal{G}$ of 
    {\em generators} is concentrated in cardinality $2$ (i.e. $\mathcal{G} = \mathcal{G}_2$). This operad 
    is {\em quadratic} if the species $\mathcal{R}$ of {\em relations} is a sub-species of 
    $\free{\mathcal{G}}^{(2)}$.
\end{defi}

\begin{exam}\label{ex_gen_com} 
    Denote by $\mathcal{R}$ the sub-species of $\ComMag^{(2)}$ generated by the {\em associativity relation}
    $s_{\set{a,\ast}}\comp^{\xi} s_{\set{b,c}} - s_{\set{c,\ast}}\comp^{\xi} s_{\set{a,b}}$. 
    Then $\Com=\Ope(\Sym_2,\mathcal{R})$ and $\Com$ is hence binary and quadratic. Remark that as a 
    consequence we have that $\mu_V$ is the image of $s_V$ under the projection $\ComMag\to\Com$.
\end{exam}

Two advantages of defining an operad by its generators and relations are that it is possible to construct 
(under some conditions) its Koszul dual and to check if the operad is Koszul. 
Koszul duality is a generalization of various dualities found in representation theory, e.g. that
of Lie algebras and commutative associative algebras. A Koszul operad is an operad with a 
subspecies of elements that act similarly to a Gröbner basis in a polynomial ring, except that it considers
Schröder trees instead of polynomials.
If $\Operad$ is a Koszul operad, it admits a Koszul dual $\Operad^!$ and the Hilbert series of $\Operad$ and $\Operad^!$ 
are related by the identity:
\begin{equation}
    \hilbert{\Operad}(-\hilbert{\Operad^!}(-t)) = t.
\end{equation}

These notions are too involved to be presented in a simple reminder and we only refer to them in 
Proposition~\ref{sp_dual} and Proposition~\ref{sp_koszul} in order to use the above equation. We refer the reader 
not versed in operad theory to Appendix~\ref{ann_koszul} for more information on these. 








\section{Extending the pre-Lie operad}\label{sec_extending_plie}
As announced, we want to extend the pre-Lie operad structure to graphs. As we will see later, it is more
natural to search an extension to multigraphs.




\subsection{Two canonical operads}\label{subsec_cano_op}

If we want to extends the above construction to multigraphs, we are faced with the problem that multigraphs 
do not have a root vertex. Our first solution is then to try to extend it to the pointed multigraphs 
$\point{\MG}$, where the pointed vertex would take the role of the roots. But this alone is not enough: 
indeed the partial composition of $\PLie$ also uses the notion 
of parent and child vertex/end which can not be defined on a general pointed multigraph. To remedy this, we 
consider oriented pointed multigraphs, where the targets play the same role as the parents ends in $\PLie$ 
and the sources the same role as the children ends.
Let then be $(g_1, \ori_1, v_1) \in (\point{\MG_{or}})'[V_1]$ and $(g_2, \ori_2, v_2) \in\point{\MG_{or}}[V_2]$ and define 
the partial composition $(g_1,\ori_1,v_1)\comp (g_2,\ori_2,v_2)$ as the sum of pointed oriented multigraphs obtained by the 
following construction and indexed by the maps $f:\texttt{t}(\ast)\to V_2$, with $\texttt{t}(\ast)$ the set of target ends of $\ast$.
\begin{enumerate}
    \item Take the union of $(g_1,\ori_1)$ and $(g_2,\ori_2)$,

    \item remove the vertex $\ast$,

    \item connect the loose source ends to $v_2$,

    \item connect the loose target ends to their images by $f$,

    \item the new pointed vertex is $v_1\comp v_2$, with the partial composition of the identity operad.
\end{enumerate}
For instance, we have:
\begin{equation}\label{canoor}
    \begin{tikzpicture}[Centering,scale=.7]
        \node[NodeGraph](s)at(0,0){$\ast$};
        \node[RootGraph](a)at(1,1){$a$};
        \node[NodeGraph](b)at(1,-1){$b$};
        \draw[ArcGraph](a)--(s);
        \draw[ArcGraph](s)--(b);
    \end{tikzpicture}
    \enspace \comp \enspace
    \begin{tikzpicture}[Centering,scale=1]
        \node[RootGraph](c)at(0,0){$c$};
        \node[NodeGraph](d)at(1,0){$d$};
        \draw[ArcGraph](c)--(d);
    \end{tikzpicture}
    \enspace = \enspace
    \begin{tikzpicture}[Centering,scale=.7]
        \node[RootGraph](a)at(1,1){$a$};
        \node[NodeGraph](b)at(1,-1){$b$};
        \node[NodeGraph](c)at(0,0){$c$};
        \node[NodeGraph](d)at(2,0){$d$};
        \draw[ArcGraph](a)--(c);
        \draw[ArcGraph](c)--(b);
        \draw[ArcGraph](c)--(d);
    \end{tikzpicture}
    \enspace + \enspace
    \begin{tikzpicture}[Centering,scale=.7]
        \node[RootGraph](a)at(1,1){$a$};
        \node[NodeGraph](b)at(1,-1){$b$};
        \node[NodeGraph](c)at(0,0){$c$};
        \node[NodeGraph](d)at(2,0){$d$};
        \draw[ArcGraph](a)--(d);
        \draw[ArcGraph](c)--(d);
        \draw[ArcGraph](c)--(b);
    \end{tikzpicture}.
\end{equation}

\begin{theorem}\label{th_op_graph_or}
    The species $\point{\MG_{or}}$, endowed with the preceding partial composition, is an operad.
\end{theorem}

We give a proof of this theorem in Section~\ref{sec_graph_operads} where we provide a general way to define 
operads on graphs and multigraphs.

It is straightforward to note that the subspecies of connected components $\point{\MG_{orc}}$ and the 
species $\point{\G_{or}}$ are sub-operads of $\point{\MG_{or}}$ and that $\point{\G_{orc}}$ is a sub-operad 
of $\point{\G_{or}}$. 
Given a rooted tree $(t,r)$, we have a natural orientation $\ori_{t,r}$ which consist of choosing the parent 
ends as sources and child ends as targets. This induces an operad monomorphism 
$(t,r)\mapsto(t,\ori_{t,r},r)$ from $\PLie$ to $\point{\G_{orc}}$ and hence makes $\PLie$ a sub-operad of 
$\point{\G}_{orc}$.

Our second solution in extending $\PLie$ is to ignore the steps of the construction of the partial 
composition which involve the root. The resulting partial composition is then much more natural than the 
previous one over $\point{\MG}_{or}$: let $g_1\in\MG'[V_1]$ and $g_2\in\MG[V_2]$ be two multigraphs and 
define the partial composition $g_1\comp g_2$ as the sum of oriented multigraphs obtained by the 
following construction and indexed by the maps $f:n(\ast)\to V_2$, with $n(\ast)$ the ends of $\ast$.
\begin{enumerate}
    \item Take the union of $g_1$ and $g_2$,

    \item remove the vertex $\ast$,

    \item connect the loose ends to their images by $f$.
\end{enumerate}
For instance, we have:
\begin{align}
    \begin{tikzpicture}[Centering,scale=1]
        \tikzset{every loop/.style={}}
        \node[NodeGraph](a)at(0,0){$a$};
        \node[NodeGraph](s)at(1,0){$\ast$};
        \draw[EdgeGraph](a)edge[bend left=40](s);
        \draw[EdgeGraph](a)edge[bend right=40](s);
        \draw[EdgeGraph](s)edge[loop above](s);
        \draw[EdgeGraph,draw=ColWhite](s)edge[loop below](s);
    \end{tikzpicture}
    &\enspace \comp \enspace
    \begin{tikzpicture}[Centering,scale=1]
        \tikzset{every loop/.style={}}
        \node[NodeGraph](b)at(0,0){$b$};
        \node[NodeGraph](c)at(1,0){$c$};
        \draw[EdgeGraph](b)--(c);
        \draw[EdgeGraph](c)edge[loop above](c);
        \draw[EdgeGraph,draw=ColWhite](s)edge[loop below](s);
    \end{tikzpicture}
    \enspace = \\ \nonumber \enspace
    &\begin{tikzpicture}[Centering,scale=1]
        \tikzset{every loop/.style={}}
        \node[NodeGraph](a)at(0,0){$a$};
        \node[NodeGraph](b)at(1,0){$b$};
        \node[NodeGraph](c)at(2,0){$c$}; 
        \draw[EdgeGraph](a)edge[bend left=40](b);
        \draw[EdgeGraph](a)edge[bend right=40](b);
        \draw[EdgeGraph](b)edge[loop below](b);
        \draw[EdgeGraph](b)--(c);
        \draw[EdgeGraph](c)edge[loop above](c);
    \end{tikzpicture}
    \enspace + \enspace
    \begin{tikzpicture}[Centering,scale=1]
        \tikzset{every loop/.style={}}
        \node[NodeGraph](a)at(0,0){$a$};
        \node[NodeGraph](b)at(1,0){$b$};
        \node[NodeGraph](c)at(2,0){$c$};
        \draw[EdgeGraph](a)edge[bend left=40](b);
        \draw[EdgeGraph](a)edge[bend right=40](b);
        \draw[EdgeGraph](c)edge[loop below](c);
        \draw[EdgeGraph](b)--(c);
        \draw[EdgeGraph](c)edge[loop above](c);
    \end{tikzpicture}
    \enspace + \enspace
    2\,
    \begin{tikzpicture}[Centering,scale=1]
        \tikzset{every loop/.style={}}
        \node[NodeGraph](a)at(0,0){$a$};
        \node[NodeGraph](b)at(1,0){$b$};
        \node[NodeGraph](c)at(2,0){$c$};
        \draw[EdgeGraph](a)edge[bend left=40](b);
        \draw[EdgeGraph](a)edge[bend right=40](b);
        \draw[EdgeGraph](b)edge[bend left=40](c);
        \draw[EdgeGraph](b)edge[bend right=40](c);
        \draw[EdgeGraph](c)edge[loop above](c);
        \draw[EdgeGraph,draw=ColWhite](c)edge[loop below](c);
    \end{tikzpicture}
    \\ \nonumber
    &+ \enspace
    2\,
    \begin{tikzpicture}[Centering,scale=1]
        \tikzset{every loop/.style={}}
        \node[NodeGraph](a)at(0,0){$a$};
        \node[NodeGraph](b)at(1,0){$b$};
        \node[NodeGraph](c)at(2,0){$c$};
        \draw[EdgeGraph](a)--(b);
        \draw[EdgeGraph](a)edge[bend right=40](c);
        \draw[EdgeGraph](b)edge[loop above](b);
        \draw[EdgeGraph](b)--(c);
        \draw[EdgeGraph](c)edge[loop above](c);
    \end{tikzpicture}
    \enspace + \enspace
    2\,
    \begin{tikzpicture}[Centering,scale=1]
        \tikzset{every loop/.style={}}
        \node[NodeGraph](a)at(0,0){$a$};
        \node[NodeGraph](b)at(1,0){$b$};
        \node[NodeGraph](c)at(2,0){$c$};
        \draw[EdgeGraph](a)--(b);
        \draw[EdgeGraph](a)edge[bend right=40](c);
        \draw[EdgeGraph](c)edge[loop below](c);
        \draw[EdgeGraph](b)--(c);
        \draw[EdgeGraph](c)edge[loop above](c);
    \end{tikzpicture}
    \enspace + \enspace
    4\,
    \begin{tikzpicture}[Centering,scale=1]
        \tikzset{every loop/.style={}}
        \node[NodeGraph](a)at(0,0){$a$};
        \node[NodeGraph](b)at(1,0){$b$};
        \node[NodeGraph](c)at(2,0){$c$};
        \draw[EdgeGraph](a)--(b);
        \draw[EdgeGraph](a)edge[bend right=40](c);
        \draw[EdgeGraph](b)edge[bend right=40](c);
        \draw[EdgeGraph](b)edge[bend left=40](c);
        \draw[EdgeGraph](c)edge[loop above](c);
    \end{tikzpicture}
    \\ \nonumber
    &+ \enspace
    \begin{tikzpicture}[Centering,scale=1]
        \tikzset{every loop/.style={}}
        \node[NodeGraph](a)at(0,0){$a$};
        \node[NodeGraph](b)at(1,0){$b$};
        \node[NodeGraph](c)at(2,0){$c$};
        \draw[EdgeGraph](a)edge[bend right=40](c);
        \draw[EdgeGraph](a)edge[bend left=40](c);
        \draw[EdgeGraph](b)edge[loop left](b);
        \draw[EdgeGraph](b)--(c);
        \draw[EdgeGraph](c)edge[loop above](c);
    \end{tikzpicture}
    \enspace + \enspace
    \begin{tikzpicture}[Centering,scale=1]
        \tikzset{every loop/.style={}}
        \node[NodeGraph](a)at(0,0){$a$};
        \node[NodeGraph](b)at(1,0){$b$};
        \node[NodeGraph](c)at(2,0){$c$};
        \draw[EdgeGraph](a)edge[bend right=40](c);
        \draw[EdgeGraph](a)edge[bend left=40](c);
        \draw[EdgeGraph](c)edge[loop below](c);
        \draw[EdgeGraph](b)--(c);
        \draw[EdgeGraph](c)edge[loop above](c);
    \end{tikzpicture}
    \enspace + \enspace
    2\,
    \begin{tikzpicture}[Centering,scale=1]
        \tikzset{every loop/.style={}}
        \node[NodeGraph](a)at(0,0){$a$};
        \node[NodeGraph](b)at(1,0){$b$};
        \node[NodeGraph](c)at(2,0){$c$};
        \draw[EdgeGraph](a)edge[bend right=40](c);
        \draw[EdgeGraph](a)edge[bend left=40](c);
        \draw[EdgeGraph](b)edge[bend left=20](c);
        \draw[EdgeGraph](b)edge[bend right=20](c);
        \draw[EdgeGraph](c)edge[loop above](c);
    \end{tikzpicture}.
\end{align}


\begin{theorem} \label{th_graphop_cano}
    The species $\MG$ endowed with the preceding partial composition is an operad.
\end{theorem}

As for Theorem~\ref{th_op_graph_or}, we give a proof of this theorem in Section~\ref{sec_graph_operads}.

This operad structure makes all the species of the diagram~\ref{graph_diagram} operads and its maps operad 
monomorphisms. In particular, we recover the Kontsevich-Willwacher operad~\cite{MV19} on $\G$. Recall now 
from Section~\ref{sec_intro} that we can identify multigraphs with polynomials. The partial composition we 
just defined can then be formally written as $g_1\comp g_2 = g_1|_{\ast \take \sum V_2}\oplus g_2$
(using the same notation for the composition of polynomials as in \eqref{pol_composition}), which can then 
be expanded as follows:
\begin{equation}\begin{split}\label{comp_multigraphs}
    g_1\comp g_2 &= g_1|_{\ast \take \sum V_2}\oplus g_2\\
    &= g_1|_{V_1} \oplus\bigoplus_{v\in n(\ast)} v(\sum V_2)\oplus((\sum V_2)^2)^{\oplus g_1(\ast\ast)}\oplus g_2 \\
    &= \sum_{f:n(\ast)\to V_2}\sum_{l:[g_1(\ast\ast)]\to V_2V_2} g_1|_{V_1}\oplus\bigoplus_{v\in n(\ast)}vf(v)\oplus\bigoplus_{i=1}^{g_1(\ast\ast)}l(i) \oplus g_2,
\end{split}\end{equation}
where $n(\ast)$ is the multiset of neighbours of $\ast$ in $g_1$ and $g_1(\ast\ast)$ is the number of loops on 
$\ast$ in $g_1$. Each of the three terms of the second line, without counting $g_2$, has a combinatorial 
interpretation: $g_1|_{V_1}$ is $g_1$ to which we removed $\ast$, $\bigoplus_{v\in n(\ast)} v(\sum V_2)$
can be understood as ``for all vertices $v$ of $n(\ast)$, sum over the ways of connecting $v$ to $g_2$'' 
and the term $((\sum V_2)^2)^{\oplus g_1(\ast\ast)}$ as ``for each loop over $\ast$, add an edge between 
any two elements of $V_2$''. This partial composition expands in a simpler way on $\G$ because of the 
absence of loops. Indeed, if $g_1$ and $g_2$ are now graphs, Equ.~\eqref{comp_multigraphs} rewrites as
\begin{equation}\begin{split}
    g_1\comp g_2 &= g_1|_{\ast \take \sum V_2}\oplus g_2\\
    &= g_1|_{V_1} \oplus\bigoplus_{v\in n(\ast)} v(\sum V_2)\oplus g_2 \\
    &= \sum_{f:n(\ast)\rightarrow V_2} g_1|_{V_1}\oplus\bigoplus_{v\in n(\ast)}vf(v)\oplus g_2.
\end{split}\end{equation}
For instance, we have:
\begin{equation}
    \begin{tikzpicture}[Centering,scale=.5]
        \node[NodeGraph](a)at(1,1){$a$};
        \node[NodeGraph](s)at(0,0){$\ast$};
        \node[NodeGraph](b)at(1,-1){$b$};
        \draw[EdgeGraph](a)--(s);
        \draw[EdgeGraph](s)--(b);
    \end{tikzpicture}
    \enspace \comp \enspace
    \begin{tikzpicture}[Centering,scale=.7]
        \node[NodeGraph](c)at(0,0){$c$};
        \node[NodeGraph](d)at(1,0){$d$};
        \draw[EdgeGraph](c)--(d);
    \end{tikzpicture}
    \enspace = \enspace
    \begin{tikzpicture}[Centering,scale=.5]
        \node[NodeGraph](a)at(1,1){$a$};
        \node[NodeGraph](b)at(1,-1){$b$};
        \node[NodeGraph](c)at(0,0){$c$};
        \node[NodeGraph](d)at(2,0){$d$};
        \draw[EdgeGraph](c)--(a);
        \draw[EdgeGraph](c)--(d);
        \draw[EdgeGraph](c)--(b);
    \end{tikzpicture}
    \enspace + \enspace
    \begin{tikzpicture}[Centering,scale=.5]
        \node[NodeGraph](a)at(1,1){$a$};
        \node[NodeGraph](b)at(1,-1){$b$};
        \node[NodeGraph](c)at(0,0){$c$};
        \node[NodeGraph](d)at(2,0){$d$};
        \draw[EdgeGraph](c)--(d);
        \draw[EdgeGraph](d)--(a);
        \draw[EdgeGraph](c)--(b);
    \end{tikzpicture}
    \enspace + \enspace
    \begin{tikzpicture}[Centering,scale=.5]
        \node[NodeGraph](a)at(1,1){$a$};
        \node[NodeGraph](b)at(1,-1){$b$};
        \node[NodeGraph](c)at(0,0){$c$};
        \node[NodeGraph](d)at(2,0){$d$};
        \draw[EdgeGraph](c)--(d);
        \draw[EdgeGraph](c)--(a);
        \draw[EdgeGraph](d)--(b);
    \end{tikzpicture}
    \enspace + \enspace
    \begin{tikzpicture}[Centering,scale=.5]
        \node[NodeGraph](a)at(1,1){$a$};
        \node[NodeGraph](b)at(1,-1){$b$};
        \node[NodeGraph](c)at(0,0){$c$};
        \node[NodeGraph](d)at(2,0){$d$};
        \draw[EdgeGraph](c)--(d);
        \draw[EdgeGraph](d)--(a);
        \draw[EdgeGraph](d)--(b);
    \end{tikzpicture}.
\end{equation}
In particular, we observe that all graphs appearing in $g_1 \comp g_2$ have $1$ as coefficient.



\subsection{Link with \texorpdfstring{$\PLie$}{PLie}}
While $\PLie$ is a sub-operad of $\point{\MG}_{or}$, the operad structure on $\MG$ does not seem to have any 
relation with $\PLie$. In fact, as we will see at the end of this subsection, there is a non-trivial 
correspondence between the four operads $\point{\MG}_{or}$, $\MG$, $\PLie$, and the Kontsevich-Willwacher 
operad $\G$. Let us begin with the following result which gives a correspondence between the sub-operad $\T$ 
of $\MG$ and $\PLie$.

\begin{prop} \label{prop_prelie}
    The monomorphism of species $\psi : \T \to \point{\T}$ defined by, for any tree $t \in \T[V]$,
    \begin{equation}
        \psi(t) =  \sum_{r \in V} (t, r),
    \end{equation}
    is a monomorphism of operads from $\T$ to $\PLie$.
\end{prop}

Before giving the proof of this proposition, we illustrate it on an example:
\begin{equation}\begin{split}
    \psi&\left(\ 
    \begin{tikzpicture}[Centering, scale=0.6]
        \node[NodeGraph, minimum size=3mm](a)at(0,0){$a$};
        \node[NodeGraph, minimum size=3mm](b)at(0.5,0.86){$b$};
        \node[NodeGraph, minimum size=3mm](s)at(1,0){$\ast$};
        \draw[EdgeGraph](a)--(b);
        \draw[EdgeGraph](b)--(s);
    \end{tikzpicture}
    \ \right) \comp
    \psi\left(\ 
    \begin{tikzpicture}[Centering, scale=0.6]
        \node[NodeGraph, minimum size=3mm](c)at(0,1){$c$};
        \node[NodeGraph, minimum size=3mm](d)at(0,0){$d$};
        \draw[EdgeGraph](c)--(d);
    \end{tikzpicture}
    \ \right)=\left(\ 
    \begin{tikzpicture}[Centering, scale=0.6]
        \node[RootGraph, minimum size=3mm](a)at(0,0){$a$};
        \node[NodeGraph, minimum size=3mm](b)at(0.5,0.86){$b$};
        \node[NodeGraph, minimum size=3mm](s)at(1,0){$\ast$};
        \draw[EdgeGraph](a)--(b);
        \draw[EdgeGraph](b)--(s);
    \end{tikzpicture}
    \  + \ 
    \begin{tikzpicture}[Centering, scale=0.6]
        \node[NodeGraph, minimum size=3mm](a)at(0,0){$a$};
        \node[RootGraph, minimum size=3mm](b)at(0.5,0.86){$b$};
        \node[NodeGraph, minimum size=3mm](s)at(1,0){$\ast$};
        \draw[EdgeGraph](a)--(b);
        \draw[EdgeGraph](b)--(s);
    \end{tikzpicture}
    \  + \ 
    \begin{tikzpicture}[Centering, scale=0.6]
        \node[NodeGraph, minimum size=3mm](a)at(0,0){$a$};
        \node[NodeGraph, minimum size=3mm](b)at(0.5,0.86){$b$};
        \node[RootGraph, minimum size=3mm](s)at(1,0){$\ast$};
        \draw[EdgeGraph](a)--(b);
        \draw[EdgeGraph](b)--(s);
    \end{tikzpicture}
    \ \right)\comp\left(\ 
    \begin{tikzpicture}[Centering, scale=0.6]
        \node[RootGraph, minimum size=3mm](c)at(0,1){$c$};
        \node[NodeGraph, minimum size=3mm](d)at(0,0){$d$};
        \draw[EdgeGraph](c)--(d);
    \end{tikzpicture}
    \  + \ 
    \begin{tikzpicture}[Centering, scale=0.6]
        \node[NodeGraph, minimum size=3mm](c)at(0,0){$c$};
        \node[RootGraph, minimum size=3mm](d)at(0,1){$d$};
        \draw[EdgeGraph](c)--(d);
    \end{tikzpicture}
    \ \right) \\
    &= 
    \begin{tikzpicture}[Centering, scale=0.6]
        \node[RootGraph, minimum size=3mm](a)at(0,0){$a$};
        \node[NodeGraph, minimum size=3mm](b)at(0.5,0.86){$b$};
        \node[NodeGraph, minimum size=3mm](s)at(1,0){$\ast$};
        \draw[EdgeGraph](a)--(b);
        \draw[EdgeGraph](b)--(s);
    \end{tikzpicture}
    \comp\left(\ 
    \begin{tikzpicture}[Centering, scale=0.6]
        \node[RootGraph, minimum size=3mm](c)at(0,1){$c$};
        \node[NodeGraph, minimum size=3mm](d)at(0,0){$d$};
        \draw[EdgeGraph](c)--(d);
    \end{tikzpicture}
    \  + \ 
    \begin{tikzpicture}[Centering, scale=0.6]
        \node[NodeGraph, minimum size=3mm](c)at(0,0){$c$};
        \node[RootGraph, minimum size=3mm](d)at(0,1){$d$};
        \draw[EdgeGraph](c)--(d);
    \end{tikzpicture}
    \ \right)
    \  + \ 
    \begin{tikzpicture}[Centering, scale=0.6]
        \node[NodeGraph, minimum size=3mm](a)at(0,0){$a$};
        \node[RootGraph, minimum size=3mm](b)at(0.5,0.86){$b$};
        \node[NodeGraph, minimum size=3mm](s)at(1,0){$\ast$};
        \draw[EdgeGraph](a)--(b);
        \draw[EdgeGraph](b)--(s);
    \end{tikzpicture}
    \comp\left(\ 
    \begin{tikzpicture}[Centering, scale=0.6]
        \node[RootGraph, minimum size=3mm](c)at(0,1){$c$};
        \node[NodeGraph, minimum size=3mm](d)at(0,0){$d$};
        \draw[EdgeGraph](c)--(d);
    \end{tikzpicture}
    \  + \ 
    \begin{tikzpicture}[Centering, scale=0.6]
        \node[NodeGraph, minimum size=3mm](c)at(0,0){$c$};
        \node[RootGraph, minimum size=3mm](d)at(0,1){$d$};
        \draw[EdgeGraph](c)--(d);
    \end{tikzpicture}
    \ \right)
    \  + \ 
    \begin{tikzpicture}[Centering, scale=0.6]
        \node[NodeGraph, minimum size=3mm](a)at(0,0){$a$};
        \node[NodeGraph, minimum size=3mm](b)at(0.5,0.86){$b$};
        \node[RootGraph, minimum size=3mm](s)at(1,0){$\ast$};
        \draw[EdgeGraph](a)--(b);
        \draw[EdgeGraph](b)--(s);
    \end{tikzpicture}
    \comp\left(\ 
    \begin{tikzpicture}[Centering, scale=0.6]
        \node[RootGraph, minimum size=3mm](c)at(0,1){$c$};
        \node[NodeGraph, minimum size=3mm](d)at(0,0){$d$};
        \draw[EdgeGraph](c)--(d);
    \end{tikzpicture}
    \  + \ 
    \begin{tikzpicture}[Centering, scale=0.6]
        \node[NodeGraph, minimum size=3mm](c)at(0,0){$c$};
        \node[RootGraph, minimum size=3mm](d)at(0,1){$d$};
        \draw[EdgeGraph](c)--(d);
    \end{tikzpicture}
    \ \right) \\
    &=
    \begin{tikzpicture}[Centering, scale=0.6]
        \node[RootGraph, minimum size=3mm](a)at(0,0){$a$};
        \node[NodeGraph, minimum size=3mm](b)at(0.5,0.86){$b$};
        \node[NodeGraph, minimum size=3mm](c)at(1,0){$c$};
        \draw[EdgeGraph](a)--(b);
        \draw[EdgeGraph](b)--(c);
        \node[NodeGraph, minimum size=3mm](d)at(1,-1){$d$};
        \draw[EdgeGraph](c)--(d);
    \end{tikzpicture}
    \  + \ 
    \begin{tikzpicture}[Centering, scale=0.6]
        \node[RootGraph, minimum size=3mm](a)at(0,0){$a$};
        \node[NodeGraph, minimum size=3mm](b)at(0.5,0.86){$b$};
        \node[NodeGraph, minimum size=3mm](c)at(1,0){$d$};
        \draw[EdgeGraph](a)--(b);
        \draw[EdgeGraph](b)--(c);
        \node[NodeGraph, minimum size=3mm](d)at(1,-1){$c$};
        \draw[EdgeGraph](c)--(d);
    \end{tikzpicture}
    \  + \ 
    \begin{tikzpicture}[Centering, scale=0.6]
        \node[NodeGraph, minimum size=3mm](a)at(0,0){$a$};
        \node[RootGraph, minimum size=3mm](b)at(0.5,0.86){$b$};
        \node[NodeGraph, minimum size=3mm](c)at(1,0){$c$};
        \draw[EdgeGraph](a)--(b);
        \draw[EdgeGraph](b)--(c);
        \node[NodeGraph, minimum size=3mm](d)at(1,-1){$d$};
        \draw[EdgeGraph](c)--(d);
    \end{tikzpicture}
    \  + \ 
    \begin{tikzpicture}[Centering, scale=0.6]
        \node[NodeGraph, minimum size=3mm](a)at(0,0){$a$};
        \node[RootGraph, minimum size=3mm](b)at(0.5,0.86){$b$};
        \node[NodeGraph, minimum size=3mm](c)at(1,0){$d$};
        \draw[EdgeGraph](a)--(b);
        \draw[EdgeGraph](b)--(c);
        \node[NodeGraph, minimum size=3mm](d)at(1,-1){$c$};
        \draw[EdgeGraph](c)--(d);
    \end{tikzpicture}
    \  + \ 
    \begin{tikzpicture}[Centering, scale=0.6]
        \node[NodeGraph, minimum size=3mm](a)at(0,0){$a$};
        \node[NodeGraph, minimum size=3mm](b)at(0.5,0.86){$b$};
        \node[RootGraph, minimum size=3mm](c)at(1,0){$c$};
        \draw[EdgeGraph](a)--(b);
        \draw[EdgeGraph](b)--(c);
        \node[NodeGraph, minimum size=3mm](d)at(1,-1){$d$};
        \draw[EdgeGraph](c)--(d);
    \end{tikzpicture}
    \  + \ 
    \begin{tikzpicture}[Centering, scale=0.6]
        \node[NodeGraph, minimum size=3mm](a)at(0,0){$a$};
        \node[NodeGraph, minimum size=3mm](b)at(0.5,0.86){$b$};
        \node[NodeGraph, minimum size=3mm](c)at(1,0){$d$};
        \draw[EdgeGraph](a)--(b);
        \draw[EdgeGraph](b)--(c);
        \node[RootGraph, minimum size=3mm](d)at(1,-1){$c$};
        \draw[EdgeGraph](c)--(d);
    \end{tikzpicture}
    \  + \ 
    \begin{tikzpicture}[Centering, scale=0.6]
        \node[NodeGraph, minimum size=3mm](a)at(0,0){$a$};
        \node[NodeGraph, minimum size=3mm](b)at(0.5,0.86){$b$};
        \node[RootGraph, minimum size=3mm](c)at(1,0){$d$};
        \draw[EdgeGraph](a)--(b);
        \draw[EdgeGraph](b)--(c);
        \node[NodeGraph, minimum size=3mm](d)at(1,-1){$c$};
        \draw[EdgeGraph](c)--(d);
    \end{tikzpicture}
    \  + \ 
    \begin{tikzpicture}[Centering, scale=0.6]
        \node[NodeGraph, minimum size=3mm](a)at(0,0){$a$};
        \node[NodeGraph, minimum size=3mm](b)at(0.5,0.86){$b$};
        \node[NodeGraph, minimum size=3mm](c)at(1,0){$c$};
        \draw[EdgeGraph](a)--(b);
        \draw[EdgeGraph](b)--(c);
        \node[RootGraph, minimum size=3mm](d)at(1,-1){$d$};
        \draw[EdgeGraph](c)--(d);
    \end{tikzpicture} \\
    &=\left(\ 
    \begin{tikzpicture}[Centering, scale=0.6]
        \node[RootGraph, minimum size=3mm](a)at(0,0){$a$};
        \node[NodeGraph, minimum size=3mm](b)at(0.5,0.86){$b$};
        \node[NodeGraph, minimum size=3mm](c)at(1,0){$c$};
        \draw[EdgeGraph](a)--(b);
        \draw[EdgeGraph](b)--(c);
        \node[NodeGraph, minimum size=3mm](d)at(1,-1){$d$};
        \draw[EdgeGraph](c)--(d);
    \end{tikzpicture}
    +
    \begin{tikzpicture}[Centering, scale=0.6]
        \node[NodeGraph, minimum size=3mm](a)at(0,0){$a$};
        \node[RootGraph, minimum size=3mm](b)at(0.5,0.86){$b$};
        \node[NodeGraph, minimum size=3mm](c)at(1,0){$c$};
        \draw[EdgeGraph](a)--(b);
        \draw[EdgeGraph](b)--(c);
        \node[NodeGraph, minimum size=3mm](d)at(1,-1){$d$};
        \draw[EdgeGraph](c)--(d);
    \end{tikzpicture}
    +
    \begin{tikzpicture}[Centering, scale=0.6]
        \node[NodeGraph, minimum size=3mm](a)at(0,0){$a$};
        \node[NodeGraph, minimum size=3mm](b)at(0.5,0.86){$b$};
        \node[RootGraph, minimum size=3mm](c)at(1,0){$c$};
        \draw[EdgeGraph](a)--(b);
        \draw[EdgeGraph](b)--(c);
        \node[NodeGraph, minimum size=3mm](d)at(1,-1){$d$};
        \draw[EdgeGraph](c)--(d);
    \end{tikzpicture}
    +
    \begin{tikzpicture}[Centering, scale=0.6]
        \node[NodeGraph, minimum size=3mm](a)at(0,0){$a$};
        \node[NodeGraph, minimum size=3mm](b)at(0.5,0.86){$b$};
        \node[NodeGraph, minimum size=3mm](c)at(1,0){$c$};
        \draw[EdgeGraph](a)--(b);
        \draw[EdgeGraph](b)--(c);
        \node[RootGraph, minimum size=3mm](d)at(1,-1){$d$};
        \draw[EdgeGraph](c)--(d);
    \end{tikzpicture}
    \ \right)
    +
    \left(\ 
    \begin{tikzpicture}[Centering, scale=0.6]
        \node[RootGraph, minimum size=3mm](a)at(0,0){$a$};
        \node[NodeGraph, minimum size=3mm](b)at(0.5,0.86){$b$};
        \node[NodeGraph, minimum size=3mm](c)at(1,0){$d$};
        \draw[EdgeGraph](a)--(b);
        \draw[EdgeGraph](b)--(c);
        \node[NodeGraph, minimum size=3mm](d)at(1,-1){$c$};
        \draw[EdgeGraph](c)--(d);
    \end{tikzpicture}
    +
    \begin{tikzpicture}[Centering, scale=0.6]
        \node[NodeGraph, minimum size=3mm](a)at(0,0){$a$};
        \node[RootGraph, minimum size=3mm](b)at(0.5,0.86){$b$};
        \node[NodeGraph, minimum size=3mm](c)at(1,0){$d$};
        \draw[EdgeGraph](a)--(b);
        \draw[EdgeGraph](b)--(c);
        \node[NodeGraph, minimum size=3mm](d)at(1,-1){$c$};
        \draw[EdgeGraph](c)--(d);
    \end{tikzpicture}
    +
    \begin{tikzpicture}[Centering, scale=0.6]
        \node[NodeGraph, minimum size=3mm](a)at(0,0){$a$};
        \node[NodeGraph, minimum size=3mm](b)at(0.5,0.86){$b$};
        \node[RootGraph, minimum size=3mm](c)at(1,0){$d$};
        \draw[EdgeGraph](a)--(b);
        \draw[EdgeGraph](b)--(c);
        \node[NodeGraph, minimum size=3mm](d)at(1,-1){$c$};
        \draw[EdgeGraph](c)--(d);
    \end{tikzpicture}
    +
    \begin{tikzpicture}[Centering, scale=0.6]
        \node[NodeGraph, minimum size=3mm](a)at(0,0){$a$};
        \node[NodeGraph, minimum size=3mm](b)at(0.5,0.86){$b$};
        \node[NodeGraph, minimum size=3mm](c)at(1,0){$d$};
        \draw[EdgeGraph](a)--(b);
        \draw[EdgeGraph](b)--(c);
        \node[RootGraph, minimum size=3mm](d)at(1,-1){$c$};
        \draw[EdgeGraph](c)--(d);
    \end{tikzpicture}
    \ \right)\\
    &= \psi\left(\ 
        \begin{tikzpicture}[Centering, scale=0.6]
        \node[NodeGraph, minimum size=3mm](a)at(0,0){$a$};
        \node[NodeGraph, minimum size=3mm](b)at(0.5,0.86){$b$};
        \node[NodeGraph, minimum size=3mm](c)at(1,0){$c$};
        \draw[EdgeGraph](a)--(b);
        \draw[EdgeGraph](b)--(c);
        \node[NodeGraph, minimum size=3mm](d)at(1,-1){$d$};
        \draw[EdgeGraph](c)--(d);
    \end{tikzpicture}
    \  + \ 
    \begin{tikzpicture}[Centering, scale=0.6]
        \node[NodeGraph, minimum size=3mm](a)at(0,0){$a$};
        \node[NodeGraph, minimum size=3mm](b)at(0.5,0.86){$b$};
        \node[NodeGraph, minimum size=3mm](c)at(1,0){$d$};
        \draw[EdgeGraph](a)--(b);
        \draw[EdgeGraph](b)--(c);
        \node[NodeGraph, minimum size=3mm](d)at(1,-1){$c$};
        \draw[EdgeGraph](c)--(d);
    \end{tikzpicture}
    \ \right) =
    \psi\left(\ 
    \begin{tikzpicture}[Centering, scale=0.6]
        \node[NodeGraph, minimum size=3mm](a)at(0,0){$a$};
        \node[NodeGraph, minimum size=3mm](b)at(0.5,0.86){$b$};
        \node[NodeGraph, minimum size=3mm](s)at(1,0){$\ast$};
        \draw[EdgeGraph](a)--(b);
        \draw[EdgeGraph](b)--(s);
    \end{tikzpicture}
    \comp
    \begin{tikzpicture}[Centering, scale=0.6]
        \node[NodeGraph, minimum size=3mm](c)at(0,1){$c$};
        \node[NodeGraph, minimum size=3mm](d)at(0,0){$d$};
        \draw[EdgeGraph](c)--(d);
    \end{tikzpicture}
    \ \right).
\end{split}\end{equation}
We can note that the case when $\ast$ is the root plays a particular role.

\begin{proof}
    For $t\in \T[V]$ a tree and $r,v\in V$, we denote by $n_t(v)$ the set of neighbours of $v$ in $t$, by 
    $c_{t,r}(v)$ the set of children of $v$ in the rooted tree $(t,r)$ and if $r\not = v$, we further denote 
    by $p_{t,r}(v)$ the parent of $v$ in~$(t,r)$.

    Let be $t_1\in \T'[V_1]$ and $t_2\in \T[V_2]$. We now make full use of the correspondence between 
    graphs and polynomials:
    \begin{align}
        \psi_{V_1}(t_1)\comp &\psi_{V_2}(t_2) 
        = \sum_{r_1\in V_1\cup\set{\ast}} (t_1,r_1) \comp \sum_{r_2\in V_2}(t_2,r_2) \\ \nonumber
        &= \sum_{r_1\in V_1\cup\set{\ast}}\sum_{r_2\in V_2} (t_1,r_1)\comp (t_2,r_2) \\ \nonumber 
        &= \sum_{r_1\in V_1}\sum_{r_2\in V_2} \left(t_1|_{V_1}\oplus p_{t_1,r_1}(\ast)r_2
        \oplus t_2 \oplus \bigoplus_{v\in c_{t_1,r_1}(\ast)} v\left(\sum V_2\right), r_1\right) \\ \nonumber
        &\quad + \sum_{r_2\in V_2}\left(t_1|_{V_1}\oplus t_2 \oplus \bigoplus_{v\in c_{t_1,\ast}(\ast)} v\left(\sum V_2\right),r_2\right) \\ \nonumber
        &= \sum_{r_1\in V_1} \left(t_1|_{V_1}\oplus p_{t_1,r_1}(\ast)\left(\sum V_2\right) \oplus t_2 \oplus \bigoplus_{v\in c_{t_1,r_1}(\ast)}v\left(\sum V_2\right), r_1\right) \\ \nonumber
        &\quad+ \sum_{r_2\in V_2} \left(t_1|_{V_1}\oplus t_2 \oplus \bigoplus_{v\in c_{t_1,\ast}(\ast)} v\left(\sum V_2\right),r_2\right) \\ \nonumber
        &= \sum_{r\in V_1+V_2}\left(t_1|_{V_1}\oplus\bigoplus_{v\in n_{t_1}(\ast)}v\left(\sum V_2\right)\oplus t_2,r\right) \\ \nonumber
        &= \sum_{r\in V_1+V_2} \left(t_1|_{\ast \take \sum V_2}\oplus t_2,r\right) \\ \nonumber
        &= \psi_{V_1+V_2}(t_1\comp t_2). \qedhere
    \end{align}
\end{proof}


The map $\psi$ naturally extends to a morphism of species $\MG_c\to \point{\MG}_c$. A natural question to 
ask is if it also possible to find an operad structure on $\point{\MG}_c$ in order to make the following 
commutative diagram of species a commutative diagram of operads:
\begin{equation}\label{diagram_psi}
    \begin{tikzcd}[column sep=huge, ampersand replacement=\&]
        \T \arrow[r, hook, "\psi"] \arrow[d, hook] \& \PLie \arrow[d, hook] \\
        \MG_c \arrow[r, hook, "\psi"] \& \point{\MG}_c
    \end{tikzcd}.
\end{equation}
But as explained before, there does not seem to be a natural way to extend $\PLie$ to $\point{\MG}_c$ and we 
must rather consider $\point{\MG}_{orc}$, from which $\PLie$ is indeed a sub-operad. By doing this we are now 
faced with the problem of finding a natural way to embed $\point{\MG}_c$ in $\point{\MG}_{orc}$ which would 
make the species $\psi(\MG_c)$ a sub-operad of $\point{\MG}_{orc}$ containing $\PLie$. This would then require 
to find a canonical orientation for each pointed multigraph compatible with $\point{\MG}_{orc}$ operad 
structure, which again does not seems possible.

Fortunately, even though it does not seems possible to make the diagram~\eqref{diagram_psi} a diagram of 
operads, we can obtain a similar result albeit with a more involved diagram. To do this let us first introduce 
three new species. 
For $g\in \MG_c[V]$, $r\in V$, and $t\in \T[V]$ a spanning tree of $g$, we denote by $\ori_{(g,t,r)}$ the 
orientation of $g$ defined by $\ori_{(g,t,r)}(e) = \ori_{t,r}(e)$ if $e\in t$ and $\ori_{(g,t,r)}(e) = (\emptyset, e)$ else. 
This orientation orients the edges in $t$ with the orientation induced by the rooted tree $(t,r)$ and make all the 
remaining ends target ends.

Let us define $\mathcal{I}\subset \Operad\subset\ST$ ($\ST$ standing for ``spanning tree'') three sub-species 
of $\point{\MG_{orc}}$ by
\begin{align}
    \ST[V] &= \K\left\{
        (g, \ori_{(g,t,r)},r) : g\in \MG_c[V], r\in V \text{ and $t$ a
    spanning tree of $g$}\right\}, \\
    \Operad[V] &= \K\left\{\sum_{r\in V} (g, \ori_{(g,t,r)},r) : 
    \begin{array}{l}
    g\in \MG_c[V] \text{ and for each $r$,} \\ 
    \text{$t(r)$ a spanning tree of $g$}
    \end{array}\right\}, \\
    \mathcal{I}[V] &=
    \K\left\{ (g, \ori_{(g,t_1,r)},r)-(g, \ori_{(g,t_2,r)},r) : 
    \begin{array}{l}
    g \in \MG_c[V],r\in V, \text{ and } \\
    \text{$t_1, t_2$ two spanning trees of $g$}
    \end{array}\right\}.
\end{align}

\begin{exam}
    Let $V=\set{a,b,c,d}$. We give example of elements in $\ST[V]$, $\Operad[V]$ and $\mathcal{I}[V]$. 
    We colored the edges of spanning trees in red.

    \begin{equation}
        \begin{tikzpicture}[Centering,scale=1.2]
            \tikzset{every loop/.style={}}
            \node[RootGraph](a)at(0,0){$a$};
            \node[NodeGraph](b)at(0.5,0.5){$b$};
            \node[NodeGraph](c)at(1,0){$c$};
            \node[NodeGraph](d)at(0.5,-0.5){$d$};
            \draw[ArcGraph, DarkRed](b)--(a);
            \draw[ArcGraph, DarkRed](c)--(b);
            \draw[ArcGraph](c)edge[bend left=20](d);
            \draw[ArcGraph](d)edge[bend right=20](c);
            \draw[ArcGraph, DarkRed](d)edge[bend left=20](c);
            \draw[ArcGraph](a)--(d);
            \draw[ArcGraph](d)--(a);
            \draw[ArcGraph](b)--(d);
            \draw[ArcGraph](d)--(b);
        \end{tikzpicture}
        \enspace \in \ST[V]
    \end{equation}
    \begin{equation}
        \begin{tikzpicture}[Centering,scale=1.2]
            \tikzset{every loop/.style={}}
            \node[RootGraph](a)at(0,0){$a$};
            \node[NodeGraph](b)at(0.5,0.5){$b$};
            \node[NodeGraph](c)at(1,0){$c$};
            \node[NodeGraph](d)at(0.5,-0.5){$d$};
            \draw[EdgeGraph, DarkRed, Stealth-](a)--(b);
            \draw[EdgeGraph, DarkRed, Stealth-](b)--(c);
            \draw[EdgeGraph](c)edge[bend left=20, Stealth-Stealth](d);
            \draw[EdgeGraph, DarkRed, Stealth-](c)edge[bend right=20](d);
            \draw[EdgeGraph, Stealth-Stealth](a)--(d);
            \draw[EdgeGraph, Stealth-Stealth](b)--(d);
        \end{tikzpicture}
        \enspace + \enspace
        \begin{tikzpicture}[Centering,scale=1.2]
            \tikzset{every loop/.style={}}
            \node[NodeGraph](a)at(0,0){$a$};
            \node[RootGraph](b)at(0.5,0.5){$b$};
            \node[NodeGraph](c)at(1,0){$c$};
            \node[NodeGraph](d)at(0.5,-0.5){$d$};
            \draw[EdgeGraph, DarkRed, -Stealth](a)--(b);
            \draw[EdgeGraph, Stealth-Stealth](b)--(c);
            \draw[EdgeGraph](c)edge[bend left=20, Stealth-Stealth](d);
            \draw[EdgeGraph, DarkRed, -Stealth](c)edge[bend right=20](d);
            \draw[EdgeGraph, Stealth-Stealth](a)--(d);
            \draw[EdgeGraph, DarkRed, Stealth-](b)--(d);
        \end{tikzpicture}
        \enspace + \enspace
        \begin{tikzpicture}[Centering,scale=1.2]
            \tikzset{every loop/.style={}}
            \node[NodeGraph](a)at(0,0){$a$};
            \node[NodeGraph](b)at(0.5,0.5){$b$};
            \node[RootGraph](c)at(1,0){$c$};
            \node[NodeGraph](d)at(0.5,-0.5){$d$};
            \draw[EdgeGraph, Stealth-Stealth](a)--(b);
            \draw[EdgeGraph, DarkRed, -Stealth](b)--(c);
            \draw[EdgeGraph](c)edge[bend left=20, Stealth-Stealth](d);
            \draw[EdgeGraph](c)edge[bend right=20, Stealth-Stealth](d);
            \draw[EdgeGraph, DarkRed, -Stealth](a)--(d);
            \draw[EdgeGraph, DarkRed, Stealth-](b)--(d);
        \end{tikzpicture}
        \enspace + \enspace
        \begin{tikzpicture}[Centering,scale=1.2]
            \tikzset{every loop/.style={}}
            \node[NodeGraph](a)at(0,0){$a$};
            \node[NodeGraph](b)at(0.5,0.5){$b$};
            \node[NodeGraph](c)at(1,0){$c$};
            \node[RootGraph](d)at(0.5,-0.5){$d$};
            \draw[EdgeGraph, DarkRed, -Stealth](a)--(b);
            \draw[EdgeGraph, DarkRed, -Stealth](b)--(c);
            \draw[EdgeGraph, DarkRed, -Stealth](c)edge[bend left=20](d);
            \draw[EdgeGraph](c)edge[bend right=20, Stealth-Stealth](d);
            \draw[EdgeGraph, Stealth-Stealth](a)--(d);
            \draw[EdgeGraph, Stealth-Stealth](b)--(d);
        \end{tikzpicture}
        \enspace \in\Operad[V]
    \end{equation}
    \begin{equation}
        \begin{tikzpicture}[Centering,scale=1.2]
            \tikzset{every loop/.style={}}
            \node[NodeGraph](a)at(0,0){$a$};
            \node[RootGraph](b)at(0.5,0.5){$b$};
            \node[NodeGraph](c)at(1,0){$c$};
            \node[NodeGraph](d)at(0.5,-0.5){$d$};
            \draw[EdgeGraph, DarkRed, -Stealth](a)--(b);
            \draw[EdgeGraph, DarkRed, Stealth-](b)--(c);
            \draw[EdgeGraph](c)edge[bend left=20, Stealth-Stealth](d);
            \draw[EdgeGraph, DarkRed, Stealth-](c)edge[bend right=20](d);
            \draw[EdgeGraph, Stealth-Stealth](a)--(d);
            \draw[EdgeGraph, Stealth-Stealth](b)--(d);
        \end{tikzpicture}
        \enspace - \enspace
        \begin{tikzpicture}[Centering,scale=1.2]
            \tikzset{every loop/.style={}}
            \node[NodeGraph](a)at(0,0){$a$};
            \node[RootGraph](b)at(0.5,0.5){$b$};
            \node[NodeGraph](c)at(1,0){$c$};
            \node[NodeGraph](d)at(0.5,-0.5){$d$};
            \draw[EdgeGraph, DarkRed, -Stealth](a)--(b);
            \draw[EdgeGraph, Stealth-Stealth](b)--(c);
            \draw[EdgeGraph](c)edge[bend left=20, Stealth-Stealth](d);
            \draw[EdgeGraph, DarkRed, -Stealth](c)edge[bend right=20](d);
            \draw[EdgeGraph, Stealth-Stealth](a)--(d);
            \draw[EdgeGraph, DarkRed, Stealth-](b)--(d);
        \end{tikzpicture} \enspace 
        \in \mathcal{I}[V]
    \end{equation}

\end{exam}

\begin{lemma} \label{lemma_op_fond}
    The following properties hold:
    \begin{enumerate}[label=(\roman*)]
        \item $\ST$ is a sub-operad of $\point{\MG}_{orc}$ isomorphic to $\MG\times\PLie$,
        \item $\Operad$ is a sub-operad of $\ST$, 
        \item $\mathcal{I}$ is an ideal of $\Operad$.
    \end{enumerate}
\end{lemma}


\begin{proof}
    Before proving these three items, we first give two equalities which will help us for the two last items.
    Let $U:\MG_{or} \rightarrow \MG$ be the forgetful functor which sends an oriented multigraph on the 
    multigraph obtained by forgetting the orientation. Let $g_1\in \MG_c'[V_1]$ and $g_2\in \MG_c[V_2]$ be 
    two connected multigraphs, $t$ a spanning tree of $g_1$ and for each $v\in V_2$, $t(v)$ a spanning tree 
    of $g_2$. When $\ast$ is the root of the spanning tree $t$, all the ends of $\ast$ are target ends.
    Since the target ends in an oriented multigraph behave the same than the normal ends in a multigraph, 
    the forgetful functor preserves the partial composition. For example we have: 
    \begin{equation}\begin{split}
        &U\times\id\left(\enspace
        \begin{tikzpicture}[Centering,scale=.6]
            \node[RootGraph, minimum size=3mm](s)at(0,0){$\ast$};
            \node[NodeGraph, minimum size=3mm](a)at(1,1){$a$};
            \node[NodeGraph, minimum size=3mm](b)at(1,-1){$b$};
            \draw[ArcGraph](a)--(s);
            \draw[ArcGraph](b)--(s);
        \end{tikzpicture}
        \enspace \comp \enspace
        \begin{tikzpicture}[Centering,scale=.6]
            \node[RootGraph, minimum size=3mm](c)at(0,0){$c$};
            \node[NodeGraph, minimum size=3mm](d)at(1,0){$d$};
            \draw[ArcGraph](d)--(c);
        \end{tikzpicture}
        \enspace\right)\\
        =&  U\times\id\left(\enspace
        \begin{tikzpicture}[Centering,scale=.6]
            \node[NodeGraph, minimum size=3mm](a)at(1,1){$a$};
            \node[NodeGraph, minimum size=3mm](b)at(1,-1){$b$};
            \node[RootGraph, minimum size=3mm](c)at(0,0){$c$};
            \node[NodeGraph, minimum size=3mm](d)at(2,0){$d$};
            \draw[ArcGraph](a)--(c);
            \draw[ArcGraph](b)--(c);
            \draw[ArcGraph](d)--(c);
        \end{tikzpicture}
        \enspace + \enspace
        \begin{tikzpicture}[Centering,scale=.6]
            \node[NodeGraph, minimum size=3mm](a)at(1,1){$a$};
            \node[NodeGraph, minimum size=3mm](b)at(1,-1){$b$};
            \node[RootGraph, minimum size=3mm](c)at(0,0){$c$};
            \node[NodeGraph, minimum size=3mm](d)at(2,0){$d$};
            \draw[ArcGraph](a)--(d);
            \draw[ArcGraph](d)--(c);
            \draw[ArcGraph](b)--(c);
        \end{tikzpicture}
        \enspace + \enspace
        \begin{tikzpicture}[Centering,scale=.6]
            \node[NodeGraph, minimum size=3mm](a)at(1,1){$a$};
            \node[NodeGraph, minimum size=3mm](b)at(1,-1){$b$};
            \node[RootGraph, minimum size=3mm](c)at(0,0){$c$};
            \node[NodeGraph, minimum size=3mm](d)at(2,0){$d$};
            \draw[ArcGraph](a)--(c);
            \draw[ArcGraph](d)--(c);
            \draw[ArcGraph](b)--(d);
        \end{tikzpicture}
        \enspace + \enspace
        \begin{tikzpicture}[Centering,scale=.6]
            \node[NodeGraph, minimum size=3mm](a)at(1,1){$a$};
            \node[NodeGraph, minimum size=3mm](b)at(1,-1){$b$};
            \node[RootGraph, minimum size=3mm](c)at(0,0){$c$};
            \node[NodeGraph, minimum size=3mm](d)at(2,0){$d$};
            \draw[ArcGraph](a)--(d);
            \draw[ArcGraph](d)--(c);
            \draw[ArcGraph](b)--(d);
        \end{tikzpicture}
        \enspace\right)\\
        =&\enspace
        \begin{tikzpicture}[Centering,scale=.6]
            \node[NodeGraph, minimum size=3mm](a)at(1,1){$a$};
            \node[NodeGraph, minimum size=3mm](b)at(1,-1){$b$};
            \node[RootGraph, minimum size=3mm](c)at(0,0){$c$};
            \node[NodeGraph, minimum size=3mm](d)at(2,0){$d$};
            \draw[EdgeGraph](a)--(c);
            \draw[EdgeGraph](b)--(c);
            \draw[EdgeGraph](d)--(c);
        \end{tikzpicture}
        \enspace + \enspace
        \begin{tikzpicture}[Centering,scale=.6]
            \node[NodeGraph, minimum size=3mm](a)at(1,1){$a$};
            \node[NodeGraph, minimum size=3mm](b)at(1,-1){$b$};
            \node[RootGraph, minimum size=3mm](c)at(0,0){$c$};
            \node[NodeGraph, minimum size=3mm](d)at(2,0){$d$};
            \draw[EdgeGraph](a)--(d);
            \draw[EdgeGraph](d)--(c);
            \draw[EdgeGraph](b)--(c);
        \end{tikzpicture}
        \enspace + \enspace
        \begin{tikzpicture}[Centering,scale=.6]
            \node[NodeGraph, minimum size=3mm](a)at(1,1){$a$};
            \node[NodeGraph, minimum size=3mm](b)at(1,-1){$b$};
            \node[RootGraph, minimum size=3mm](c)at(0,0){$c$};
            \node[NodeGraph, minimum size=3mm](d)at(2,0){$d$};
            \draw[EdgeGraph](a)--(c);
            \draw[EdgeGraph](d)--(c);
            \draw[EdgeGraph](b)--(d);
        \end{tikzpicture}
        \enspace + \enspace
        \begin{tikzpicture}[Centering,scale=.6]
            \node[NodeGraph, minimum size=3mm](a)at(1,1){$a$};
            \node[NodeGraph, minimum size=3mm](b)at(1,-1){$b$};
            \node[RootGraph, minimum size=3mm](c)at(0,0){$c$};
            \node[NodeGraph, minimum size=3mm](d)at(2,0){$d$};
            \draw[EdgeGraph](a)--(d);
            \draw[EdgeGraph](d)--(c);
            \draw[EdgeGraph](b)--(d);
        \end{tikzpicture}\\
        =&\enspace 
        \begin{tikzpicture}[Centering,scale=.6]
            \node[RootGraph, minimum size=3mm](s)at(0,0){$\ast$};
            \node[NodeGraph, minimum size=3mm](a)at(1,1){$a$};
            \node[NodeGraph, minimum size=3mm](b)at(1,-1){$b$};
            \draw[EdgeGraph](a)--(s);
            \draw[EdgeGraph](b)--(s);
        \end{tikzpicture}
        \enspace \comp \enspace
        \begin{tikzpicture}[Centering,scale=.6]
            \node[RootGraph, minimum size=3mm](c)at(0,0){$c$};
            \node[NodeGraph, minimum size=3mm](d)at(1,0){$d$};
            \draw[EdgeGraph](d)--(c);
        \end{tikzpicture} 
        \enspace = 
        U\times\id\left(\enspace
        \begin{tikzpicture}[Centering,scale=.6]
            \node[RootGraph, minimum size=3mm](s)at(0,0){$\ast$};
            \node[NodeGraph, minimum size=3mm](a)at(1,1){$a$};
            \node[NodeGraph, minimum size=3mm](b)at(1,-1){$b$};
            \draw[ArcGraph](a)--(s);
            \draw[ArcGraph](b)--(s);
        \end{tikzpicture}
        \enspace\right) \comp U\times\id\left(\enspace
        \begin{tikzpicture}[Centering,scale=.6]
            \node[RootGraph, minimum size=3mm](c)at(0,0){$c$};
            \node[NodeGraph, minimum size=3mm](d)at(1,0){$d$};
            \draw[ArcGraph](d)--(c);
        \end{tikzpicture}
        \enspace\right).
    \end{split}\end{equation}
    More formally, for $r\in V_2$, we have:
    \begin{equation} \begin{split}\label{lemma_op_eq1}
        U\times\id&\left((g_1, \ori_{(g_1,t,\ast)},\ast)\comp (g_2, \ori_{(g_2,t(r),r)},r)\right) \\
        &= \left(g_1|_{V_1} \oplus\bigoplus_{v\in n(\ast)}v\left(\sum V_2\right)\oplus((\sum V_2)^2)^{\oplus g_1(\ast\ast)}\oplus g_2, r\right) \\
        &= (g_1\comp g_2, r).
    \end{split} \end{equation}
    Let now $r$ be a vertex in $V_1$. Denote by $p$ the parent of $\ast$ in the rooted tree $(t,r)$, by 
    $c(\ast)$ its children, by $n_{g_1\setminus t}(\ast)$ the multiset of neighbours of $\ast$ in 
    $g_1\setminus t$ and by $n(\ast)$ the multiset of neighbours of $\ast$ in $g_1$, so that 
    $n(\ast)=n_{g_1\setminus t}(\ast)\cup c(\ast)\cup \{p\}$. We then have
    \begin{align} \label{lemma_op_eq2}
        U\times\id&\left((g_1,\ori_{(g_1,t,r)},r)\comp \sum_{v\in V_2}(g_2,\ori_{(g_2,t(v),v)},v) \right) \\ \nonumber
        &= \sum_{v\in V_2} U\times\id\left((g_1,\ori_{(g_1,t,r)},r)\comp (g_2,\ori_{(g_2,t(v),v)},v) \right) \\ \nonumber
        &= \sum_{v\in V_2}\Bigg(g_1|_{V_1}\oplus pv \oplus \bigoplus_{v\in c(\ast)}v\left(\sum V_2\right) \oplus \\ \nonumber
        & \phantom{= \sum_{v\in V_2}\quad } \bigoplus_{v\in n_{g_1\setminus t}(\ast)}v\left(\sum V_2\right) \oplus \Big(\Big(\sum V_2\Big)^2\Big)^{\oplus g_1(\ast\ast)}\oplus g_2, r\Bigg)\\ \nonumber
        &=\Bigg(g_1|_{V_1}\oplus p\left(\sum V_2\right) \oplus \bigoplus_{v\in c(\ast)}v \left(\sum V_2\right) \oplus \\ \nonumber
        & \phantom{= \quad } \bigoplus_{v\in n_{g_1\setminus t}(\ast)}v\left(\sum V_2\right) \oplus \Big(\Big(\sum V_2\Big)^2\Big)^{\oplus g_1(\ast\ast)}\oplus g_2, r\Bigg) \\ \nonumber
        &= \left(g_1|_{V_1}\oplus \bigoplus_{v\in n(\ast)}\left(\sum V_2\right)\oplus \Big(\Big(\sum V_2\Big)^2\Big)^{\oplus g_1(\ast\ast)}\oplus g_2, r\right) \\ \nonumber
        &= (g_1\comp g_2, r).
        \end{align}

    {\em Proof of (i)} The species morphism from $\MG\times \PLie$ to $\point{\MG_{orc}}$ given by 
    $(g,t,r) \mapsto (g\sqcup t,\ori_{(g\sqcup t,t,r)}, r)$ is an operad morphism and hence its image $\ST$ is a sub-operad 
    of $\point{\MG_{orc}}$.

    \begin{figure}[htbp]
        \begin{equation*}
            \begin{tikzpicture}[Centering,scale=1.3]
                \tikzset{every loop/.style={}}
                \node[NodeGraph, minimum size=3mm](a)at(0,0){$a$};
                \node[NodeGraph, minimum size=3mm](b)at(0.5,0.5){$b$};
                \node[NodeGraph, minimum size=3mm](c)at(1,0){$c$};
                \node[NodeGraph, minimum size=3mm](d)at(0.5,-0.5){$d$};
                \draw[EdgeGraph](a)--(b);
                \draw[EdgeGraph](c)edge[bend left=20](d);
                \draw[EdgeGraph](c)edge[bend right=20](d);
            \end{tikzpicture}
            \enspace \otimes \enspace
            \begin{tikzpicture}[Centering,scale=1.3]
                \tikzset{every loop/.style={}}
                \node[NodeGraph, minimum size=3mm](a)at(0,0){$a$};
                \node[NodeGraph, minimum size=3mm](b)at(0.5,0.5){$b$};
                \node[RootGraph, minimum size=3mm](c)at(1,0){$c$};
                \node[NodeGraph, minimum size=3mm](d)at(0.5,-0.5){$d$};
                \draw[EdgeGraph, DarkRed, -Stealth](b)--(c);
                \draw[EdgeGraph, DarkRed, -Stealth](a)--(d);
                \draw[EdgeGraph, DarkRed, Stealth-](b)--(d);
            \end{tikzpicture}
            \enspace \mapsto \enspace
            \begin{tikzpicture}[Centering,scale=1.3]
                \tikzset{every loop/.style={}}
                \node[NodeGraph, minimum size=3mm](a)at(0,0){$a$};
                \node[NodeGraph, minimum size=3mm](b)at(0.5,0.5){$b$};
                \node[RootGraph, minimum size=3mm](c)at(1,0){$c$};
                \node[NodeGraph, minimum size=3mm](d)at(0.5,-0.5){$d$};
                \draw[ArcGraph](a)--(b);
                \draw[ArcGraph](b)--(a);
                \draw[ArcGraph, DarkRed](b)--(c);
                \draw[ArcGraph](c)edge[bend left=30](d);
                \draw[ArcGraph](c)edge[bend right=30](d);
                \draw[ArcGraph](d)edge[bend left=30](c);
                \draw[ArcGraph](d)edge[bend right=30](c);
                \draw[ArcGraph, DarkRed](a)--(d);
                \draw[ArcGraph, DarkRed](d)--(b);
            \end{tikzpicture}
        \end{equation*}
        \caption{An example of the isomorphism of item \textit{i}.}
    \end{figure}

    {\em Proof of (ii)} Let $V_1$ and $V_2$ be two disjoint sets, $g_1\in \MG'_c[V_1]$ and $g_2\in \MG_c[V_2]$
    be two connected multigraphs and for each $v\in V_1\cup\set{\ast}$, $t(v)$ a spanning tree of $g_1$ and for 
    each $v\in V_2$, $t(v)$ a spanning tree of $g_2$. We have
    \begin{align}
        \sum_{r_1\in V_1\cup\set{\ast}} &(g_1,\ori_{(g_1,t(r_1),r_1)}, r_1) \comp \sum_{r_2\in V_2} (g_2,\ori_{(g_2,t(r_2),r_2)}, r_2) \\ \nonumber
        &= \sum_{r_1\in V_1\cup\set{\ast}}\sum_{r_2\in V_2} (g_1,\ori_{(g_1,t(r_1),r_1)}, r_1)\comp (g_2,\ori_{(g_2,t(r_2),r_2)}, r_2) 
        %&= \sum_{r_1\in V_1\cup\set{\ast}}\sum_{r_2\in V_2} \left(\overrightarrow{g_1}^{(t(r_1),r_1)}|_{\ast_{s}\take r_{2s},\, \ast_{t}\take(\sum V_2)_{t}}
        %\oplus \overrightarrow{g_2}^{(t(r_2),r_2)}, r_1|_{\ast\take r_2}\right).
    \end{align}
    Then from \eqref{lemma_op_eq1} and \eqref{lemma_op_eq2} we know that applying $U\times\id$ to the 
    preceding sum gives us:
    \begin{equation}
        \sum_{r\in V_1+V_2} (g_1\comp g_2,r).
    \end{equation}
    To conclude note that $\Operad[V]$ is the reciprocal image of $\K\{\sum_{v\in V}(g,v)\,|\,g\in \MG_c[V]\}$
    by  $U\times\id: \ST\rightarrow \point{\MG}$.

    {\em Proof of (iii)} It is easy to see that $\mathcal{I}$ is a left ideal of $\ST$ and hence of $\Operad$.
    Let be $g_1\in \MG'_c[V_1]$, $g_2\in \MG_c[V_2]$, $r\in V_1$, $t$ a spanning tree of $g_1$ and for every 
    $v\in V_2$, $t(v)$ a spanning tree of $g_2$. Then from \eqref{lemma_op_eq1} and \eqref{lemma_op_eq2} we 
    know that $U\times\id((g_1,\ori_{(g,t,r)}, r)\comp \sum_{v\in V_2} (g_2,\ori_{(g_2,t(v),v))}, v)$ is of the form
    $(g_1\comp g_2, r)$ if $r\not = \ast$, and of the form $\sum_{v\in V_2} (g_1\comp g_2, v)$ otherwise. In 
    both cases it does not depend on $t$. This concludes this proof since $\mathcal{I}[V]$ is the kernel of 
    $(U\times\id)_V: \ST[V]\rightarrow \MG_c^{\bullet}[V]$.
\end{proof}

We can see $\PLie$ as a sub-operad of $\ST$ by the monomorphism $(t,r)\mapsto (t, \ori_{t,r},r)$. The image of the
operad morphism $\psi$ of Proposition~\ref{prop_prelie} is then $\Operad\cap \PLie$ and we have that 
$\mathcal{I}\cap \PLie = \set{0}$ and hence $\Operad\cap \PLie/\mathcal{I}\cap \PLie = \Operad\cap \PLie$.


\begin{prop}
    The operad isomorphism $\psi: \T \to \PLie\cap\Operad$ extends into an operad isomorphism 
    $\psi: \MG_c \to \Operad/\mathcal{I}$ satisfying,
    for any $g \in \MG_c[V]$,
    \begin{equation}
        \psi(g) = \sum_{r\in V} (g,\ori_{(g,t(r),r)}, r),
    \end{equation}
    where for each $r\in V$, $t(r)$ is a spanning tree of $g$. Furthermore, this isomorphism restricts 
    itself to an isomorphism  $\G_c \to \Operad\cap\point{\G_{orc}}/\mathcal{I}\cap\point{\G_{orc}}$.
\end{prop}

\begin{proof}
    This statement is a direct consequence of Lemma~\ref{lemma_op_fond} and its proof.
\end{proof}

The last results are summarized in the following commutative diagram of operad morphisms.
\begin{equation}
    \begin{tikzcd}
        \T \arrow[r, "\sim"] \arrow[d, hook]
        & \PLie\cap\K\Operad/\mathcal{I} \arrow[r, equal] \arrow[d, hook]
        & \PLie\cap\Operad \arrow[d,hook] \arrow[r, hook]
        & \PLie \arrow[d,hook]\\
        \G_c \arrow[r, "\sim"] \arrow[d, hook]
        & \Operad\cap\point{\G}_{orc}/\mathcal{I}\cap\point{\G_{orc}} \arrow[d, hook]
        & \point{\G}_{orc}\cap\Operad \arrow[l, two heads] \arrow[r, hook] \arrow[d, hook]
        & \point{\G}_{orc}\cap\ST \arrow[d, hook] \\
        \MG_c \arrow[r, "\sim"]
        & \Operad/\mathcal{I}
        & \Operad \arrow[l, two heads] \arrow[r, hook]
        & \MG\times\PLie
    \end{tikzcd}
\end{equation}





\subsection{The \texorpdfstring{$\MG$}{MG} operad}
As shown in the previous subsection, $\MG$ and its sub-operads have some interesting properties and links 
with  $\PLie$. While $\MG$ was defined by explicitly describing $\MG[V]$ for each $V$, the operad $\PLie$ 
was first defined by its generator and the pre-Lie relation and then is was proven that it is an operad on 
rooted tree in \cite{CL01}. It is then natural to search for generators and relations of $\MG$ and its 
sub-operads. 


To describe a generating family of $\G$, we currently have no better method than performing an exhaustive 
search. With the aid of a computer, we therefore construct this generating family iteratively, arity by arity, 
incorporating a simple graph if it cannot be expressed as a linear combination of partial compositions and 
applications of the symmetric group action on elements of the previously computed family.

Through this process, we have computed the following family of simple graphs with fewer than six vertices:
\begin{equation}\begin{split} \label{equ:generators_G}
    &\begin{tikzpicture}[Centering,scale=.6]
        \node[UnlabeledNodeGraph](a)at(0,0){};
        \node[UnlabeledNodeGraph](b)at(1,0){};
    \end{tikzpicture},
    \enspace
    \begin{tikzpicture}[Centering,scale=.6]
        \node[UnlabeledNodeGraph](a)at(0,0){};
        \node[UnlabeledNodeGraph](b)at(1,0){};
        \draw[EdgeGraph](a)--(b);
    \end{tikzpicture},
    \qquad \quad
    \begin{tikzpicture}[Centering,scale=.3]
        \node[UnlabeledNodeGraph](a)at(-1,-1){};
        \node[UnlabeledNodeGraph](b)at(1,-1){};
        \node[UnlabeledNodeGraph](c)at(0,.707){};
        \draw[EdgeGraph](a)--(b);
        \draw[EdgeGraph](a)--(c);
        \draw[EdgeGraph](b)--(c);
    \end{tikzpicture},
    \qquad \quad
    \begin{tikzpicture}[Centering,scale=.4]
        \node[UnlabeledNodeGraph](a)at(0,0){};
        \node[UnlabeledNodeGraph](b)at(1,0){};
        \node[UnlabeledNodeGraph](c)at(2,0){};
        \node[UnlabeledNodeGraph](d)at(3,0){};
        \draw[EdgeGraph](a)--(b);
        \draw[EdgeGraph](b)--(c);
        \draw[EdgeGraph](c)--(d);
    \end{tikzpicture},
    \enspace
    \begin{tikzpicture}[Centering,scale=.4]
        \node[UnlabeledNodeGraph](a)at(-1,-1){};
        \node[UnlabeledNodeGraph](b)at(-1,1){};
        \node[UnlabeledNodeGraph](c)at(0,0){};
        \node[UnlabeledNodeGraph](d)at(-2,0){};
        \draw[EdgeGraph](a)--(b);
        \draw[EdgeGraph](a)--(c);
        \draw[EdgeGraph](a)--(d);
        \draw[EdgeGraph](b)--(d);
    \end{tikzpicture},
    \enspace
    \begin{tikzpicture}[Centering,scale=.4]
        \node[UnlabeledNodeGraph](a)at(-1,-1){};
        \node[UnlabeledNodeGraph](b)at(-1,1){};
        \node[UnlabeledNodeGraph](c)at(0,0){};
        \node[UnlabeledNodeGraph](d)at(-2,0){};
        \draw[EdgeGraph](a)--(b);
        \draw[EdgeGraph](a)--(c);
        \draw[EdgeGraph](a)--(d);
        \draw[EdgeGraph](b)--(c);
        \draw[EdgeGraph](b)--(d);
    \end{tikzpicture},
    \enspace
    \begin{tikzpicture}[Centering,scale=.4]
        \node[UnlabeledNodeGraph](a)at(-1,-1){};
        \node[UnlabeledNodeGraph](b)at(-1,1){};
        \node[UnlabeledNodeGraph](c)at(0,0){};
        \node[UnlabeledNodeGraph](d)at(-2,0){};
        \draw[EdgeGraph](a)--(b);
        \draw[EdgeGraph](a)--(c);
        \draw[EdgeGraph](a)--(d);
        \draw[EdgeGraph](b)--(c);
        \draw[EdgeGraph](b)--(d);
        \draw[EdgeGraph](c)--(d);
    \end{tikzpicture},
    \\
    &
    \begin{tikzpicture}[Centering,scale=.3]
        \node[UnlabeledNodeGraph](a)at(-2,0){};
        \node[UnlabeledNodeGraph](b)at(0,0){};
        \node[UnlabeledNodeGraph](c)at(-2,2){};
        \node[UnlabeledNodeGraph](d)at(0,2){};
        \node[UnlabeledNodeGraph](e)at(-1,-2){};
        \draw[EdgeGraph](a)--(b);
        \draw[EdgeGraph](a)--(c);
        \draw[EdgeGraph](a)--(e);
        \draw[EdgeGraph](b)--(d);
        \draw[EdgeGraph](b)--(e);
    \end{tikzpicture},
    \enspace
    \begin{tikzpicture}[Centering,scale=.3]
        \node[UnlabeledNodeGraph](a)at(-2,0){};
        \node[UnlabeledNodeGraph](b)at(0,-2){};
        \node[UnlabeledNodeGraph](c)at(-2,2){};
        \node[UnlabeledNodeGraph](d)at(0,0){};
        \node[UnlabeledNodeGraph](e)at(-2,-2){};
        \draw[EdgeGraph](a)--(c);
        \draw[EdgeGraph](a)--(d);
        \draw[EdgeGraph](a)--(e);
        \draw[EdgeGraph](b)--(d);
        \draw[EdgeGraph](b)--(e);
    \end{tikzpicture},
    \enspace
    \begin{tikzpicture}[Centering,scale=.3]
        \node[UnlabeledNodeGraph](a)at(-1,-1){};
        \node[UnlabeledNodeGraph](b)at(-1.25,.75){};
        \node[UnlabeledNodeGraph](c)at(0,2){};
        \node[UnlabeledNodeGraph](d)at(1.25,.75){};
        \node[UnlabeledNodeGraph](e)at(1,-1){};
        \draw[EdgeGraph](a)--(b);
        \draw[EdgeGraph](b)--(c);
        \draw[EdgeGraph](c)--(d);
        \draw[EdgeGraph](d)--(e);
        \draw[EdgeGraph](e)--(a);
    \end{tikzpicture},
    \enspace
    \begin{tikzpicture}[Centering,scale=.3]
        \node[UnlabeledNodeGraph](a)at(-2,0){};
        \node[UnlabeledNodeGraph](b)at(0,0){};
        \node[UnlabeledNodeGraph](c)at(-2,2){};
        \node[UnlabeledNodeGraph](d)at(0,2){};
        \node[UnlabeledNodeGraph](e)at(-1,-2){};
        \draw[EdgeGraph](a)--(b);
        \draw[EdgeGraph](a)--(c);
        \draw[EdgeGraph](a)--(d);
        \draw[EdgeGraph](a)--(e);
        \draw[EdgeGraph](b)--(d);
        \draw[EdgeGraph](b)--(e);
    \end{tikzpicture},
    \enspace
    \begin{tikzpicture}[Centering,scale=.4]
        \node[UnlabeledNodeGraph](a)at(0,0){};
        \node[UnlabeledNodeGraph](b)at(1,-1){};
        \node[UnlabeledNodeGraph](c)at(-1,1){};
        \node[UnlabeledNodeGraph](d)at(1,1){};
        \node[UnlabeledNodeGraph](e)at(-1,-1){};
        \draw[EdgeGraph](a)--(b);
        \draw[EdgeGraph](a)--(c);
        \draw[EdgeGraph](a)--(d);
        \draw[EdgeGraph](a)--(e);
        \draw[EdgeGraph](b)--(d);
        \draw[EdgeGraph](c)--(e);
    \end{tikzpicture},
    \enspace
    \begin{tikzpicture}[Centering,scale=.4]
        \node[UnlabeledNodeGraph](a)at(0,0){};
        \node[UnlabeledNodeGraph](b)at(1,-1){};
        \node[UnlabeledNodeGraph](c)at(-1,1){};
        \node[UnlabeledNodeGraph](d)at(1,1){};
        \node[UnlabeledNodeGraph](e)at(-1,-1){};
        \draw[EdgeGraph](a)--(b);
        \draw[EdgeGraph](a)--(d);
        \draw[EdgeGraph](a)--(e);
        \draw[EdgeGraph](b)--(d);
        \draw[EdgeGraph](b)--(e);
        \draw[EdgeGraph](c)--(e);
    \end{tikzpicture},
    \begin{tikzpicture}[Centering,scale=.4]
        \node[UnlabeledNodeGraph](a)at(0,0){};
        \node[UnlabeledNodeGraph](b)at(1,-1){};
        \node[UnlabeledNodeGraph](c)at(-1,1){};
        \node[UnlabeledNodeGraph](d)at(1,1){};
        \node[UnlabeledNodeGraph](e)at(-1,-1){};
        \draw[EdgeGraph](a)--(c);
        \draw[EdgeGraph](a)--(d);
        \draw[EdgeGraph](a)--(e);
        \draw[EdgeGraph](b)--(d);
        \draw[EdgeGraph](b)--(e);
        \draw[EdgeGraph](c)--(e);
    \end{tikzpicture},
    \enspace
    \begin{tikzpicture}[Centering,scale=.3]
        \node[UnlabeledNodeGraph](a)at(0,-1){};
        \node[UnlabeledNodeGraph](b)at(3,-1){};
        \node[UnlabeledNodeGraph](c)at(1.5,2){};
        \node[UnlabeledNodeGraph](d)at(1.5,.5){};
        \node[UnlabeledNodeGraph](e)at(1.5,-2.5){};
        \draw[EdgeGraph](a)--(b);
        \draw[EdgeGraph](a)--(c);
        \draw[EdgeGraph](a)--(d);
        \draw[EdgeGraph](a)--(e);
        \draw[EdgeGraph](b)--(c);
        \draw[EdgeGraph](b)--(d);
        \draw[EdgeGraph](b)--(e);
    \end{tikzpicture},
    \enspace
    \begin{tikzpicture}[Centering,scale=.3]
        \node[UnlabeledNodeGraph](a)at(0,-1.5){};
        \node[UnlabeledNodeGraph](b)at(3,-1){};
        \node[UnlabeledNodeGraph](c)at(3,-2){};
        \node[UnlabeledNodeGraph](d)at(1.5,0){};
        \node[UnlabeledNodeGraph](e)at(1.5,-3){};
        \draw[EdgeGraph](a)--(b);
        \draw[EdgeGraph](a)--(c);
        \draw[EdgeGraph](a)--(d);
        \draw[EdgeGraph](a)--(e);
        \draw[EdgeGraph](b)--(c);
        \draw[EdgeGraph](b)--(d);
        \draw[EdgeGraph](c)--(e);
    \end{tikzpicture},
    \enspace
    \begin{tikzpicture}[Centering,scale=.6]
        \node[UnlabeledNodeGraph](a)at(0,0){};
        \node[UnlabeledNodeGraph](b)at(1,1){};
        \node[UnlabeledNodeGraph](c)at(.5,-1){};
        \node[UnlabeledNodeGraph](d)at(0,1){};
        \node[UnlabeledNodeGraph](e)at(1,0){};
        \draw[EdgeGraph](a)--(b);
        \draw[EdgeGraph](a)--(c);
        \draw[EdgeGraph](a)--(d);
        \draw[EdgeGraph](a)--(e);
        \draw[EdgeGraph](b)--(d);
        \draw[EdgeGraph](b)--(e);
        \draw[EdgeGraph](c)--(e);
    \end{tikzpicture},
    \\
    &
    \quad
    \begin{tikzpicture}[Centering,scale=.6]
        \node[UnlabeledNodeGraph](a)at(0,0){};
        \node[UnlabeledNodeGraph](b)at(1,1){};
        \node[UnlabeledNodeGraph](c)at(2,-1){};
        \node[UnlabeledNodeGraph](d)at(0,1){};
        \node[UnlabeledNodeGraph](e)at(1,0){};
        \draw[EdgeGraph](a)--(c);
        \draw[EdgeGraph](a)--(d);
        \draw[EdgeGraph](a)--(e);
        \draw[EdgeGraph](b)--(c);
        \draw[EdgeGraph](b)--(d);
        \draw[EdgeGraph](b)--(e);
        \draw[EdgeGraph](c)--(e);
    \end{tikzpicture},
    \enspace
    \begin{tikzpicture}[Centering,scale=.6]
        \node[UnlabeledNodeGraph](a)at(0,0){};
        \node[UnlabeledNodeGraph](b)at(1,1){};
        \node[UnlabeledNodeGraph](c)at(2,-1){};
        \node[UnlabeledNodeGraph](d)at(0,1){};
        \node[UnlabeledNodeGraph](e)at(1,0){};
        \draw[EdgeGraph](a)--(b);
        \draw[EdgeGraph](a)--(c);
        \draw[EdgeGraph](a)--(d);
        \draw[EdgeGraph](a)--(e);
        \draw[EdgeGraph](b)--(c);
        \draw[EdgeGraph](b)--(d);
        \draw[EdgeGraph](b)--(e);
        \draw[EdgeGraph](c)--(e);
    \end{tikzpicture},
    \enspace
    \begin{tikzpicture}[Centering,scale=.6,rotate=90]
        \node[UnlabeledNodeGraph](a)at(0,0){};
        \node[UnlabeledNodeGraph](b)at(1,1){};
        \node[UnlabeledNodeGraph](c)at(-1,-1){};
        \node[UnlabeledNodeGraph](d)at(0,1){};
        \node[UnlabeledNodeGraph](e)at(1,0){};
        \draw[EdgeGraph](a)--(b);
        \draw[EdgeGraph](a)--(c);
        \draw[EdgeGraph](a)--(d);
        \draw[EdgeGraph](a)--(e);
        \draw[EdgeGraph](b)--(d);
        \draw[EdgeGraph](b)--(e);
        \draw[EdgeGraph](c)--(d);
        \draw[EdgeGraph](c)--(e);
    \end{tikzpicture},
    \enspace
    \begin{tikzpicture}[Centering,scale=.6,rotate=90]
        \node[UnlabeledNodeGraph](a)at(0,0){};
        \node[UnlabeledNodeGraph](b)at(1,1){};
        \node[UnlabeledNodeGraph](c)at(-1,-1){};
        \node[UnlabeledNodeGraph](d)at(0,1){};
        \node[UnlabeledNodeGraph](e)at(1,0){};
        \draw[EdgeGraph](a)--(b);
        \draw[EdgeGraph](a)--(c);
        \draw[EdgeGraph](a)--(d);
        \draw[EdgeGraph](a)--(e);
        \draw[EdgeGraph](b)--(d);
        \draw[EdgeGraph](b)--(e);
        \draw[EdgeGraph](c)--(d);
        \draw[EdgeGraph](c)--(e);
        \draw[EdgeGraph](d)--(e);
    \end{tikzpicture},
    \enspace
    \begin{tikzpicture}[Centering,scale=.35]
        \node[UnlabeledNodeGraph](a)at(-1,-1){};
        \node[UnlabeledNodeGraph](b)at(-1.25,.75){};
        \node[UnlabeledNodeGraph](c)at(0,2){};
        \node[UnlabeledNodeGraph](d)at(1.25,.75){};
        \node[UnlabeledNodeGraph](e)at(1,-1){};
        \draw[EdgeGraph](a)--(b);
        \draw[EdgeGraph](a)--(c);
        \draw[EdgeGraph](a)--(d);
        \draw[EdgeGraph](a)--(e);
        \draw[EdgeGraph](b)--(c);
        \draw[EdgeGraph](b)--(d);
        \draw[EdgeGraph](b)--(e);
        \draw[EdgeGraph](c)--(d);
        \draw[EdgeGraph](c)--(e);
        \draw[EdgeGraph](d)--(e);
    \end{tikzpicture}.
\end{split}\end{equation}
Due to the symmetric group action on $\G$, only the knowledge of the shapes of the graphs is significant. 
While~\eqref{equ:generators_G} does not provide to us any particular insight on a possible characterisation 
of the generators, it does suggest that any graph with ``enough'' edges must be a generator. This is confirmed
by the following lemma.

\begin{lemma} \label{lem:infinite_number_generators_G}
    Let $\spe{S}$ be a sub-species of $\G$ and let $g$ be a graph in $\G[V]$ with at least $\binom{n-1}{2} +2$ 
    edges, where $n = \card{V}$. Then $g$ belongs to the sub-operad generated by $\spe{S}$ if and only if $g\in\spe{S}[V]$.
\end{lemma}

\begin{proof}
    Suppose that $g\not\in\spe{S}[V]$. It is sufficient to show that $g$ cannot appear in the support of any 
    vector of the form $g_1\comp g_2$ for any $g_1$ and $g_2$ different of $g$. Hence let $g_1\in \G'[V_1]$ 
    and $g_2\in \G[V_2]$ be two graphs, and denote by $e_1$ the number of edges of $g_1$ and by $e_2$ the 
    number of edges of $g_2$. Then the graphs in the sum $g_1\comp g_2$ have $e_1+e_2$ edges. This is maximal
    when $g_1$ and $g_2$ are both complete graphs and is then equal to 
    $\binom{x+1}{2}+\binom{n-x}{2} = x^2-(n-1)x+\binom{n}{2}$ where $0\leq x=\card{V_1}\leq n-1$.

    If $x = 0$ then necessarily $g_1 =\emptyset_{\ast}$ and $g_1\comp g_2 = g_2$ and $g$ appears in the sum 
    $g_1\comp g_2$ if and only if $g=g_2$. This is impossible, hence $x\not = 0$. Similarly we have 
    $x\not = n-1$. The expression $x^2-(n-1)x+\binom{n}{2}$ is then maximal for $x=1$ or $x=n-2$ and is 
    equal in both cases to $\binom{n-1}{2}+1<\binom{n-1}{2} +2$. This implies that $g$ can not be part of the 
    sum of graphs $g_1\comp g_2$ and concludes the proof.
\end{proof}

\begin{prop} \label{gc} 
    The operad $\G$ is not free and has an infinite number of generators.
\end{prop}

\begin{proof}
    The fact that $\G$ has an infinite number of generators is a direct consequence of 
    Lemma~\ref{lem:infinite_number_generators_G}. Moreover, the relation
    \begin{equation}\begin{split} \label{nf}
        \begin{tikzpicture}[Centering,scale=.7]
            \node[NodeGraph](a)at(0,0){$a$};
            \node[NodeGraph](s)at(1,0){$\ast$};
            \draw[EdgeGraph](a)--(s);
        \end{tikzpicture}
        &
        \comp
        \begin{tikzpicture}[Centering,scale=.7]
            \node[NodeGraph](b)at(0,0){$b$};
            \node[NodeGraph](c)at(1,0){$c$};
            \draw[EdgeGraph](b)--(c);
        \end{tikzpicture}
        \enspace + \enspace
        \begin{tikzpicture}[Centering,scale=.7]
            \node[NodeGraph](c)at(0,0){$c$};
            \node[NodeGraph](s)at(1,0){$\ast$};
            \draw[EdgeGraph](c)--(s);
        \end{tikzpicture}
        \comp
        \begin{tikzpicture}[Centering,scale=.7]
            \node[NodeGraph](b)at(0,0){$b$};
            \node[NodeGraph](a)at(1,0){$a$};
            \draw[EdgeGraph](b)--(a);
        \end{tikzpicture}
        \enspace - \enspace
        \begin{tikzpicture}[Centering,scale=.7]
            \node[NodeGraph](b)at(0,0){$b$};
            \node[NodeGraph](s)at(1,0){$\ast$};
            \draw[EdgeGraph](b)--(s);
        \end{tikzpicture}
        \comp
        \begin{tikzpicture}[Centering,scale=.7]
            \node[NodeGraph](a)at(0,0){$a$};
            \node[NodeGraph](c)at(1,0){$c$};
            \draw[EdgeGraph](a)--(c);
        \end{tikzpicture}
        \enspace
        - 2 \,
        \begin{tikzpicture}[Centering,scale=.7]
            \node[NodeGraph](a)at(0,0){$a$};
            \node[NodeGraph](b)at(1,0){$b$};
            \node[NodeGraph](c)at(2,0){$c$};
            \draw[EdgeGraph](a)--(b);
            \draw[EdgeGraph](b)--(c);
        \end{tikzpicture}
        \\[.5em]
        & =
        \begin{tikzpicture}[Centering,scale=.7]
            \node[NodeGraph](a)at(0,0){$a$};
            \node[NodeGraph](b)at(1,0){$b$};
            \node[NodeGraph](c)at(2,0){$c$};
            \draw[EdgeGraph](a)--(b);
            \draw[EdgeGraph](b)--(c);
        \end{tikzpicture}
        \enspace + \enspace
        \begin{tikzpicture}[Centering,scale=.7]
            \node[NodeGraph](b)at(0,0){$b$};
            \node[NodeGraph](c)at(1,0){$c$};
            \node[NodeGraph](a)at(2,0){$a$};
            \draw[EdgeGraph](b)--(c);
            \draw[EdgeGraph](c)--(a);
        \end{tikzpicture}
        \enspace + \enspace
        \begin{tikzpicture}[Centering,scale=.7]
            \node[NodeGraph](c)at(0,0){$c$};
            \node[NodeGraph](b)at(1,0){$b$};
            \node[NodeGraph](a)at(2,0){$a$};
            \draw[EdgeGraph](c)--(b);
            \draw[EdgeGraph](b)--(a);
        \end{tikzpicture}
        \enspace + \enspace
        \begin{tikzpicture}[Centering,scale=.7]
            \node[NodeGraph](b)at(0,0){$b$};
            \node[NodeGraph](a)at(1,0){$a$};
            \node[NodeGraph](c)at(2,0){$c$};
            \draw[EdgeGraph](b)--(a);
            \draw[EdgeGraph](a)--(c);
        \end{tikzpicture}
        \\[.5em]
        & \qquad - \enspace
        \begin{tikzpicture}[Centering,scale=.7]
            \node[NodeGraph](b)at(0,0){$b$};
            \node[NodeGraph](a)at(1,0){$a$};
            \node[NodeGraph](c)at(2,0){$c$};
            \draw[EdgeGraph](b)--(a);
            \draw[EdgeGraph](a)--(c);
        \end{tikzpicture}
        \enspace - \enspace
        \begin{tikzpicture}[Centering,scale=.7]
            \node[NodeGraph](a)at(0,0){$a$};
            \node[NodeGraph](c)at(1,0){$c$};
            \node[NodeGraph](b)at(2,0){$b$};
            \draw[EdgeGraph](a)--(c);
            \draw[EdgeGraph](c)--(b);
        \end{tikzpicture}
        - 2 \,
        \begin{tikzpicture}[Centering,scale=.7]
            \node[NodeGraph](a)at(0,0){$a$};
            \node[NodeGraph](b)at(1,0){$b$};
            \node[NodeGraph](c)at(2,0){$c$};
            \draw[EdgeGraph](a)--(b);
            \draw[EdgeGraph](b)--(c);
        \end{tikzpicture}
        \\[.5em]
        & = 0
    \end{split}\end{equation}
    shows that $\G$ is not free.
\end{proof}

As a consequence of Proposition~\ref{gc}, it seems particularly involved to find a definition of $\G$ by 
generators and relations. We did not have any more success in describing $\T$ in such a way even exploiting 
the monomorphism $\psi:\T\to\PLie$. We at least managed to compute a family of generators for trees over less 
than $7$ vertices:
\begin{equation}\begin{split}
    \begin{tikzpicture}[Centering,scale=.6]
        \node[UnlabeledNodeGraph](a)at(0,0){};
        \node[UnlabeledNodeGraph](b)at(1,0){};
        \draw[EdgeGraph](a)--(b);
    \end{tikzpicture},
    \qquad \quad
    \begin{tikzpicture}[Centering,scale=.35]
        \node[UnlabeledNodeGraph](a)at(0,0){};
        \node[UnlabeledNodeGraph](b)at(1,0){};
        \node[UnlabeledNodeGraph](c)at(2,0){};
        \node[UnlabeledNodeGraph](d)at(3,0){};
        \draw[EdgeGraph](a)--(b);
        \draw[EdgeGraph](b)--(c);
        \draw[EdgeGraph](c)--(d);
    \end{tikzpicture},
    \qquad \quad
    \begin{tikzpicture}[Centering,scale=.35]
        \node[UnlabeledNodeGraph](a)at(0,0){};
        \node[UnlabeledNodeGraph](b)at(-1,0){};
        \node[UnlabeledNodeGraph](c)at(0,1){};
        \node[UnlabeledNodeGraph](d)at(1,0){};
        \node[UnlabeledNodeGraph](e)at(-2,0){};
        \draw[EdgeGraph](a)--(b);
        \draw[EdgeGraph](a)--(c);
        \draw[EdgeGraph](a)--(d);
        \draw[EdgeGraph](b)--(e);
    \end{tikzpicture},
    \qquad \quad
    \begin{tikzpicture}[Centering,scale=.35]
        \node[UnlabeledNodeGraph](a)at(0,0){};
        \node[UnlabeledNodeGraph](b)at(-1,0){};
        \node[UnlabeledNodeGraph](c)at(0,1){};
        \node[UnlabeledNodeGraph](d)at(1,0){};
        \node[UnlabeledNodeGraph](e)at(0,-1){};
        \node[UnlabeledNodeGraph](f)at(-2,0){};
        \draw[EdgeGraph](a)--(b);
        \draw[EdgeGraph](a)--(c);
        \draw[EdgeGraph](a)--(d);
        \draw[EdgeGraph](a)--(e);
        \draw[EdgeGraph](b)--(f);
    \end{tikzpicture},
    \enspace
    \begin{tikzpicture}[Centering,scale=.4]
        \node[UnlabeledNodeGraph](a)at(0,0){};
        \node[UnlabeledNodeGraph](b)at(-1,0){};
        \node[UnlabeledNodeGraph](c)at(1,-1){};
        \node[UnlabeledNodeGraph](d)at(1,1){};
        \node[UnlabeledNodeGraph](e)at(-2,-1){};
        \node[UnlabeledNodeGraph](f)at(-2,1){};
        \draw[EdgeGraph](a)--(b);
        \draw[EdgeGraph](a)--(c);
        \draw[EdgeGraph](a)--(d);
        \draw[EdgeGraph](b)--(e);
        \draw[EdgeGraph](b)--(f);
    \end{tikzpicture},
    \enspace
    \begin{tikzpicture}[Centering,scale=.4]
        \node[UnlabeledNodeGraph](a)at(0,0){};
        \node[UnlabeledNodeGraph](b)at(1,0){};
        \node[UnlabeledNodeGraph](c)at(2,0){};
        \node[UnlabeledNodeGraph](d)at(3,0){};
        \node[UnlabeledNodeGraph](e)at(4,0){};
        \node[UnlabeledNodeGraph](f)at(2,1){};
        \draw[EdgeGraph](a)--(b);
        \draw[EdgeGraph](b)--(c);
        \draw[EdgeGraph](c)--(d);
        \draw[EdgeGraph](d)--(e);
        \draw[EdgeGraph](c)--(f);
    \end{tikzpicture}.
\end{split}\end{equation}
As with the previous family~\eqref{equ:generators_G}, this does not give to us any insight on a possible family of generators.




\subsection{Finitely generated sub-operads}
Since finding family of generators seems out of reach, let us now directly focus on finitely generated 
sub-operads of $\MG$. In particular we will study the operads generated by:
\begin{enumerate}
\item 
\begin{math}
\left\{
\begin{tikzpicture}[Centering,scale=.6]
    \node[NodeGraph](a)at(0,0){$a$};
    \node[NodeGraph](b)at(1,0){$b$};
\end{tikzpicture}
\right\}
\end{math} which we denote by $\G_{\emptyset}$ and which is isomorphic to $\Com$,
\item \begin{math}\left\{
\begin{tikzpicture}[Centering,scale=.6]
    \node[NodeGraph](a)at(0,0){$a$};
    \node[NodeGraph](b)at(1,0){$b$};
    \draw[EdgeGraph](a)--(b);
\end{tikzpicture}
\right\}
\end{math} which we denote by $\Seg$ and which is isomorphic to $\ComMag$,
\item 
\begin{math}
\left\{
\begin{tikzpicture}[Centering,scale=.6]
    \node[NodeGraph](a)at(0,0){$a$};
    \node[NodeGraph](b)at(1,0){$b$};
\end{tikzpicture},
\begin{tikzpicture}[Centering,scale=.6]
    \node[NodeGraph](a)at(0,0){$a$};
    \node[NodeGraph](b)at(1,0){$b$};
    \draw[EdgeGraph](a)--(b);
\end{tikzpicture}
\right\}
\end{math} which we denote by $\SP$,
\item \begin{math}
\left\{
\begin{tikzpicture}[Centering,scale=.6]
    \tikzset{every loop/.style={}}
    \node[NodeGraph](a)at(0,0){$a$};
    \draw[EdgeGraph](a)edge[loop](a);
\end{tikzpicture},
\begin{tikzpicture}[Centering,scale=.6]
    \node[NodeGraph](a)at(0,0){$a$};
    \node[NodeGraph](b)at(1,0){$b$};
\end{tikzpicture}
\right\}
\end{math} which we denote by $\LP$.
\end{enumerate}

First, note that the sub-operad $\G_{\emptyset}$ generated by
\begin{math}
\left\{
\begin{tikzpicture}[Centering,scale=.6]
    \node[NodeGraph](a)at(0,0){$a$};
    \node[NodeGraph](b)at(1,0){$b$};
\end{tikzpicture}
\right\}
\end{math}
is isomorphic to the commutative operad $\Com$. Indeed, recall from Example~\ref{ex_gen_com} that $\Com$ is 
the quotient of the free operad over one symmetric element of size two $\Sym_2$ by the associativity relation.
By definition $\G_{\emptyset}$ is also generated by one symmetric element of size $2$, and furthermore we have:
\begin{equation}\label{eq_iso_com}
    \begin{tikzpicture}[Centering,scale=.7]
        \node[NodeGraph](a)at(0,0){$a$};
        \node[NodeGraph](s)at(1,0){$\ast$};
    \end{tikzpicture}
    \enspace \comp \enspace
    \begin{tikzpicture}[Centering,scale=.7]
        \node[NodeGraph](b)at(0,0){$b$};
        \node[NodeGraph](c)at(1,0){$c$};
    \end{tikzpicture}
    \enspace = \enspace
    \begin{tikzpicture}[Centering,scale=.7]
        \node[NodeGraph](a)at(0,0){$a$};
        \node[NodeGraph](b)at(1,0){$b$};
        \node[NodeGraph](c)at(2,0){$c$};
    \end{tikzpicture}
    \enspace = \enspace
    \begin{tikzpicture}[Centering,scale=.7]
        \node[NodeGraph](s)at(0,0){$\ast$};
        \node[NodeGraph](c)at(1,0){$c$};
    \end{tikzpicture}
    \enspace \comp \enspace
    \begin{tikzpicture}[Centering,scale=.7]
        \node[NodeGraph](a)at(0,0){$a$};
        \node[NodeGraph](b)at(1,0){$b$};
    \end{tikzpicture},
\end{equation}
which is the associativity relation. Hence $\G_{\emptyset}\cong\Com$. This could also be observed from the 
fact that we clearly have $\G_{\emptyset}[V]=\K\emptyset_V$ and hence the map $\emptyset_V\mapsto\mu_V$ 
implies an isomorphism from $\G_{\emptyset}$ to $\Com$. 

The other three cases are more involved.
\begin{prop} \label{prop_commag}
    The sub-operad $\Seg$ of $\G$ generated by
    \begin{math}
        \left\{
        \begin{tikzpicture}[Centering,scale=.6]
            \node[NodeGraph](a)at(0,0){$a$};
            \node[NodeGraph](b)at(1,0){$b$};
            \draw[EdgeGraph](a)--(b);
        \end{tikzpicture}
        \right\}
    \end{math}
    is isomorphic to $\ComMag$.
\end{prop}

\begin{proof}
    We know from Proposition~\ref{prop_prelie} that $\Seg$ is isomorphic to the sub-operad of $\PLie$ 
    generated by 
    \begin{equation}
        \left\{
        \begin{tikzpicture}[Centering,scale=.6]
            \node[RootGraph](a)at(0,0){$a$};
            \node[NodeGraph](b)at(0,-1){$b$};
            \draw[EdgeGraph](a)--(b);
        \end{tikzpicture}
        +
        \begin{tikzpicture}[Centering,scale=.6]
            \node[RootGraph](a)at(0,0){$b$};
            \node[NodeGraph](b)at(0,-1){$a$};
            \draw[EdgeGraph](a)--(b);
        \end{tikzpicture}
        \right\}
    \end{equation}
    Then~\cite{BL11} gives us that the map which sends the above element to $s_{\set{a,b}}\in\Sym_2[\set{a,b}]$
    induces an isomorphism between the sub-operad and $\ComMag$. This concludes the proof.
\end{proof}

Remark that while $\ComMag\cong\Seg$, the image of an element in the canonical basis of $\ComMag[V]$, i.e. a 
binary tree with $V$ as set of leaves, is not a single graph but a sum of trees with vertex set $V$.

This result suggests the following more general conjecture.
\begin{conjecture}
    The sub-operad of $\G$ generated by the complete graphs $K_V$ is isomorphic to $\free{\Sym_n}$, the free
    operad over one symmetric element of size $n=\card{V}$.
\end{conjecture}
Proving this would require showing that there are no relations involving only complete graphs of size $n$ in 
$\G$ which is highly non-trivial. In fact this was avoided in the proof of Proposition~\ref{prop_commag} by 
using the results of \cite{BL11} which is somewhat equivalent in the case of $n=2$.
    
The isomorphisms $\Com\cong\G_{\emptyset}$ and $\ComMag\cong\Seg$ allow us to see $\Com$ and $\ComMag$ as 
disjoint sub-operads of $\G$ which gives us a natural way to define the smallest operad containing these two 
as disjoint sub-operads. Denote by $\mathcal{G}$ the sub-species of $\G$ generated by
\begin{math}
    \left\{
    \begin{tikzpicture}[Centering,scale=.6]
        \node[NodeGraph](a)at(0,0){$a$};
        \node[NodeGraph](b)at(1,0){$b$};
    \end{tikzpicture},
    \begin{tikzpicture}[Centering,scale=.6]
        \node[NodeGraph](a)at(0,0){$a$};
        \node[NodeGraph](b)at(1,0){$b$};
        \draw[EdgeGraph](a)--(b);
    \end{tikzpicture}
    \right\}
\end{math}
and $\SP$ the sub-operad of $\mathcal{G}$ generated by these two elements. This operad has some interesting 
properties. Recall from Remark~\ref{sec_operads} that we use the notation $\comp^{\xi}$ for the grafting of 
tree in a free operad and that we denote the equivalence class of $x$ by $\iso{x}$. 


\begin{prop}\label{prop_sp}
    The three following operads are isomorphic
    \begin{itemize}
        \item $\SP$
        \item $\Ope(\mathcal{G},\mathcal{R})$ where $\mathcal{R}$ is the subspecies of $\free{\mathcal{G}}$ generated by
        \begin{subequations}
            \begin{equation} \label{equ:rel_1}
                \Points{c}{\ast} \comp^{\xi} \Points{a}{b}
                \enspace - \enspace
                \Points{a}{\ast} \comp^{\xi} \Points{b}{c},
            \end{equation}
            \begin{center}
                and
            \end{center}
            \begin{equation} \label{equ:rel_2}
                \Segment{a}{\ast} \comp^{\xi} \Points{b}{c}
                \enspace - \enspace
                \Points{c}{\ast} \comp^{\xi} \Segment{a}{b}
                \enspace - \enspace
                \Points{b}{\ast} \comp^{\xi} \Segment{a}{c}.
            \end{equation}
        \end{subequations}
        \item $\Com(\ComMag)$, the assemblies of $\ComMag$.
    \end{itemize}
    In particular, $\SP$ is binary and quadratic.
\end{prop}

\begin{figure}[htbp]
    \begin{equation*}
        \left(\left(
        \begin{tikzpicture}[Centering,xscale=0.6,yscale=0.35]
            \node[NodeFree](x)at(0,0){$\mu$};
            \node[NodeFree](y)at(-1,2){$\mu$};
            \node[NodeFree, draw=white](1)at(-2,4){$1$};
            \node[NodeFree, draw=white](2)at(0,4){$2$};
            \node[NodeFree, draw=white](3)at(1,2){$3$};
            \draw[EdgeFree](x)--(y);
            \draw[EdgeFree](x)--(3);
            \draw[EdgeFree](1)--(y);
            \draw[EdgeFree](2)--(y);
        \end{tikzpicture}
        \enspace \comp[1]^\xi \enspace
        \begin{tikzpicture}[Centering,xscale=0.6,yscale=0.35]
            \node[NodeFree](x)at(0,0){$s$};
            \node[NodeFree, draw=white](a)at(-1,2){$a$};
            \node[NodeFree, draw=white](b)at(1,2){$b$};
            \draw[EdgeFree](x)--(a);
            \draw[EdgeFree](x)--(b);
        \end{tikzpicture}
        \enspace \right) \comp[2]^\xi \enspace
        \begin{tikzpicture}[Centering,scale=0.6]
            \node[NodeFree, draw=white](a)at(0,0){$c$};
        \end{tikzpicture}
        \enspace \right) \comp[3]^\xi \enspace
        \begin{tikzpicture}[Centering,xscale=0.6,yscale=0.35]
            \node[NodeFree](x)at(0,0){$s$};
            \node[NodeFree](y)at(1,2){$s$};
            \node[NodeFree, draw=white](d)at(0,4){$d$};
            \node[NodeFree, draw=white](e)at(2,4){$e$};
            \node[NodeFree, draw=white](f)at(-1,2){$f$};
            \draw[EdgeFree](x)--(y);
            \draw[EdgeFree](x)--(f);
            \draw[EdgeFree](d)--(y);
            \draw[EdgeFree](e)--(y);
        \end{tikzpicture}
        = 
    \end{equation*}
    \begin{equation*}
        \begin{tikzpicture}[Centering,xscale=0.6,yscale=0.35]
            \node[NodeFree](x)at(0,0){$\mu$};
            \node[NodeFree](y)at(-1,2){$\mu$};
            \node[NodeFree, draw=white](2)at(-0.5,4){$c$};
            \node[NodeFree, draw=white](3)at(1,2){$3$};
            \node[NodeFree](1)at(-2,4){$s$};
            \node[NodeFree, draw=white](a)at(-3,6){$a$};
            \node[NodeFree, draw=white](b)at(-1,6){$b$};
            \draw[EdgeFree](1)--(a);
            \draw[EdgeFree](1)--(b);
            \node[NodeFree](3)at(1,2){$s$};
            \node[NodeFree](4)at(2,4){$s$};
            \node[NodeFree, draw=white](d)at(1,6){$d$};
            \node[NodeFree, draw=white](e)at(3,6){$e$};
            \node[NodeFree, draw=white](f)at(0.5,4){$f$};
            \draw[EdgeFree](3)--(4);
            \draw[EdgeFree](3)--(f);
            \draw[EdgeFree](d)--(4);
            \draw[EdgeFree](e)--(4);
            \draw[EdgeFree](x)--(y);
            \draw[EdgeFree](x)--(3);
            \draw[EdgeFree](1)--(y);
            \draw[EdgeFree](2)--(y);
        \end{tikzpicture}
    \end{equation*}
    \caption{An element in the generating family of $\Ope(\mathcal{G}, \mathcal{R})[\set{a,b,c,d,e,f}]$.}
    \label{fig_gen_fam}
\end{figure}

\begin{proof}
    The element $\iso{\Points{a}{b}}$ of $\Ope(\mathcal{G},\mathcal{R})$ is symmetric of size 2 and follows 
    the associativity relation~\eqref{equ:rel_1}. Hence the sub-operad of $\Ope(\mathcal{G},\mathcal{R})$ 
    generated by $\iso{\Points{a}{b}}$ is equal to $\Com$ and all the trees in $\free{\mathcal{G}}$ 
    with $V$ as leaves and whose labels are all empty graphs over two points are sent over $\mu_V$ when 
    passing to the quotient. In the same way, the element $\iso{\Segment{a}{b}}$ of $\Ope(\mathcal{G},
    \mathcal{R})$ is symmetric of size 2 and does not follow any relation involving only itself, hence the 
    sub-operad of $\Ope(\mathcal{G},\mathcal{R})$ generated by $\iso{\Segment{a}{b}}$ is equal to $\ComMag$ 
    and $\iso{\Segment{a}{b}}=s_{\set{a,b}}\in\ComMag[\set{a,b}]$.

    There is a natural epimorphism $\phi$ from $\free{\mathcal{G}}$ to $\SP$ which is the identity on 
    \begin{math}
        \begin{tikzpicture}[Centering,scale=.6]
            \node[NodeGraph](a)at(0,0){$a$};
            \node[NodeGraph](b)at(1,0){$b$};
        \end{tikzpicture}
    \end{math}
    and
    \begin{math}
        \begin{tikzpicture}[Centering,scale=.6]
            \node[NodeGraph](a)at(0,0){$a$};
            \node[NodeGraph](b)at(1,0){$b$};
            \draw[EdgeGraph](a)--(b);
        \end{tikzpicture}
    \end{math}
    and which sends a partial composition $g_1\comp^{\xi} g_2$ on the partial composition $g_1\comp g_2$. 
    We already proved by \eqref{eq_iso_com} that the vector~\eqref{equ:rel_1} is in the kernel $\phi$. The 
    case of \eqref{equ:rel_2} is also straightforward:
    \begin{equation}\begin{split}
        \begin{tikzpicture}[Centering,scale=.7]
            \node[NodeGraph](a)at(0,0){$a$};
            \node[NodeGraph](s)at(1,0){$\ast$};
            \draw[EdgeGraph](a)--(s);
        \end{tikzpicture}
        \enspace \comp \enspace
        \begin{tikzpicture}[Centering,scale=.7]
            \node[NodeGraph](b)at(0,0){$b$};
            \node[NodeGraph](c)at(1,0){$c$};
        \end{tikzpicture}
        \enspace &= \enspace
        \begin{tikzpicture}[Centering,scale=.7]
            \node[NodeGraph](a)at(0,0){$a$};
            \node[NodeGraph](b)at(1,0){$b$};
            \node[NodeGraph](c)at(2,0){$c$};
            \draw[EdgeGraph](a)--(b);
        \end{tikzpicture}
        \enspace + \enspace
        \begin{tikzpicture}[Centering,scale=.7]
            \node[NodeGraph](a)at(0,0){$a$};
            \node[NodeGraph](c)at(1,0){$c$};
            \node[NodeGraph](b)at(2,0){$b$};
            \draw[EdgeGraph](a)--(c);
        \end{tikzpicture}\\
        \enspace &= \enspace
        \begin{tikzpicture}[Centering,scale=.7]
            \node[NodeGraph](c)at(0,0){$c$};
            \node[NodeGraph](s)at(1,0){$\ast$};
        \end{tikzpicture}
        \enspace \comp \enspace
        \begin{tikzpicture}[Centering,scale=.7]
            \node[NodeGraph](a)at(0,0){$a$};
            \node[NodeGraph](b)at(1,0){$b$};
            \draw[EdgeGraph](a)--(b);
        \end{tikzpicture}
        \enspace + \enspace
        \begin{tikzpicture}[Centering,scale=.7]
            \node[NodeGraph](b)at(0,0){$b$};
            \node[NodeGraph](s)at(1,0){$\ast$};
        \end{tikzpicture}
        \enspace \comp \enspace
        \begin{tikzpicture}[Centering,scale=.7]
            \node[NodeGraph](a)at(0,0){$a$};
            \node[NodeGraph](c)at(1,0){$c$};
            \draw[EdgeGraph](a)--(c);
        \end{tikzpicture}.
    \end{split}\end{equation}
    To conclude that $\SP\cong \Ope(\mathcal{G},\mathcal{R})$, we must now show that for any 
    $w\in\Ope(\mathcal{G}, \mathcal{R})[V]$, $\phi(w)=0$ implies $w=0$.
    To do this, we first prove the bijection $\Ope(\mathcal{G},\mathcal{R})\cong \Com(\ComMag)$. Because of 
    \eqref{equ:rel_2}, we have that the vector $s_{\set{a,\ast}}\comp^{\xi}\mu_{\set{b,c}}$ is equal to the 
    vector $\mu_{\set{b,\ast}}\comp^{\xi}s_{\set{a,c}} + \mu_{\set{c,\ast}}\comp^{\xi}s_{\set{a,b}}$ in 
    $\Ope(\mathcal{G}, \mathcal{R})$. Hence, by iterating this process, we get that all elements of 
    $\Ope(\mathcal{G}, \mathcal{R})$ can be written as a sum of equivalence classes of trees where no 
    $s$ vertex has a $\mu$ vertex as descendent. This means that $\Ope(\mathcal{G}, \mathcal{R})[V]$ has the 
    following generating family (cf. Figure~\ref{fig_gen_fam} for an example): 
    \begin{equation}\label{sp_eq_com_commag}
        \set{\mu_{\pi}\comp[]^\xi (t_1, t_2,\dots, t_k)\,|\, \mu_{\pi}\in\Com[\pi],\, t_i\in \ComMag[V_i]}_{\pi = \set{V_1,\dots, V_k}\text{ partition of $V$}},
    \end{equation}
    where $\mu_{\pi}\comp[]^\xi (t_1, t_2,\dots, t_k)$ stands for 
    $(\dots((\mu_{\pi}\comp[V_1]^{\xi} t_1)\comp[V_2]^{\xi} t_2)\dots)\comp[V_k]^{\xi}t_k$.
    
    To conclude that $\Ope(\mathcal{G},\mathcal{R})\cong \Com(\ComMag)$ just remark that the partial composition 
    of $\Ope(\mathcal{G},\mathcal{R})$ acts the same way than the partial composition over assemblies of $\ComMag$.

    Let now be $w$ of the form $\sum_{i=1}^{l}a_iw_i$ where for each $1\leq i\leq l$, $a_i\in \K$ and there 
    is a partition $\pi_i=\set{V_{i,1},\dots, V_{i,k_i}}$ of $V$ such that 
    $w_i=\mu_{\pi_i}\comp[ ]^\xi (t_{i,1},t_{i,2},\dots, t_{i,k_i})$ with $t_{i,j}\in\ComMag[V_{i,j}]$.
    For $i\in[l]$, the image of $\mu_{\pi_i}$ by $\phi$ is the empty graph over $\pi_i$, and so the image 
    of $w_i$ is equal to: 
    \begin{math}
        \bigsqcup_{j=1}^{k_i}\phi(t_{i,j}).
    \end{math}
    Hence, for $i\not=j$ two indices, if $\pi_i\not= \pi_j$ the supports of $\phi(w_i)$ and $\phi(w_j)$ are 
    disjoint. We can then restrict ourselves to the case where all the $w_i$ are on the same partition of 
    $V$, i.e. $\pi_i=\set{V_1,\dots, V_k}$ for all indices $i$.

    Denote by $\G[V_1,\dots,V_k]$ the vector space $\K\set{g_1\sqcup\dots\sqcup g_k\,|\,g_i\in\G[V_i]}$. Then 
    there is an isomorphism from $\G[V_1,\dots,V_k]$ to $\G[V_1]\otimes\dots\otimes\G[V_k]$ defined by 
    $g_1\sqcup\dots\sqcup g_k \mapsto g_1\otimes\dots\otimes g_k$ which sends $\phi(w)$ on 
    $\sum_{i=1}^l a_i \bigotimes_{j=1}^{k}\phi(t_{i,j})$. By definition, $\phi$ sends the elements $t_{i,j}$ 
    on images of the basis elements of $\ComMag[V_i]\hookrightarrow\G[V_i]$ and hence form a free family in 
    $\G[V_i]$. Hence the tensor products $\bigotimes_{j=1}^{k}\phi(t_{i,j})$ also form a free family of 
    $\G[V_1]\otimes\dots\otimes\G[V_k]$ and so $\phi(w)=0$ if and only if $w=0$. This concludes the proof.
\end{proof}

\begin{rema}
    The elements of $\Com(\ComMag)$ are assemblies of $\ComMag$ which can be interpreted as forest of 
    binary trees.
\end{rema}

From now on we identify $\Points{a}{b}$ and $\Segment{a}{b}$ with their respective image in $\Com$ and 
$\ComMag$: $\mu_{\set{a,b}}$ and $s_{\set{a,b}}$. We now exhibit the Koszul dual of $\SP$. 

\begin{prop}\label{sp_dual}
    The operad $\SP$ admits as Koszul dual the operad $\SP^!$ which is isomorphic to the operad $\Ope((\mathcal{G})^{\vee}, \mathcal{R})$ where 
    $\mathcal{R}$ is the subspecies of $\free{\mathcal{G}^\vee}$ generated by
    \begin{subequations}
        \begin{equation} \label{equ:rel_dual_1}
            \Segment{a}{\ast}^{\vee} \comp^{\xi} \Segment{b}{c}^{\vee},
        \end{equation}
        \begin{equation} \label{equ:rel_dual_2}
            \Points{a}{\ast}^{\vee} \comp^{\xi} \Segment{b}{c}^{\vee}
            \enspace + \enspace
            \Segment{c}{\ast}^{\vee} \comp^{\xi} \Points{a}{b}^{\vee}
            \enspace + \enspace
            \Segment{b}{\ast}^{\vee} \comp^{\xi} \Points{a}{c}^{\vee},
        \end{equation}
        \begin{equation} \label{equ:rel_dual_3}
            \Points{a}{\ast}^{\vee} \comp^{\xi} \Points{b}{c}^{\vee}
            +
            \Points{c}{\ast}^{\vee} \comp^{\xi} \Points{a}{b}^{\vee}
            +
            \Points{b}{\ast}^{\vee} \comp^{\xi} \Points{c}{a}^{\vee}.
        \end{equation}
    \end{subequations}
\end{prop}

\begin{proof}
    Let us respectively denote by $r_1$ and $r_2$ and $r'_1$, $r'_2$, and $r'_3$ the vectors \eqref{equ:rel_1}, 
    \eqref{equ:rel_2}, \eqref{equ:rel_dual_1}, \eqref{equ:rel_dual_2}, and~\eqref{equ:rel_dual_3}. Denote by 
    $\mathcal{I}$ the operad ideal generated by~$r_1$ and~$r_2$. As a vector space, 
    $\mathcal{I}[\set{a,b,c}]$ is then the linear span of the set
    \begin{equation}
        \set{r_1, (ab)\cdot r_1, r_2, (abc) \cdot r_2, (acb) \cdot r_2},
    \end{equation}
    where $\cdot$ is the action of the symmetric group, \textit{e.g} 
    $r_1\cdot (ab) = \free{\mathcal{G}}[(ab)](r_1)$. This space is a sub-space of dimension $5$ of 
    $\free{\mathcal{G}}[\set{a,b,c}]$, which is of dimension $12$. Hence, since as a vector space we have
    \begin{equation}
        \free{\mathcal{G}^\vee}[\set{a,b,c}]
        \cong \free{\mathcal{G}^*}[\set{a,b,c}]\cong
        \free{\mathcal{G}}[\set{a,b,c}],
    \end{equation}
    we conclude that $\mathcal{I}^{\bot}[\set{a,b,c}]$ must be of dimension 7.

    Denote by $\mathcal{J}$ the ideal generated by $r_1'$, $r_2'$ and $r_3'$. As a vector space, 
    $\mathcal{J}[\set{a,b,c}]$ is then the linear span of the set
    \begin{equation}
        \set{ r_1', (ab) \cdot r_1', (ac) \cdot r_1',r_2', (abc)\cdot r_2', (acb)\cdot r_2', r_3'}.
    \end{equation}
    This vector space is of dimension 7. To conclude, we need to show that for any elements 
    $f\in\mathcal{J}[\set{a,b,c}]$ and $x\in\mathcal{I}[\set{a,b,c}]$ we have $\scalar{f}{x}=0$. For every 
    pair $\alpha < \beta$ in $\set{a,b,c,\ast}$ ordered with the alphabetical order ($c<\ast$) denote 
    by $s_{\alpha\beta}^{\vee}$ the dual of $s_{\set{\alpha,\beta}}$ and by $\mu_{\alpha\beta}^{\vee}$ the 
    dual of $\mu_{\set{\alpha,\beta}}$. Among the 21 cases to check, we have for example:
    \begin{equation}\begin{split}
        \scalar{ r_1'}{r_1 } &=
        \scalar{ s_{a\ast}^{\vee}\comp^{\xi} s_{bc}^{\vee}}{\mu_{\set{\ast,c}}\comp^{\xi}\mu_{\set{a,b}} -
        \mu_{\set{a,\ast}}\comp^{\xi} \mu_{\set{b,c}} } \\
        &= \scalar{ s_{a\ast}^{\vee}\comp^{\xi} s_{bc}^{\vee}}{ \mu_{\set{\ast,c}}\comp^{\xi}\mu_{\set{a,b}} } - \scalar{ s_{a\ast}^{\vee}\comp^{\xi} s_{bc}^{\vee}}{\mu_{\set{a,\ast}}\comp^{\xi} \mu_{\set{b,c}}} \\
        &= s_{a\ast}^{\vee}(\mu_{\set{\ast,c}})s_{bc}^{\vee}(\mu_{\set{a,b}}) -
        s_{a\ast}^{\vee}(\mu_{\set{a,\ast}})s_{bc}^{\vee}(\mu_{\set{b,c}}) = 0,
    \end{split}\end{equation}
    and
    \begin{align}
        &\scalar{ (abc)\cdot r_2'}{r_2}  = \\ \nonumber
        &= \Big\langle
        \begin{array}{c}%
            \mu_{b\ast}^{\vee}\comp^{\xi}s_{ca}^{\vee} + s_{a\ast}^{\vee}\comp^{\xi}\mu_{bc}^{\vee} \\
            {}+{} s_{c\ast}^{\vee}\comp^{\xi}\mu_{ab}^{\vee}
        \end{array}
        \Bigm|
        \begin{array}{c}%
            s_{\set{a,\ast}}{\comp} \mu_{\set{b,c}} - \mu_{\set{c,\ast}}{\comp} s_{\set{a,b}} \\
            {}-{} \mu_{\set{b,\ast}}{\comp} s_{\set{c,a}}
        \end{array}%
        \Big\rangle  \\ \nonumber
        &= \mu_{b\ast}^{\vee}(s_{\set{a,\ast}})s_{ca}^{\vee}(\mu_{\set{b,c}}) 
        -\mu_{b\ast}^{\vee}(\mu_{\set{c,\ast}})s_{ca}^{\vee}(s_{\set{a,b}}) 
        -\mu_{b\ast}^{\vee}(\mu_{\set{b,\ast}})s_{ca}^{\vee}(s_{\set{c,a}}) \\ \nonumber
        &+ s_{a\ast}^{\vee}(s_{\set{a,\ast}})\mu_{bc}^{\vee}(\mu_{\set{b,c}}) 
        - s_{a\ast}^{\vee}(\mu_{\set{c,\ast}})\mu_{bc}^{\vee}(s_{\set{a,b}}) 
        - s_{a\ast}^{\vee}(\mu_{\set{b,\ast}})\mu_{bc}^{\vee}(s_{\set{c,a}}) \\ \nonumber
        &+ s_{c\ast}^{\vee}(s_{\set{a,\ast}})\mu_{ab}^{\vee}(\mu_{\set{b,c}}) 
        - s_{c\ast}^{\vee}(\mu_{\set{c,\ast}})\mu_{ab}^{\vee}(s_{\set{a,b}}) 
        - s_{c\ast}^{\vee}(\mu_{\set{b,\ast}})\mu_{ab}^{\vee}(s_{\set{c,a}}) \\ \nonumber
        &= -1 +1 = 0.
    \end{align}
    We leave the verification of the 19 remaining cases as an exercise to the interested reader.
\end{proof}

In order to compute the Hilbert series of $\SP^!$ we need to use identity \eqref{hdual} and hence to prove 
that the operad $\SP$ is Koszul.

\begin{prop}\label{sp_koszul}
    The operad $\SP$ is Koszul.
\end{prop}

\begin{proof}
Let $\mathcal{R}$ be the species defined as in Proposition~\ref{prop_sp} so that 
$\SP\cong\Ope(\K\mathcal{G},\mathcal{R})$. Denote by $\mathfrak{p}(a,b)$ and $\mathfrak{s}(a,b)$ the 
elements of $\mathcal{G}^{\mathscr{F}}[\set{a,b},ab]$. Then the following vectors form a basis 
$\mathcal{B}$ of $\mathcal{R}^{\mathscr{F}}[\set{a,b,c},abc]$:
\begin{align}
    \mathfrak{v}_1 =\mathfrak{p}(\mathfrak{p}(a,b),c) - \mathfrak{p}(a,\mathfrak{p}(b,c))  \enspace &,\enspace 
    \mathfrak{v}_2=\mathfrak{p}(\mathfrak{p}(a,c),b) - \mathfrak{p}(a,\mathfrak{p}(b,c))\label{vec_points}\\
    \mathfrak{v'}_1=\mathfrak{s}(\mathfrak{p}(a,b),c) - &\mathfrak{p}(\mathfrak{s}(a,c),b) - \mathfrak{p}(a,\mathfrak{s}(b,c)) \label{vec_seg1} \\
    \mathfrak{v'}_2=\mathfrak{s}(\mathfrak{p}(a,c),b) - &\mathfrak{p}(\mathfrak{s}(a,b),c) - \mathfrak{p}(a, \mathfrak{s}(b,c))  \label{vec_seg2}\\
    \mathfrak{v'}_3=\mathfrak{s}(a,\mathfrak{p}(b,c)) - &\mathfrak{p}(\mathfrak{s}(a,b),c) - \mathfrak{p}(\mathfrak{s}(a,c),b).\label{vec_seg3}
\end{align}
We need to show that it is a Gröbner bases of $(\mathcal{R}^{\mathscr{F}})$. Let us now consider the 
path-lexicographic ordering presented presented page 27 in Appendix~\ref{sec_koszul} with $s>p$. Then the leading terms 
of $\mathfrak{v}_1$, $\mathfrak{v}_2$, $\mathfrak{v'}_1$, $\mathfrak{v'}_2$ and $\mathfrak{v'}_3$ are 
respectively $\mathfrak{p}(\mathfrak{p}(a,b),c)$, $\mathfrak{p}(\mathfrak{p}(a,c),b)$, 
$\mathfrak{s}(\mathfrak{p}(a,b),c)$, $\mathfrak{s}(\mathfrak{p}(a,c),b)$ and $\mathfrak{s}(a,p(b,c))$. 
We conclude with Proposition~\ref{prop_spol}. Indeed, it is shown in \cite{DK10} that the S-polynomials of 
pairs of elements in $\set{\mathfrak{v}_1,\mathfrak{v}_2}$ are congruent to zero modulo $\mathcal{B}$. We 
show for example that the S-polynomial of $\mathfrak{v}_1$ and $\mathfrak{v'}_1$ corresponding to 
$\mathfrak{c}=\mathfrak{s}(\mathfrak{p}(\mathfrak{p}(a,b),c),d)\in\free{\mathcal{G}}^{sh}[\set{a,b,c,d},abcd]$ is congruent to zero. We have
\begin{equation}\begin{split}
    m_{\mathfrak{c},\text{lt}(\mathfrak{v}_1)}(\mathfrak{v}_1) &= \mathfrak{c}-\mathfrak{s}(\mathfrak{p}(a,\mathfrak{p}(b,c)),d)\\
    m_{\mathfrak{c},\text{lt}(\mathfrak{v'}_1)}(\mathfrak{v'}_1) &= \mathfrak{c} - \mathfrak{p}(\mathfrak{s}(\mathfrak{p}(a,b),d),c)
    -\mathfrak{p}(\mathfrak{p}(a,b),\mathfrak{s}(c,d)),
\end{split}\end{equation}
which gives us
\begin{equation}\label{proof_spol}
    s_{\mathfrak{c}}(\mathfrak{v}_1,\mathfrak{v'}_1) =  \mathfrak{p}(\mathfrak{s}(\mathfrak{p}(a,b),d),c) + \mathfrak{p}(\mathfrak{p}(a,b),\mathfrak{s}(c,d))
    -\mathfrak{s}(\mathfrak{p}(a,\mathfrak{p}(b,c)),d).
\end{equation}
Let us look how each of the terms of $s_{\mathfrak{c}}(\mathfrak{v}_1,\mathfrak{v'}_1)$ reduces 
modulo $\mathcal{B}$:
\begin{equation}\begin{split}
    \mathfrak{p}(\mathfrak{s}(\mathfrak{p}(a,b),d),c)
    &\equiv_{\mathfrak{v'}_1} \mathfrak{p}(\mathfrak{p}(\mathfrak{s}(a,d),b),c) + \mathfrak{p}(\mathfrak{p}(a,\mathfrak{s}(b,d)),c) \\
    &\equiv_{\mathfrak{v}_1} \mathfrak{p}(\mathfrak{s}(a,d),\mathfrak{p}(b,c)) + \mathfrak{p}(a, \mathfrak{p}(\mathfrak{s}(b,d),c), \\
    & \\
    \mathfrak{p}(\mathfrak{p}(a,b),\mathfrak{s}(c,d)) &\equiv_{\mathfrak{v}_1} \mathfrak{p}(a, \mathfrak{p}(b,\mathfrak{s}(c,d)), \\
    & \\
    \mathfrak{s}(\mathfrak{p}(a,\mathfrak{p}(b,c)),d) &\equiv_{\mathfrak{v'}_1} \mathfrak{p}(\mathfrak{s}(a,d),\mathfrak{p}(b,c)) +
    \mathfrak{p}(a,\mathfrak{s}(\mathfrak{p}(b,c),d))\\
    &\equiv_{\mathfrak{v'}_1} \mathfrak{p}(\mathfrak{s}(a,d),\mathfrak{p}(b,c)) + \mathfrak{p}(a,\mathfrak{p}(\mathfrak{s}(b,d),c)) +
    \mathfrak{p}(a,\mathfrak{p}(b,\mathfrak{s}(c,d))).
\end{split}\end{equation}
Putting this together in \eqref{proof_spol} gives us that $s_{\mathfrak{c}}(\mathfrak{v}_1,\mathfrak{v'}_1)$ 
reduces to $0$ modulo $\mathcal{B}$. We leave the verification of the other cases to the interested reader.
\end{proof}

\begin{prop} \label{prop_sp_hilbert}
    The Hilbert series of $\SP^{!}$ is given
    \begin{equation}
        \hilbert{\SP^!}(x) = \dfrac{(1-\log(1-x))^2-1}{2}.
    \end{equation}
\end{prop}

\begin{proof}
    The Hilbert series of $\ComMag$ is $\hilbert{\ComMag}(x) = 1-\sqrt{1-2x}$ 
    hence the Hilbert series of $\SP\cong \Com(\ComMag)$ is $\mathcal{H}_{\SP}(x)=e^{1-\sqrt{1-2x}}-1$, 
    where the $-1$ comes from the fact that we consider positive species. 
    We deduce the Hilbert series of $\SP^!$ from $\hilbert{\SP}$ and the identity \eqref{hdual}.
\end{proof}
The first dimensions $\dim \SP^![[n]]$ for $n\geq 1$ are
\begin{equation}
    1, 2, 5, 17, 74, 394, 2484, 18108, 149904.
\end{equation}
This is sequence~\OEIS{A000774} of~\cite{Slo}. This sequence is in particular linked to some
pattern avoiding signed permutations and mesh patterns.

Before ending this section let us mention the sub-operad $\LP$ of $\MG$ generated by
\begin{equation}
    \left\{
    \begin{tikzpicture}[Centering,scale=.6]
        \tikzset{every loop/.style={}}
        \node[NodeGraph](a)at(0,0){$a$};
        \draw[EdgeGraph](a)edge[loop](a);
    \end{tikzpicture},
    \begin{tikzpicture}[Centering,scale=.6]
        \node[NodeGraph](a)at(0,0){$a$};
        \node[NodeGraph](b)at(1,0){$b$};
    \end{tikzpicture}
    \right\}.
\end{equation}
This operad seems particularly interesting to us since its two generators can be considered as minimal
elements in the sense that a partial composition with the two isolated vertices adds exactly
one vertex and no edge, while a partial composition with the loop adds exactly one edge and
no vertex. A natural question to ask at this point concerns the description of the
multigraphs generated by these two minimal elements.

\begin{prop}
The following properties hold:
\begin{itemize}
\item the operad $\SP$ is a sub-operad of $\LP$;
\item the operad $\LP$ is a strict sub-operad of $\MG$. In particular, the multigraph
\begin{equation}
    \begin{tikzpicture}[Centering,scale=.8]
        \node[NodeGraph](a)at(0,0){$a$};
        \node[NodeGraph](b)at(1,0){$b$};
        \node[NodeGraph](c)at(2,0){$c$};
        \draw[EdgeGraph](a)--(b);
        \draw[EdgeGraph](b)edge[bend left=40](c);
        \draw[EdgeGraph](b)edge[bend right=40](c);
    \end{tikzpicture}
\end{equation}
is in $\MG$ but not in $\LP$.
\end{itemize}
\end{prop}

\begin{proof}
    \begin{itemize}
        \item The following identities show that 
        \begin{math}
        \begin{tikzpicture}[Centering,scale=.6]
            \node[NodeGraph](a)at(0,0){$a$};
            \node[NodeGraph](b)at(1,0){$b$};
            \draw[EdgeGraph](a)--(b);
        \end{tikzpicture}
        \end{math}
        is in $\LP[\set{a,b}]$ and hence that $\SP$ is a sub-operad of $\LP$:
        \begin{equation*}
            \Points{\ast}{b} \enspace \comp
            \begin{tikzpicture}[Centering,scale=.6]
                \tikzset{every loop/.style={}}
                \node[NodeGraph](a)at(0,0){$a$};
                \draw[EdgeGraph](a)edge[loop](a);
            \end{tikzpicture}
            \enspace = \enspace
            \begin{tikzpicture}[Centering,scale=.6]
                \tikzset{every loop/.style={}}
                \node[NodeGraph](a)at(0,0){$a$};
                \node[NodeGraph](b)at(1,0){$b$};
                \draw[EdgeGraph](a)edge[loop](a);
            \end{tikzpicture}
        \end{equation*}
        \begin{equation*}
            \begin{tikzpicture}[Centering,scale=.6]
                \tikzset{every loop/.style={}}
                \node[NodeGraph](a)at(0,0){$\ast$};
                \draw[EdgeGraph](a)edge[loop](a);
            \end{tikzpicture}
            \enspace \comp \enspace
            \Points{a}{b} 
            \enspace - \enspace
            \begin{tikzpicture}[Centering,scale=.6]
                \tikzset{every loop/.style={}}
                \node[NodeGraph](a)at(0,0){$a$};
                \node[NodeGraph](b)at(1,0){$b$};
                \draw[EdgeGraph](a)edge[loop](a);
            \end{tikzpicture}
            \enspace - \enspace
            \begin{tikzpicture}[Centering,scale=.6]
                \tikzset{every loop/.style={}}
                \node[NodeGraph](b)at(0,0){$a$};
                \node[NodeGraph](a)at(1,0){$b$};
                \draw[EdgeGraph](a)edge[loop](a);
            \end{tikzpicture}
            \enspace = \enspace
            2 \enspace \Segment{a}{b}.
        \end{equation*}
        \item Using computer algebra, one generates all vectors in $\LP[\set{a,b,c}]$ with three edges
        and shows that the announced multigraph is not a linear combination of these. \qedhere
    \end{itemize}
\end{proof}







\section{Graph insertion operads}\label{sec_graph_operads}
The goal of this section is to give a general construction of operads on multigraphs where the partial 
composition of two elements $g_1\comp g_2$ is given by
\begin{enumerate}
    \item taking the union of $g_1$ and $g_2$,
    \item removing the vertex $\ast$,
    \item \underline{independently} connecting loose ends with vertices of $g_2$.
\end{enumerate}
What we mean by independently is that the way we connect one end does not depend on how we 
connect the others. Note that we may connect a loose end to more than one vertex and hence increase the number
of edges.




\subsection{Constructions on species and operads}\label{sec_constructions}

We begin by defining new constructions on species and operads. We define three constructions:
the augmentation, the semi-direct product and the maps from a set to an operad.

\begin{defi}
    Let $A$ be a set and $\spe{S}$ be a species. An {\em $A$-augmentation} of $\spe{S}$ is a species
     $\augm{A}{\spe{S}}$ such that $\augm{A}{\spe{S}}[V] \cong \spe{S}[A\times V]$ for every finite set $V$.
\end{defi}

\begin{exam}\label{ex_deform}
    Let $A$ be a set.
    \begin{itemize}
        \item Instead of considering an $A$-augmented multigraph on $V$ as a multigraph on $V\times A$, 
        we consider them as multigraphs on $V$ where the ends are labelled with elements of 
        $A$. In particular, the species of oriented multigraphs $\MG_{or}$ is in bijection with the species 
        of $\set{\mathtt{s,t}}$-augmented multigraphs $\augm{\set{\mathtt{s,t}}}{\MG}$. For $g\in\MG$ and 
        $\ori$ an orientation of $g$, the pair $(g,\ori)$ is sent on augmented the multigraph obtained by respectively 
        labeling by $\mathtt{s}$ and $\mathtt{t}$ its sources and targets ends.
        \item Instead of seeing the elements in $\augm{A}{\Poly}[V]$ as polynomials with set of variables the 
        couples $(v,a)\in V\times A$, we consider them as polynomials with set of variables 
        $\set{v_a\,|\,v\in V,a\in A}$ of elements of $V$ indexed by elements of $A$.
    \end{itemize}
\end{exam}

\begin{rema} \label{rem_bij_graphpol}
    For $\spe{S}$ and $\spe{R}$ any two species and $f:\spe{R}\to\spe{S}$ a morphism, $f$ extends to a 
    morphism between any two $A$-augmentations of $\spe{R}$ and $\spe{S}$ by 
    $\augm{A}{\spe{R}}[V]\cong\spe{R}[A\times V]\stackrel{f}{\to}\spe{S}[A\times V]\cong\augm{A}{\spe{S}}$. 
    In particular the morphism $\MG\to\Poly$ given in subsubsection~\ref{graphpol} extends to a morphism 
    which sends an edge $(u_a,v_b)$ on the monomial $u_av_b$.
\end{rema}

In the following proposition, we give an operad structure to a Hadamard product $\spe{S}\times\Operad$ 
where $\spe{S}$ is a species and $\Operad$ is an operad. 
\begin{prop}\label{prop_semiprod}
    Let $\spe{S}$ be a linear species and $\Operad$ an operad. Let $\varphi$ be a morphism from 
    $\spe{S}'\cdot(\spe{S}\times\Operad)$ to $\spe{S}$ and denote by $x\comp^f y=\varphi(x\otimes y\otimes f)$.
    Suppose that $\varphi$ satisfies the following hypotheses. 
    \begin{description}
        \item[Commutativity.] For $x$ an element of $\spe{S}''$ and $y\otimes f$ and $z\otimes g$ two elements 
        of $\spe{S}\times\Operad$,
        \begin{equation}\label{prod_sd_eq1}
            (x\comp[\ast_1]^f y)\comp[\ast_2]^g z  = (x\comp[\ast_2]^g z)\comp[\ast_1]^f y.
        \end{equation}
        \item[Associativity.] For $x$ an element of $\spe{S}'$, $y\otimes f$ an element of 
        $(\spe{S}\times\Operad)'$ and $z\otimes h$ an element of $\spe{S}\times\Operad$,
        \begin{equation}\label{prod_sd_eq2} 
            (x\comp[\ast_1]^f y)\comp[\ast_2]^g z = x\comp[\ast_1]^{f\comp[\ast_2] g}(y\comp[\ast_2]^g z).
        \end{equation}
        \item[Unity.] There exists a map $e:\spe{X}\rightarrow \spe{S}$ such that 
        \begin{equation}\label{prod_sd_eq3}
            x\comp^{e_{\Operad}(v)}e(v) = \spe{S}[\sigma](x) \quad \text{ and }\quad e(\ast)\comp^f x = x,
        \end{equation} 
        where $e_{\Operad}$ is the unit of $\Operad$ and $\sigma$ is the bijection which sends $\ast$ on 
        $v$ and is the identity on the rest of the set on which $x$ is defined.
    \end{description}
    Then the partial composition $\comp^{\varphi}$ defined by
    \begin{equation}\begin{split}
        \comp^{\varphi}: (\spe{S}\times\Operad)'\cdot \spe{S}\times\Operad &\to \spe{S}\times\Operad \\
        (x\otimes f)\otimes (y\otimes g) &\mapsto x\comp^g y\otimes f\comp g
    \end{split}\end{equation}
    makes $\spe{S}\times\Operad$ an operad with unit $e$. We call this operad the {\em semi-direct product 
    of $\spe{S}$ and $\Operad$ over $\varphi$} and we denote it by $\spe{S}\ltimes_{\varphi}\Operad$.
\end{prop}

\begin{proof}
    We must verify that the three diagrams~\eqref{def_op} commute.
    \begin{itemize} 
        \item Let $V_1,V_2,V_3$ be three disjoint sets and $x\otimes f\in(\spe{S}\times\Operad)''[V_1]$, 
        $y\otimes g\in\spe{S}\times\Operad[V_2]$ and $z\otimes h\in\spe{S}\times\Operad[V_3]$. We then have
        \begin{equation}\begin{split}
            \big((x\otimes f) \comp[\ast_1]^{\varphi} (y\otimes g)\big)\comp[\ast_2]^{\varphi} (z\otimes h) 
            &= \big( (x \comp[\ast_1]^g y) \comp[\ast_2]^h z \big)  \otimes  \big((f\comp[\ast_1]g)\comp[\ast_2]h \big) \\
            &= \big( (x \comp[\ast_2]^h z) \comp[\ast_1]^g y \big)  \otimes  \big((f\comp[\ast_2]h)\comp[\ast_1]g \big) \\
            &= \big( (x\otimes f) \comp[\ast_1]^{\varphi} (z\otimes h) \big) \comp[\ast_2]^{\varphi} (y\otimes g),
        \end{split}\end{equation}
        where the second equality follows from~\eqref{prod_sd_eq1} and the fact that $\Operad$ is an operad.
        \item Let $V_1,V_2,V_3$ be three disjoint sets and $x\otimes f\in(\spe{S}\times\Operad)'[V_1]$, 
        $y\otimes g\in(\spe{S}\times\Operad)'[V_2]$ and $z\otimes h\in\spe{S}\times\Operad[V_3]$. We then have
        \begin{equation}\begin{split}
            \big((x\otimes f) \comp[\ast_1]^{\varphi} (y\otimes g)\big)\comp[\ast_2]^{\varphi} (z\otimes h) 
            &= \big( (x\comp[\ast_1]^g y)\comp[\ast_2]^h z \big)  \otimes  \big( (f\comp[\ast_1]g)\comp[\ast_2]h \big) \\
            &= \big( x\comp[\ast_1]^{g\comp[\ast_1]h}( y\comp[\ast_2]^h z) \big)  \otimes  f\comp[\ast_1](g\comp[\ast_2]h) \big) \\
            &= (x\otimes f)\comp[\ast_1]^{\varphi}\big( (y\otimes g)\comp[\ast_2]^{\varphi}(z\otimes h)\big),
        \end{split}\end{equation} 
        where the second equality follows from~\eqref{prod_sd_eq2} and the fact that $\Operad$ is an operad.
        \item Let be $x\otimes f\in(\spe{S}\times\Operad)'[V]$ and $v\not\in V$. We then have 
        \begin{equation}
            (x\otimes f)\comp^{\varphi} (e(v)\otimes e_{\Operad}(v))= x\comp^{e_{\Operad}(v)}e(v) \otimes f\comp e(v) = \spe{S}[\sigma](x\otimes f)
        \end{equation}
        where $\sigma$ is the bijection which sends $\ast$ to $v$ and is the identity on $V$. The last 
        equality follows from the first equality from \eqref{prod_sd_eq3} and the fact that $\Operad$ is an 
        operad. Let now be $x\otimes f\in\spe{S}\times\Operad[V]$. Then
        \begin{equation}
            (e(\ast)\otimes e_{\Operad}(\ast))\comp (x\otimes f) = e(\ast)\comp^f x \otimes e_{\Operad}(\ast)\comp f = x\otimes f
        \end{equation}
        where the last equality comes from second equality from \eqref{prod_sd_eq3} and the fact that 
        $\Operad$ is an operad. \qedhere
    \end{itemize}
\end{proof}

When it is clear from the context, we do not mention $\varphi$ and just write semi-direct product of 
$\spe{S}$ and $\Operad$ and denote it by $\spe{S}\ltimes \Operad$. In practice, the operad structure 
$\Operad$ of $\spe{S}\ltimes\Operad$ is transparent and we are just interested in what happens on 
$\spe{S}$, that it is to say the ``pseudo partial composition'' $x\comp^g y$.

\begin{exam}
    For $C$ a finite set, let $\spe{C}$ be the trivial species given by $\spe{C}[V]=\K C$ for each set $V$ 
    and let $\mathcal{C}=\spe{X}+\spe{C}_{2+}$. This second species has an operad structure given defined by, 
    for $c_1\in \mathcal{C}'[V_1]$ and $c_2\in \mathcal{C}[V_2]$: $c_1\comp c_2 = c_1$ if $V_1\noemp$ and 
    $\ast\comp c_2 = c_2$ when $V_1=\emptyset$ and $\ast\in\spe{X}[\set{\ast}]$. 

    Let $\mathcal{F}^C = \spe{X} + \mathcal{F}_{2+}^C$ be the species of maps with co-domain $C$: 
    $\mathcal{F}^C[V]=\K\set{f:V\to C}$ for $\card{V} > 1$. We define a semi-direct product structure 
    $\mathcal{F}^C\ltimes_{\varphi} \mathcal{C}$. Suppose that $\card{V_1\cup\set{\ast}},\card{V_2}>1$ and let 
    be $f\in \mathcal{F}^C[V_1\cup\set{\ast}]$ and $g\otimes x\in \mathcal{F}^C\times \mathcal{C}[V_2]$. We 
    then have $f\comp^c g = 0$ if $f(\ast) \not = c$ and 
    $f\comp^c g(v) = \left\{\begin{array}{rl}
    f(v)  & \text{if $v\in V_1$}    \\ 
    g(v)  & \text{if $v\in V_2$}\end{array}\right.$ 
    else. When $V_1=\emptyset$ or $V_2$ is a singleton, the action of $\varphi$ is implied by the 
    unit hypothesis. 
   
    We call this operad the {\em $C$-coloration} operad. When this operad is considered alone, one can see 
    an element $(f,c)\in\mathcal{F}^C\ltimes C[V]$ as a corolla on $V$ with its root colored by $c$ and its 
    leaves $v\in V$ colored by $f(v)$. For example, for $C=\set{{\color{Red} Red}, {\color{Green} Green}, {\color{Blue} Blue}}$,
    $V=\set{a,b,c,d}$ and $f\in\mathcal{F}^C[V]$ defined by $f(a) = {\color{Red} Red}$, $f(b) = {\color{Red} Red}$, $f(c) = {\color{Blue} Blue}$, 
    and $f(d) = {\color{Green} Green}$, the element $(f, {\color{Blue} Blue})$ would be represented by the following corolla: 
    \begin{equation}
        \begin{tikzpicture}[Centering, scale=0.6]
            \node[NodeGraph, draw=Blue, fill=Blue!20](n)at(0,0){$ $};
            \node[NodeGraph, draw=Red, fill=Red!20](a)at(-1,1){$a$};
            \node[NodeGraph, draw=Red, fill=Red!20](b)at(-0.125,1.5){$b$};
            \node[NodeGraph, draw=Blue, fill=Blue!20](s)at(0.875,1.25){$c$};
            \node[NodeGraph, draw=Green, fill=Green!20](c)at(1.5,0.75){$d$};
            \draw[EdgeGraph, black](n)--(a);
            \draw[EdgeGraph, black](n)--(b);
            \draw[EdgeGraph, black](n)--(s);
            \draw[EdgeGraph, black](n)--(c);
        \end{tikzpicture}
    \end{equation}
    The partial composition consists then in grafting two corollas if the 
    root and the leaf on which it must be grafted share the same colors. For instance we have:
    \begin{equation}
        \begin{tikzpicture}[Centering, scale=0.6]
            \node[NodeGraph, draw=Red, fill=Red!20](n)at(0,0){$ $};
            \node[NodeGraph, draw=Red, fill=Red!20](a)at(-1,1){$a$};
            \node[NodeGraph, draw=Green, fill=Green!20](b)at(-0.125,1.5){$b$};
            \node[NodeGraph, draw=Blue, fill=Blue!20](s)at(0.875,1.25){$\ast$};
            \node[NodeGraph, draw=Magenta, fill=Magenta!20](c)at(1.5,0.75){$c$};
            \draw[EdgeGraph, black](n)--(a);
            \draw[EdgeGraph, black](n)--(b);
            \draw[EdgeGraph, black](n)--(s);
            \draw[EdgeGraph, black](n)--(c);
        \end{tikzpicture}
        \enspace \comp \enspace
        \begin{tikzpicture}[Centering, scale=0.6]
            \node[NodeGraph, draw=Red, fill=Red!20](n)at(0,0){$ $};
            \node[NodeGraph, draw=Blue, fill=Blue!20](a)at(-0.5,1){$1$};
            \node[NodeGraph, draw=Magenta, fill=Magenta!20](b)at(0.5,1.25){$2$};
            \draw[EdgeGraph, black](n)--(a);
            \draw[EdgeGraph, black](n)--(b);
        \end{tikzpicture}
        \enspace = \enspace 0 \text{ and}
    \end{equation}
    \begin{equation}
        \begin{tikzpicture}[Centering, scale=0.6]
            \node[NodeGraph, draw=Red, fill=Red!20](n)at(0,0){$ $};
            \node[NodeGraph, draw=Red, fill=Red!20](a)at(-1,1){$a$};
            \node[NodeGraph, draw=Green, fill=Green!20](b)at(-0.125,1.5){$b$};
            \node[NodeGraph, draw=Blue, fill=Blue!20](s)at(0.875,1.25){$\ast$};
            \node[NodeGraph, draw=Magenta, fill=Magenta!20](c)at(1.5,0.75){$c$};
            \draw[EdgeGraph, black](n)--(a);
            \draw[EdgeGraph, black](n)--(b);
            \draw[EdgeGraph, black](n)--(s);
            \draw[EdgeGraph, black](n)--(c);
        \end{tikzpicture}
        \enspace \comp \enspace
        \begin{tikzpicture}[Centering, scale=0.6]
            \node[NodeGraph, draw=Blue, fill=Blue!20](n)at(0,0){$ $};
            \node[NodeGraph, draw=Blue, fill=Blue!20](a)at(-0.5,1){$1$};
            \node[NodeGraph, draw=Magenta, fill=Magenta!20](b)at(0.5,1.25){$2$};
            \draw[EdgeGraph, black](n)--(a);
            \draw[EdgeGraph, black](n)--(b);
        \end{tikzpicture}
        \enspace = \enspace \left(
        \begin{tikzpicture}[Centering, scale=0.6]
            \node[NodeGraph, draw=Red, fill=Red!20](n)at(0,0){$ $};
            \node[NodeGraph, draw=Red, fill=Red!20](a)at(-1,1){$a$};
            \node[NodeGraph, draw=Green, fill=Green!20](b)at(-0.125,1.5){$b$};
            \node[NodeGraph, draw=Blue, fill=Blue!20](s)at(0.875,1.25){$\ast$};
            \node[NodeGraph, draw=Magenta, fill=Magenta!20](c)at(1.5,0.75){$c$};
            \node[NodeGraph, draw=Blue, fill=Blue!20](1)at(0.375,2.25){$1$};
            \node[NodeGraph, draw=Magenta, fill=Magenta!20](2)at(1.375,2.5){$2$};
            \draw[EdgeGraph, black](n)--(a);
            \draw[EdgeGraph, black](n)--(b);
            \draw[EdgeGraph, black](n)--(s);
            \draw[EdgeGraph, black](n)--(c);
            \draw[EdgeGraph, black](s)--(1);
            \draw[EdgeGraph, black](s)--(2);
        \end{tikzpicture} \right)
        \enspace = \enspace
            \begin{tikzpicture}[Centering, scale=0.6]
            \node[NodeGraph, draw=Red, fill=Red!20](n)at(0,0){$ $};
            \node[NodeGraph, draw=Red, fill=Red!20](a)at(-1,1){$a$};
            \node[NodeGraph, draw=Green, fill=Green!20](b)at(-0.25,1.5){$b$};
            \node[NodeGraph, draw=Blue, fill=Blue!20](1)at(0.5,1.25){$1$};
            \node[NodeGraph, draw=Magenta, fill=Magenta!20](2)at(1.25,0.875){$2$};
            \node[NodeGraph, draw=Magenta, fill=Magenta!20](c)at(1.5,0.25){$c$};
            \draw[EdgeGraph, black](n)--(a);
            \draw[EdgeGraph, black](n)--(b);
            \draw[EdgeGraph, black](n)--(c);
            \draw[EdgeGraph, black](n)--(1);
            \draw[EdgeGraph, black](n)--(2);
        \end{tikzpicture}
    \end{equation}
    A way to define {\em colored operads}~(see \cite{DY16} for more details on the theory of colored operads)
    is then to define them as any Hadamard product of a $C$-coloration operad with another operad. In particular,
    when doing this with a free operad, the elements of the product can be seen as Schröder trees with colored
    vertices and the partial composition as the grafting of trees if the root and the leaf on which it must grafted
    share the same color.
\end{exam}

Let us now define our last construction.

\begin{defi}
    Let $A$ be a set and $\spe{S}$ be a species. The set species of {\em functions from $A$ to $\spe{S}$} 
    is defined by $\func{A}{\spe{S}}[V] = \K\set{f:A\to \spe{S}[V]}$.
\end{defi}

The following proposition then tells us that if $\Operad$ has an operad structure, it naturally reflects 
on $\func{A}{\Operad}$.

\begin{prop}\label{prop_func}
    If $\Operad$ is an operad with unit $e$, $\func{A}{\Operad}$ has an operad structure with the elements 
    $e_v:A \to \set{e(v)}\in\func{A}{\Operad}[\set{v}]$ as units and partial composition defined by 
    $f_1\comp f_2(a) = f_1(a)\comp f_2(a)$.
\end{prop}

\begin{proof}
    We must verify that the diagrams~\eqref{def_op} are indeed commutative.
    \begin{itemize}
        \item Let be $f_1\in(\func{A}{\Operad})''[V_1]$, $f_2\in\func{A}{\Operad}[V_2]$ and 
        $f_3\in\func{A}{\Operad}[V_3]$. Then for all $a\in A$, we have
        \begin{equation}\begin{split}
            \big((f_1\comp[\ast_1] f_2)\comp[\ast_2] f_3\big)(a) &= (f_1\comp[\ast_1] f_2(a))\comp[\ast_2] f_3(a) \\
            &= \big(f_1(a)\comp[\ast_1] f_2(a)\big)\comp f_3(a) \\
            &= f_1(a)\comp[\ast_2] f_3(a)\comp f_2(a) \\
            &= (f_1\comp[\ast_2] f_3(a))\comp[\ast_1] f_2(a) \\
            &= \big((f_1\comp[\ast_2] f_3)\comp[\ast_1] f_2\big)(a),
        \end{split}\end{equation}
        where the third equality follows from the fact that $\Operad$ is an operad and the other equalities 
        from the definition of the partial composition on $\func{A}{\Operad}$. Hence, we have 
        $(f_1\comp[\ast_1]f_2)\comp[\ast_2] f_3 = (f_1\comp[\ast_2]f_3)\comp[\ast_1]f_2$.
        
        \item Let be $f_1\in(\func{A}{\Operad})'[V_1]$, $f_2\in(\func{A}{\Operad})'[V_2]$ and 
        $f_3\in\func{A}{\Operad}[V_3]$. Then for all $a\in A$:
        \begin{equation}\begin{split}
            \big((f_1\comp[\ast_1] f_2)\comp[\ast_2] f_3\big)(a) &= (f_1\comp[\ast_1] f_2(a))\comp[\ast_2] f_3(a) \\
            &= f_1(a)\comp[\ast_1] f_2(a)\comp f_3(a) \\
            &= f_1(a)\comp[\ast_2] f_3(a)\comp f_2(a) \\
            &= (f_1\comp[\ast_2] f_3(a))\comp[\ast_1] f_2(a) \\
            &= \big((f_1\comp[\ast_2] f_3)\comp[\ast_1] f_2\big)(a),
        \end{split}\end{equation}
        where the third equality comes from the fact that $\Operad$ is an operad. Hence, we have 
        $(f_1\comp[\ast_1]f_2)\comp[\ast_2] f_3 = (f_1\comp[\ast_2] f_3)\comp[\ast_1] f_2$.
        
        \item Let be $f\in(\func{A}{\Operad})'[V]$ and $v\not\in V$. Then for all $a\in A$, we have 
        the equalities $f\comp e_v(a)= f(a)\comp e(v) = \Operad[\sigma](f(a))$ and so 
        $f\comp e_v=\Operad[\sigma](f)$, where $\sigma$ is the bijection which sends $\ast$ to $v$ and is 
        the identity over $V$. If now $f\in\func{A}{\Operad}$, we have for all 
        $a\in A$: $e_{\ast}\comp f(a)= e(\ast)\comp f(a) = f(a)$ and so $e_{\ast}\comp f = f$. \qedhere
    \end{itemize}
\end{proof}

Note that if $A$ is a singleton then $\func{A}{\Operad} \cong \Operad$. Let $A,B,C,D$ be four sets such 
that $A$ and $B$ are disjoint and let $f:A\to C$ and $g:B\to D$ be two maps. We denote by $f\uplus g$ the map 
from $A\cup B$ to $C\cup D$ defined by $f\uplus g(a) = f(a)$ for $a\in A$ and $f\uplus g(b) = g(b)$ for 
$b\in B$. 

\begin{prop}
    Let $A$ and $B$ be two disjoint sets and $\Operad_1$ and $\Operad_2$ be two operads. Then the species 
    $\func{A,B}{\Operad_1,\Operad_2}$ defined by 
    $\func{A,B}{\Operad_1,\Operad_2}[V] = \set{f\uplus g\,|\, f\in\func{A}{\Operad_1}, g\in\func{B}{\Operad_2}}$ 
    is an operad with same partial composition as in Proposition~\ref{prop_func}.
\end{prop}

\begin{proof}
    We remark that since $A$ and $B$ are disjoint, 
    $f_1\uplus f_2\comp g_1\uplus g_2 = (f_1\comp g_1)\uplus (f_2\comp g_2)$. To conclude we apply what 
    was already shown in the proof of Proposition~\ref{prop_func}.
\end{proof}




\subsection{Application to multigraphs}\label{subsec_graph_op}

We now use the construction of the previous subsection to define operad structures on multigraphs. 

Recall from Remark~\ref{rem_bij_graphpol} that there is a monomorphism from $\augm{A}{\MG}$ to 
$\augm{A}{\Poly}$. We now make use of this monomorphism and consider the elements of $\augm{A}{\MG}$ as 
both multigraphs and polynomials. Let $p\in\Poly[V]$ be a sum of polynomials. Then for $A$ a set and $a\in A$, 
we denote by $p_a\in\augm{A}{\Poly}$ the sum of polynomials obtained by indexing all the variables in $p$ 
by $a$. Let now $p$ be any polynomial and $x_1,\dots, x_n$ a subset of its variables. Then for 
$q_1,\dots, q_n$ $n$ polynomials, we expand the notation introduced in equation \eqref{pol_composition} 
and denote by $p|_{\set{x_i\take q_i}}$ the polynomial 
$\big(\dots((p_1\comp[x_1] q_1)\comp[x_2] q_2)\dots\big)\comp[x_n] q_n$. This notation generalizes to sum 
of polynomials by recalling that the multiplication and addition of polynomials act as bilinear maps.

\begin{exam}
    For $A$ the singleton $\set{a}$ and $p$ the polynomial $xy\oplus x^2 + zy\in\Poly[\set{x,y,z}]$ we 
    have $p_a= x_ay_a\oplus x_a^2 + z_ay_a$. Let now $p=x\oplus yz$, $q_1 = u+v$ and
    $q_2=x_1\oplus x_2$. Then
    \begin{equation}\begin{split}
        p|_{x\take q_1, y\take q_2} &= q_1\oplus q_2z = (u+v)\oplus (x_1\oplus x_2)z \\
        &= u\oplus x_1z\oplus x_2z + v\oplus x_1z\oplus x_2z.
    \end{split}\end{equation}
\end{exam}

\begin{theorem}\label{th_op_ins}
    Let $A$ be a set and $\varphi$ be the morphism from 
    $\augm{A}{\MG}\cdot(\augm{A}{\MG}\times\func{A}{\K\Poly^1})$ to $\augm{A}{\MG}$ given by
    \begin{equation}
        \varphi(g_1\otimes g_2 \otimes f) = g_1\comp^f g_2 = g_1|_{\set{\ast_a\take f(a)_a}}\oplus g_2.
    \end{equation}
    Then $\varphi$ satisfies the hypotheses of Proposition~\ref{prop_semiprod} and we can consider the 
    semidirect product of $\augm{A}{\MG}$ and $\mathcal{F}_A^{\K \Poly^1}$ over $\varphi$.
\end{theorem}

\begin{proof}
    We need to check that $\varphi$ satisfies the three hypotheses of Proposition~\ref{prop_semiprod}. The 
    first two are simply computations over polynomials.
    \begin{description}
        \item[Commutativity.] Let $g_1$ be an element of $\augm{A}{\MG}''$ and $g_2\otimes f$ and 
        $g_3\otimes h$ two elements of $\augm{A}{\MG}\times\func{A}{\K\Poly^1}$. Then
        \begin{equation}\begin{split}
            (g_1\comp[\ast_1]^f g_2)\comp[\ast_2]^h g_3 
            &= (g_1|_{\set{\ast_{1a}\take f(a)_a}}\oplus g_2)|_{\{\ast_{2a}\take h(a)_a\}}\oplus g_3 \\
            &= g_1|_{\set{\ast_{1a}\take f(a)_a}}|_{\set{\ast_{2a}\take h(a)_a}}\oplus g_2\oplus g_3 \\
            &= g_1|_{\set{\ast_{2a}\take h(a)_a}}|_{\set{\ast_{1a}\take f(a)_a}}\oplus g_3\oplus g_2 \\
            &= (g_1|_{\set{\ast_{2a}\take h(a)_a}}\oplus g_3)|_{\set{\ast_{1a}\take f(a)_a}}\oplus g_2 \\
            &= (g_1\comp[\ast_2]^h g_3)\comp[\ast_1]^f g_2 .
        \end{split}\end{equation}
        \item[Associativity.] Let be $g_1$ an element of $\augm{A}{\MG}'$, $g_2\otimes f$ an element of 
        $(\augm{A}{\MG}\times\func{A}{\K\Poly^1})'$ and $g_3\otimes h$ an element of 
        $\augm{A}{\MG}\times\func{A}{\K\Poly^1}$. Then
        \begin{equation}\begin{split} 
            (g_1\comp[\ast_1]^f g_2)\comp[\ast_2]^h g_3
            &= (g_1|_{\set{\ast_{1a}\take f(a)_a}}\oplus g_2)|_{\set{\ast_{2a}\take h(a)_a}}\oplus g_3 \\
            &= g_1|_{\set{\ast_{1a}\take f(a)_a}}|_{\set{\ast_{2a}\take h(a)_a}}\oplus g_2|_{\set{\ast_{2a}\take h(a)_a}}\oplus g_3 \\
            &= g_1|_{\set{\ast_{1a}\take f(a)_a|_{\set{\ast_{2a}\take h(a)_a}}}}\oplus g_2|_{\set{\ast_{2a}\take h(a)_a}}\oplus g_3 \\
            &= g_1|_{\set{\ast_{1a}\take f\comp[\ast_2] h(a)_a}}\oplus g_2|_{\set{\ast_{2a}\take h(a)_a}}\oplus g_3 \\
            &= g_1\comp[\ast_1]^{f\comp[\ast_1] g}(g_2\comp[\ast_2]^h g_3).
        \end{split}\end{equation}
        \item[Unity.] Let $e: \spe{X}\to \augm{A}{\MG}$ be defined by $e(v) = \emptyset_{\set{v}}$ and 
        let $e_{\mathcal{F}}$ the unit of $\func{A}{\K\Poly^1}$. Let be $g\in\augm{A}{\MG}'[V]$ and 
        $v\not\in V$. Then
        \begin{equation}\begin{split}
            g\comp^{e_{\mathcal{F}}(v)}e(v) 
            &= g|_{\set{\ast_a \take e_{\mathcal{F}}(v)(a)_a}}\oplus\emptyset_{\set{v}} \\
            &= g|_{\set{\ast_a \take v_a}} = \augm{A}{\MG}[\sigma](g),
        \end{split}\end{equation}
        where $\sigma$ is the bijection which sends $\ast$ to $v$ and is the identity over $V$. If now 
        $g\in\augm{A}{\MG}[V]$, then for any $f$,
        \begin{equation}
            e(\ast)\comp^f g = \emptyset_{\set{\ast}}|_{\set{\ast_a \take f(a)_a}}\oplus g = g. \qedhere
        \end{equation}
    \end{description}
\end{proof}

\begin{exam}
    Let $A$ be the set of shapes $\set{\blacklozenge, \blacksquare, \blacktriangle}$ and let be 
    $g_1\otimes f_1 \in (\augm{A}{\MG}\ltimes_{\varphi}\func{A}{\K\Poly^1})'[\set{a,b}]$
    and $g_2\otimes f_2 \in \augm{A}{\MG}\ltimes_{\varphi}\func{A}{\K\Poly^1}[\set{c,d,e}]$ defined by:
    \begin{equation}
        \begin{tikzpicture}[Centering, scale=1]
            \node[NodeGraph] (a) at (0,0) {$a$};
            \node[NodeGraph] (s) at (1,0) {$\ast$};
            \node[NodeGraph] (b) at (1,1) {$b$};
            
            \draw[EdgeGraph, Square-Diamond] (a) -- (s);
            \draw[EdgeGraph, Triangle-Triangle] (b) to[bend left=45] (s); % Bend left
            \draw[EdgeGraph, Triangle-Square] (b) to[bend right=45] (s); % Bend right
        \end{tikzpicture}
        \otimes
        \begin{cases}
            \blacklozenge\mapsto a \oplus \ast\\
            \blacksquare\mapsto b\\
            \blacktriangle\mapsto b+\ast
        \end{cases}
        \quad\quad,\quad\quad
        \begin{tikzpicture}[Centering, scale=1]
            \node[NodeGraph] (e) at (0,0) {$e$};
            \node[NodeGraph] (d) at (1,0) {$d$};
            \node[NodeGraph] (c) at (0.5,0.866) {$c$};
            
            \draw[EdgeGraph, Square-Square] (e) -- (d);
            \draw[EdgeGraph, Triangle-Diamond] (d) -- (c);
            \draw[EdgeGraph, Triangle-Triangle] (c) -- (e);
        \end{tikzpicture}
        \otimes
        \begin{cases}
            \blacklozenge\mapsto c\oplus e + c\\
            \blacksquare\mapsto d\\
            \blacktriangle\mapsto c + d 
        \end{cases}
    \end{equation}
    The partial composition $(g_1\otimes f_1)\comp^{\varphi}(g_2\otimes f_2)$ is then:
    \begin{align}%
        \Bigg(
        \begin{tikzpicture}[Centering, scale=1]
            \node[NodeGraph] (e) at (0,0) {$e$};
            \node[NodeGraph] (d) at (1,0) {$d$};
            \node[NodeGraph] (c) at (0.5,0.866) {$c$};
            \node[NodeGraph] (a) at (-0.5,0.866) {$a$};
            \node[NodeGraph] (b) at (1.5,0.866) {$b$};
            \draw[EdgeGraph, Square-Square] (e) -- (d);
            \draw[EdgeGraph, Triangle-Diamond] (d) -- (c);
            \draw[EdgeGraph, Triangle-Triangle] (c) -- (e);
            \draw[EdgeGraph, Square-Diamond] (a) -- (c);
            \draw[EdgeGraph, Square-Diamond] (a) -- (e);
            \draw[EdgeGraph, Triangle-Triangle] (b) -- (c);
            \draw[EdgeGraph, Triangle-Square] (b) -- (d);
        \end{tikzpicture}
        +
        \begin{tikzpicture}[Centering, scale=1]
            \node[NodeGraph] (e) at (0,0) {$e$};
            \node[NodeGraph] (d) at (1,0) {$d$};
            \node[NodeGraph] (c) at (0.5,0.866) {$c$};
            \node[NodeGraph] (a) at (-0.5,0.866) {$a$};
            \node[NodeGraph] (b) at (1.5,0.866) {$b$};
            \draw[EdgeGraph, Square-Square] (e) -- (d);
            \draw[EdgeGraph, Triangle-Diamond] (d) -- (c);
            \draw[EdgeGraph, Triangle-Triangle] (c) -- (e);
            \draw[EdgeGraph, Square-Diamond] (a) -- (c);
            \draw[EdgeGraph, Square-Diamond] (a) -- (e);
            \draw[EdgeGraph, Triangle-Triangle] (b) to[bend right=35] (d);
            \draw[EdgeGraph, Triangle-Square] (b) to[bend left=35] (d);
        \end{tikzpicture}
        +&
        \begin{tikzpicture}[Centering, scale=1]
            \node[NodeGraph] (e) at (0,0) {$e$};
            \node[NodeGraph] (d) at (1,0) {$d$};
            \node[NodeGraph] (c) at (0.5,0.866) {$c$};
            \node[NodeGraph] (a) at (-0.5,0.866) {$a$};
            \node[NodeGraph] (b) at (1.5,0.866) {$b$};
            \draw[EdgeGraph, Square-Square] (e) -- (d);
            \draw[EdgeGraph, Triangle-Diamond] (d) -- (c);
            \draw[EdgeGraph, Triangle-Triangle] (c) -- (e);
            \draw[EdgeGraph, Square-Diamond] (a) -- (c);
            \draw[EdgeGraph, Triangle-Triangle] (b) -- (c);
            \draw[EdgeGraph, Triangle-Square] (b) -- (d);
        \end{tikzpicture}
        +
        \begin{tikzpicture}[Centering, scale=1]
            \node[NodeGraph] (e) at (0,0) {$e$};
            \node[NodeGraph] (d) at (1,0) {$d$};
            \node[NodeGraph] (c) at (0.5,0.866) {$c$};
            \node[NodeGraph] (a) at (-0.5,0.866) {$a$};
            \node[NodeGraph] (b) at (1.5,0.866) {$b$};
            \draw[EdgeGraph, Square-Square] (e) -- (d);
            \draw[EdgeGraph, Triangle-Diamond] (d) -- (c);
            \draw[EdgeGraph, Triangle-Triangle] (c) -- (e);
            \draw[EdgeGraph, Square-Diamond] (a) -- (c);
            \draw[EdgeGraph, Triangle-Triangle] (b) to[bend right=35] (d);
            \draw[EdgeGraph, Triangle-Square] (b) to[bend left=35] (d);
        \end{tikzpicture}
        \Bigg) \\ \nonumber
        &\otimes
        \begin{cases}
            \blacklozenge\mapsto a \oplus c\oplus e + a\oplus c\\
            \blacksquare\mapsto b\\
            \blacktriangle\mapsto b+c+d
        \end{cases}.
    \end{align}
    The sum $g_1\comp^{f_2} g_2$ is obtained removing $\ast$ from $g_1$ and reconnecting the loose ends by looking at the images of their 
    shapes by $f_2$. The map $f_1\comp f_2$ is obtained by composing the images of $f_1$ with the images of $f_2$ as in Proposition~\ref{prop_func}.
    Remark that the two first elements of the sum had one of their edge duplicated: this happens when reconnecting the end labeled by $\blacklozenge$
    to $c\oplus e$. 
\end{exam}

In all the following, we only consider this semi-direct product and we drop the $\varphi$ index.

\begin{defi}
    We call {\em graph insertion operad} any sub-operad of $\augm{A}{\MG}\ltimes\func{A}{\K\Poly^1}$,
    for $A$ a set and $\varphi$ as defined in Theorem~\ref{th_op_ins}.
\end{defi}

\begin{rema}
    This notion of graph insertion operad is different than the one mentioned in \cite{Kreimer:2000ja}, in 
    the context of Feynman graph insertions in quantum field theory.
\end{rema}


While defining a graph insertion operad seems involved, it essentially is equivalent to finding sub-operads of $\K\Poly^1$,
i.e. sums of degree 1 polynomials stable by composition.
Let us give two simple examples of graph insertion operads. Recall from 
Section~\ref{sec_intro} that we have a natural embedding of $\spe{\id}$ in $\K\Poly$, and three natural 
embeddings of $\Sym_+$ in $\K\Poly$.

\begin{exam}
    $\point{\G}$ has a natural operad structure given by 
    $\point{\G} \cong \G\times \spe{\id} \cong \augm{\set{0}}{\G}\ltimes\func{\set{0}}{\spe{\id}}$.
    For $(g_1,v_1)$ and $(g_2,v_2)$ two pointed graphs, the partial composition $(g_1,v_1)\comp(g_2,v_2)$ 
    is then equal to $(g_3,v_1|_{\ast\take v_2})$ where $g_3$ is the graph obtained by connecting all the 
    ends on $\ast$ to $v_2$. More formally,
    \begin{equation}\begin{split}
        (g_1,v_1)\comp (g_2,v_2) &= (g_1|_{\ast\take v_2}\oplus g_2, v_1|_{\ast\take v_2}) \\
        &= (\G[\sigma](g_1)\oplus g_2,v_1|_{\ast\take v_2}),
    \end{split}\end{equation}
    where $\sigma$ is the bijection which sends $\ast$ on $v_2$ and which is the identity on the rest of 
    its domain. For instance, we have:
    \begin{equation}
        \begin{tikzpicture}[Centering,scale=.7]
            \node[NodeGraph](s)at(0,0){$\ast$};
            \node[RootGraph](a)at(1,1){$a$};
            \node[NodeGraph](b)at(1,-1){$b$};
            \draw[EdgeGraph](a)--(s);
            \draw[EdgeGraph](s)--(b);
        \end{tikzpicture}
        \enspace \comp \enspace
        \begin{tikzpicture}[Centering,scale=1]
            \node[RootGraph](c)at(0,0){$c$};
            \node[NodeGraph](d)at(1,0){$d$};
            \draw[EdgeGraph](c)--(d);
        \end{tikzpicture}
        \enspace = \enspace
        \begin{tikzpicture}[Centering,scale=.7]
            \node[RootGraph](a)at(1,1){$a$};
            \node[NodeGraph](b)at(1,-1){$b$};
            \node[NodeGraph](c)at(0,0){$c$};
            \node[NodeGraph](d)at(2,0){$d$};
            \draw[EdgeGraph](a)--(c);
            \draw[EdgeGraph](c)--(b);
            \draw[EdgeGraph](c)--(d);
        \end{tikzpicture}.
    \end{equation}
    Remark that the operad $\NAP$~\cite{Liv06} is a sub-operad of the operad above and hence is a graph 
    insertion operad.
\end{exam}


\begin{exam}
    $\G$ has a natural operad structure given by 
    $\G\cong \G\times \Sym_+\cong \augm{\set{0}}{\G}\ltimes\func{\set{0}}{\Sym}$,
    where we consider here the embedding $V\mapsto \bigoplus_{v\in V} v$.
    For $g_1$ and $g_2$ two graphs, the partial composition $g_1\circ g_2$ is then the graph obtained by 
    adding an edge between each neighbour of $\ast$ and each vertex of $g_2$. More formally, for 
    $g_1\in G'[V_1]$ and $g_2\in G[V_2]$:
    \begin{equation}\begin{split}
        g_1\comp g_2 &= g_1|_{\ast \take \bigoplus V_2}\oplus g_2\\
        &= g_1|_{V_1}\oplus\bigoplus_{v\in n(\ast)} v\bigoplus_{v\in V_2}\oplus g_2 \\
        &= g_1|_{V_1}\oplus\bigoplus_{v_1\in n(\ast), v_2\in V_2} v_1v_2\oplus g_2,
    \end{split}\end{equation}
    where $n(\ast)$ is the set of neighbours of $\ast$. Each of terms of the last line have a 
    combinatorial interpretation: $g_1|_{V_1}$ is $g_1$ to which we removed $\ast$, the term 
    $\bigoplus_{v_1\in n(\ast),v_2\in V_2}v_1v_2$ means that we add an edge between any element in 
    $n(\ast)$ and any element in $V_2$, and finally the term $g_2$ means that we keep all the edges of $g_2$.
    For instance, we have:
    \begin{equation}
        \begin{tikzpicture}[Centering,scale=.7]
            \node[NodeGraph](s)at(0,0){$\ast$};
            \node[NodeGraph](a)at(1,1){$a$};
            \node[NodeGraph](b)at(1,-1){$b$};
            \draw[EdgeGraph](a)--(s);
            \draw[EdgeGraph](s)--(b);
        \end{tikzpicture}
        \enspace \comp \enspace
        \begin{tikzpicture}[Centering,scale=1]
            \node[NodeGraph](c)at(0,0){$c$};
            \node[NodeGraph](d)at(1,0){$d$};
            \draw[EdgeGraph](c)--(d);
        \end{tikzpicture}
        \enspace = \enspace
        \begin{tikzpicture}[Centering,scale=.7]
            \node[NodeGraph](a)at(1,1){$a$};
            \node[NodeGraph](b)at(1,-1){$b$};
            \node[NodeGraph](c)at(0,0){$c$};
            \node[NodeGraph](d)at(2,0){$d$};
            \draw[EdgeGraph](a)--(c);
            \draw[EdgeGraph](c)--(b);
            \draw[EdgeGraph](c)--(d);
            \draw[EdgeGraph](a)--(d);
            \draw[EdgeGraph](b)--(d);
        \end{tikzpicture}.
    \end{equation}
\end{exam}

Let us now prove that the two partial composition introduced in subsection~\ref{subsec_cano_op} do indeed 
make $\point{\MG}_{or}$ and $\MG$ operads.

\begin{proof}[Proof of Theorem~\ref{th_op_graph_or}]
    Recall from Example~\ref{ex_deform} that we have a bijection 
    $\MG_{or}\cong \augm{\set{\mathtt{s,t}}}{\MG}$. The isomorphism 
    $\point{\MG_{or}}\cong \augm{\set{\mathtt{s,t}}}{\MG}\times\spe{\id}\times\Sym_+ \cong \augm{\set{\mathtt{s,t}}}{\MG}\ltimes\func{\set{s},\set{t}}{\spe{\id},\Sym_+}$, 
    give the desired operad structure on $\point{\MG_{or}}$ when considering the embedding 
    $V\mapsto \sum_{v\in V}v$ of $\Sym_+$ in $\K\Poly^1$.
\end{proof}

\begin{proof}[Proof of Theorem~\ref{th_graphop_cano}]
    This is the operad structure on $\MG$ given by 
    $\MG\cong \MG\times\Sym_+\cong \augm{\set{0}}{\MG}\ltimes\func{\set{0}}{\Sym_+}$
    when considering the embedding $V\mapsto \sum_{v\in V}v$ of $\Sym_+$ in $\K\Poly^1$.
\end{proof}

We end this section by mentioning that while we restricted ourselves to multigraphs in this paper, all the 
work done in this section naturally generalizes to the very general framework of multi-hypergraphs, whose 
edges can contain any non-null number of vertices and can appear more than once. This is done by replacing 
$\MG$ by $\MHG$, the species of multi-hypergraphs, and $\K\Poly^1$ by $\K\Poly$ in Theorem~\ref{th_op_ins}.
In this more general context, we can also connect a loose end to a group of vertices, which also allows to 
increase the number of vertices of the edges. 

\appendix


\section{Koszul duality and Koszul operads}\label{ann_koszul}
As said at the end of Section~\ref{sec_intro}, 
two advantages of defining an operad by its generators and relations are
that it is possible to construct (under some conditions) its Koszul dual and 
that it is possible to check if the operad is Koszul.


\subsection*{Koszul duality}
Let us begin by defining the Koszul dual of an operad.
To do this, we consider from now on that we have an arbitrary order for every finite set $V$ in order to consider the signature of a bijection 
between two different sets. 

For $\spe{S}$ a linear species, we denote by $\spe{S}^*$ the {\em dual} species of $\spe{S}$ which is defined by $\spe{S}^*[V]=\spe{S}[V]^*$ and 
$\spe{S}^*[\sigma](f) = f\circ\spe{S}[\sigma^{-1}]$. We denote by $\spe{S}^\vee$ the species defined by $\spe{S}^\vee[V]=\spe{S}^*[V]$ and
$\spe{S}^\vee[\sigma](f) = \sign(\sigma)f\circ\spe{S}[\sigma^{-1}]$, where $\sign(\sigma)$ is the signature of $\sigma$. 

\begin{defi}
    Let $\Operad =\Ope(\mathcal{G},\mathcal{R})$ be a binary quadratic operad. Define the linear form $\scalar{ -}{ -}$ on 
    $\free{\mathcal{G}^{\vee}}^{(2)}\times\free{\mathcal{G}}^{(2)}$ by
    \begin{equation}
        \scalar{f_1\comp f_2}{x_1\comp x_2} = f_1(x_1)f_2(x_2),
    \end{equation}
    The {\em Koszul dual} of $\Operad$ is then the operad $\Operad^!=\Ope(\mathcal{G}^{\vee},\mathcal{R}^{\bot})$
    where $\mathcal{R}^{\bot}$ is the orthogonal of $\mathcal{R}$ for $\scalar{-}{-}$.
\end{defi}


\begin{exam}\label{ex_kdual}
    The {\em Lie operad} $\Lie$ is the quotient of the free operad over one antisymmetric generator by the Jacobi relation. More formally, denote by $[a,b]$
    the generating element of $\Sym_2^{\vee}[\set{a,b}]$ and by $[b,a]= (ab)\cdot [a,b]$ so that $[b,a] = -[a,b]$. Then $\Lie=\Ope(\Sym_2^\vee,\mathcal{R})$ with 
    $\mathcal{R}$ the sub-species of $\free{\Sym_2^\vee}^{(2)}$ generated by the {\em Jacobi relation} $[a,[b,c]]+[c,[a,b]]+[b,[c,a]$, where $[a,[b,c]]$ stands 
    for $[a,\ast]\comp[\ast] [b,c]$ (one can check that this is indeed stable under the action of $\symgrp_{\set{a,b,c}}$ and hence $\mathcal{R}$ is indeed a species). 
    This operad is the Koszul dual of $\Com$. Indeed, for every pair $\alpha < \beta$  in $\set{a,b,c,\ast}$ ordered with the alphabetical order ($c<\ast$) denote 
    by $s_{\alpha\beta}^{\vee}= [\alpha,\beta]$ the dual of $s_{\set{\alpha,\beta}}\in\ComMag[\set{\alpha, \beta}]$. Then we have, for example:
    \begin{align}%
        \scalar{&
            s^\vee_{a\ast}\comp s^\vee_{bc} + s^\vee_{c\ast}\comp s^\vee_{ab} + s^\vee_{b\ast}\comp s^\vee_{ca}
            }{
            s_{\set{a,\ast}}\comp s_{\set{b,c}} - s_{\set{c,\ast}}\comp s_{\set{a,b}}
            }\\ \nonumber
        &= \scalar{
            s^\vee_{a\ast}{\comp} s^\vee_{bc}
            }{
            s_{\set{a,\ast}}{\comp} s_{\set{b,c}}
            } 
        +  \scalar{
            s^\vee_{c\ast}{\comp} s^\vee_{ab}
            }{
            s_{\set{a,\ast}}{\comp} s_{\set{b,c}}
            } \\ \nonumber
        & \hspace{1.65in} + \scalar{
            s^\vee_{b\ast}{\comp} s^\vee_{ca}
            }{
            s_{\set{a,\ast}}{\comp} s_{\set{b,c}}
            } \\ \nonumber
        &- \scalar{
            s^\vee_{a\ast}{\comp} s^\vee_{bc}
            }{
            s_{\set{c,\ast}}{\comp} s_{\set{a,b}}
            } 
        - \scalar{
            s^\vee_{c\ast}{\comp} s^\vee_{ab}
            }{
            s_{\set{c,\ast}}{\comp} s_{\set{a,b}}
            } \\ \nonumber
        & \hspace{1.64in} - \scalar{
            s^\vee_{b\ast}{\comp} s^\vee_{ca}
            }{
            s_{\set{c,\ast}}{\comp} s_{\set{a,b}}
            } \\ \nonumber
        &= s^\vee_{a\ast}(s_{\set{a,\ast}})s^{\vee}_{bc}(s_{\set{b,c}}) 
        + s^\vee_{c\ast}(s_{\set{a,\ast}})s^{\vee}_{ab}(s_{\set{b,c}}) 
        + s^\vee_{b\ast}(s_{\set{a,\ast}})s^{\vee}_{ca}(s_{\set{b,c}}) \\ \nonumber
        &- s^\vee_{a\ast}(s_{\set{c,\ast}})s^{\vee}_{bc}(s_{\set{a,b}}) 
        - s^\vee_{c\ast}(s_{\set{c,\ast}})s^{\vee}_{ab}(s_{\set{a,b}}) 
        - s^\vee_{b\ast}(s_{\set{c,\ast}})s^{\vee}_{ca}(s_{\set{a,b}}) \\ \nonumber
        &= 1 +0+0-1-0-0 = 0.
    \end{align}
\end{exam}




\subsection*{Koszul operads}\label{sec_koszul}
Koszulity is an important aspect of operad theory. We only give here a very quick overview of Koszulity and Gröbner bases for operads which hides a 
lot of the theory. We do not give the general results but only restricted versions which suffice for our use.
We refer the reader to the literature; for a broader approach of the topic, see for example \cite{LV12, Men15, Hof10, DK10}.
In particular, all the examples presented here come from \cite{DK10}.

In order to give the characterisation which interests us, we need to introduce the concepts of $\mathscr{L}$-species, shuffle operads and Gröbner bases. 
Informally, we can see these objects as the same as species and operads, except with a total order on every set of vertices.


\noindent\textbf{$\mathscr{L}$-species.} A {\em linear positive $\mathscr{L}$-species}  consists of the following data:
\begin{itemize}
   \item for each finite set $V$ and total order $l$ on $V$, a 
   vector space $\spe{S}[V,l]$, such that $\spe{S}[\emptyset,\emptyset]=\set{0}$.
   \item For each increasing bijection $\sigma: (V,l)\to (V',l')$, a linear map $\spe{S}[\sigma]:\spe{S}[V,l]\to \spe{S}[V',l']$.
   These maps should be such that $\spe{S}[\sigma_1\circ\sigma_2] = \spe{S}[\sigma_1]\circ \spe{S}[\sigma_2]$ and $\spe{S}[\id] = \id$.
\end{itemize}


In the sequel, we write order to designate a total order and $\mathscr{L}$-species to designate linear positive $\mathscr{L}$-species. For $\spe{S}$ an
$\mathscr{L}$-species and $l$ an order on $V$, we also denote by $\spe{S}[l]=\spe{S}[V,l]$. We can do this since the data of $V$ is included in $l$.
As for species, we denote by $\spe{X}$ the $\mathscr{L}$-species defined $\spe{X}[V,l]=\set{0}$ if $V$ is not a singleton and $\spe{X}[\set{v},v]= \K{v}$ else.


As in the case of classical species, $\mathscr{L}$-species also have constructions on them, but before giving them, let us give some notations.
\begin{itemize}
    \item For $l$ an order on $V$ and $W\subseteq V$ a subset of $V$, we denote by $l_W$ the order on $W$ induced by $l$.
    \item For $l=l_1\dots l_n$ an order on a set $V$ of size $n$, $i\in[n]$ and $\ast\not\in V$, we denote by $l\stackrel{i}{\take}\ast$ the order 
    $l_1\dots l_{i-1}\ast l_i\dots l_n$ on $V\cup\set{\ast}$.
    \item For $l^1,\dots l^k$ $k$ orders on pairwise disjoint sets $V_1,\dots, V_k$, we denote by $sh(l^1,\dots,l^k)=\set{w\,|\, w_{V_i}=l^i}$ the set of 
    {\em shuffles} of $l^1,\dots, l^k$. Note that this is an ``associative operation'' in the sense that for $l^1,l^2,l^3$ three orders, the union of
    the shuffles of $l^1$ with the elements of $sh(l^2,l^3)$ is exactly $sh(l^1,l^2,l^3)$.
    \item  The {\em shuffle compositions} $\Comp_{sh}[V,l]$ of an ordered set $(V,l)$ are the compositions $P=P_1,\dots,P_k$ of $V$ such that, for every 
    $1\leq i<j\leq k$, $\min_l P_i<\min_l P_j$.
\end{itemize}

Let $\spe{R}$ and $\spe{S}$ be two linear $\mathscr{L}$-species and $l$ a total order on $V$. Denote by $n=\card{V}$ and let be $i\in[n]$. We define the following operations.

\begin{alignat*}{2}
&Product\quad & \spe{R}\cdot\spe{S}[V,l] = & \ \bigoplus_{l\in sh(l',l'')} \spe{R}[l']\otimes\spe{S}[l''], \\
&i\text{-}th\ derivative\quad & \spe{S}^i[V,l] = & \ \spe{S}[V\cup\set{\ast}, l\stackrel{i}{\take}\ast], \\
&Composition\quad & \spe{R}(\spe{S})[V,l] = & \ \bigoplus_{P\in\Comp_{sh}[V,l]} \bigg(\spe{R}[\set{P_1,\dots,P_k},P]\otimes 
\bigotimes_{i=1}^k \spe{S}[P_i, l_{P_i}]\bigg).
\end{alignat*}
%\spe{S}[P_1,l_{P_1}]\otimes\dots\otimes\spe{S}[P_k,l_{P_k}].$$

Since we have the $\mathscr{L}$-species $\spe{X}$ and the notion of composition, we can define Schröder tree on $\mathscr{L}$-species in the same way as
for species: if $\spe{S} = \spe{X} + \spe{S}_{2^+}$, then $\stree{\spe{S}} = \spe{X} + \spe{S}_{2^+}(\stree{\spe{S}})$. 
In this case, for $\spe{S}$ a $\mathscr{L}$-species, $\stree{\spe{S}}[V]$ is the vector space of rooted {\em planar} trees with internal vertices decorated 
with elements of $\spe{S}$ and set of leaves $V$.
This is because of the orders: instead of indexing with partitions in the composition we use shuffle compositions.

For $\spe{S}$ a $\mathscr{L}$-species and $l$ an order on $V$, the fact that we have an order enables us an easier notation of the elements of $\stree{\spe{S}}[V,l]$
as operations, e.g. $\alpha(l_1,\dots, l_n)$.

\begin{exam}\label{ex_free_shuffle}
    Let $\spe{S}$ be the $\mathscr{L}$-species defined by $\spe{S}[V,l]=\set{0}$ when $\card{V}\not=2$ and $\spe{S}[V,l]=\K\set{1,\dots n}$ otherwise, for $n$
    an integer greater than 1. Then the generators of $\stree{\spe{S}}^{sh}[[3],123]$ are of the form
    \begin{equation}
        \begin{tikzpicture}[Centering,xscale=0.6,yscale=0.4]
            \node[NodeFree](x)at(0,0){$x$};
            \node[NodeFree](y)at(-1,2){$y$};
            \node[NodeFree, draw=white](1)at(-2,4){$1$};
            \node[NodeFree, draw=white](2)at(0,4){$2$};
            \node[NodeFree, draw=white](3)at(1,2){$3$};
            \draw[EdgeFree](x)--(y);
            \draw[EdgeFree](x)--(3);
            \draw[EdgeFree](1)--(y);
            \draw[EdgeFree](2)--(y);
        \end{tikzpicture}
        \enspace , \enspace
        \begin{tikzpicture}[Centering,xscale=0.6,yscale=0.4]
            \node[NodeFree](x)at(0,0){$x$};
            \node[NodeFree](y)at(-1,2){$y$};
            \node[NodeFree, draw=white](1)at(-2,4){$1$};
            \node[NodeFree, draw=white](3)at(0,4){$3$};
            \node[NodeFree, draw=white](2)at(1,2){$2$};
            \draw[EdgeFree](x)--(y);
            \draw[EdgeFree](x)--(2);
            \draw[EdgeFree](1)--(y);
            \draw[EdgeFree](3)--(y);
        \end{tikzpicture}
        \enspace\text{and}\enspace
        \begin{tikzpicture}[Centering,xscale=0.6,yscale=0.4]
            \node[NodeFree](x)at(0,0){$x$};
            \node[NodeFree](y)at(1,2){$y$};
            \node[NodeFree, draw=white](1)at(-1,2){$1$};
            \node[NodeFree, draw=white](2)at(0,4){$2$};
            \node[NodeFree, draw=white](3)at(2,4){$3$};
            \draw[EdgeFree](x)--(y);
            \draw[EdgeFree](x)--(1);
            \draw[EdgeFree](3)--(y);
            \draw[EdgeFree](2)--(y);
        \end{tikzpicture},
    \end{equation}
    where $1\leq x,y\leq n$. This is because the only shuffle compositions of size 2 of $[3]$ are $(12,3)$, $(13,2)$ and $(1,23)$. 
    These trees can be denoted in a more compact way by $\mu_x(\mu_y(1,2),3)$, $\mu_x(\mu_y(1,3),2)$ and $\mu_x(1,\mu_y(2,3))$.
\end{exam}

Let us end this brief presentation of $\mathscr{L}$-species by giving their link to classical species. Let $\mathscr{F}$ be the forgetful functor which send a  
species $\spe{S}$ on the $\mathscr{L}$-species $\spe{S}^{\mathscr{F}}$ defined by:
\begin{itemize}
    \item for $l$ an order on $V$, $\spe{S}^{\mathscr{F}}[V,l]=\spe{S}[V]$. 
    \item For $\sigma:(V,l)\to (V',l')$ an increasing bijection, $\spe{S}^{\mathscr{F}}[\sigma]$ is given by
    \begin{equation}
        \spe{S}^{\mathscr{F}}[V,l]=\spe{S}[V] \stackrel{\sigma}{\to} \spe{S}[W]= \spe{S}^{\mathscr{F}}[W,l'].
    \end{equation}
    \item For $f:\spe{S}\to\spe{R}$ a species morphism, $f^{\mathscr{F}}$ is the $\mathscr{L}$-species morphism given by
        \begin{equation}
        f^{\mathscr{F}}_{V,l}:\spe{S}^{\mathscr{F}}[V,l]= \spe{S}[V] \stackrel{f}{\to} \spe{R}[V]= \spe{R}^{\mathscr{F}}[V,l].
    \end{equation}
\end{itemize}
This is a forgetful functor in the sense that we forget the action $\sigma_V$ on $\spe{S}[V]$. We have the following fundamental proposition from
\cite{DK10}.

\begin{prop}[{\cite[Proposition 3]{DK10}}]\label{prop_dk10}
    Let $\spe{S}$ and $\spe{R}$ be two species. Then 
    \begin{equation}
        (\spe{R}(\spe{S}))^{\mathscr{F}}=\spe{R}^\mathscr{F}(\spe{S}^\mathscr{F}).
    \end{equation}
\end{prop}


\noindent\textbf{Shuffle operads.} We would like to define shuffle operads as $\mathscr{L}$-species satisfying the same axioms that a linear species must satisfy to be an operad. For now we can
not do this because we do not have the notion of derivative of a shuffle operad $\Operad'$. Fortunately, there is a way to make sense of the different 
species appearing in the diagrams \eqref{def_op}. 

For $l$ an order on $V$ and $v\in V$, denote by $\text{arg}_lv$ the index of $v$ in $l$: $l_{\text{arg}_lv}=v$. 
For $\spe{R}$ and $\spe{S}$ two $\mathscr{L}$-species we define the $\mathscr{L}$-species $\spe{R}'\cdot\spe{S}$ and $\spe{R}''\cdot\spe{S}\cdot\spe{S}$ as follow:
\begin{equation}\begin{split}
    \spe{R}'\cdot\spe{S}[V,l] &= \bigotimes_{l\in sh(l^1,l^2)} \spe{R}^{\text{arg}_l l_1^2}[l^1]\otimes\spe{S}[l^2], \\
    \spe{R}''\cdot\spe{S}^2[V,l] &= \bigotimes_{l\in sh(l^1,l^2,l^3)} \spe{R}^{\text{arg}_l l^2_1, \text{arg}_l l^3_1}[l^1]\otimes\spe{S}[l^2]\otimes\spe{S}[l^3].
\end{split}\end{equation}
%Note that the notation $\spe{R}'\cdot\spe{R}'\cdot\spe{S}$ is non ambiguous, that is to say, $\spe{R}'\cdot(\spe{R}'\cdot\spe{S})=(\spe{R}'\cdot\spe{R}')\cdot\spe{S}$.

A {\em shuffle operad} is then a $\mathscr{L}$-species $\Operad$ with a unity $e$ and a partial composition $\comp^{sh}:\Operad'\cdot\Operad\to\Operad$ 
such that the diagrams \eqref{def_op} commutes. 
For $\spe{S}$ a $\mathscr{L}$-species, the {\em free shuffle operad over $\spe{S}$} is denoted by $\free{\spe{S}}^{sh}$ and defined in the same way as the free
operad over a species. The same goes with the {\em ideal of a shuffle operad} and the notation $\Ope_{sh}(\mathcal{G},\mathcal{R})$.

\begin{rema}
    With the notations of elements of the free operad as operations, the partial composition of the free operad is then the composition of operation:
    $\mu_x(\dots, \ast,\dots)\comp^\xi \mu_y(\dots)= \mu_x(\dots,\mu_y(\dots),\dots)$.
\end{rema}

We have the following corollary from Proposition~\ref{prop_dk10}.
\begin{theorem}[{\cite[Corollary 1]{DK10}}]
    Let $\mathcal{G}$ be a species. The image by $\mathscr{F}$ of the free operad generated by $\mathcal{G}$ is isomorphic to the free shuffle operad generated
    by $\mathscr{F}(\mathcal{G})$. For $\mathcal{R}$ a sub-species of $\free{\mathcal{G}}$, the image by $\mathscr{F}$ of the operad ideal generated by $\mathcal{R}$ is 
    isomorphic to the shuffle operad ideal generated  by $\mathscr{F}(\mathcal{R})$.
    This writes as $\free{\mathcal{G}}^{\mathscr{F}} \cong \free{\mathcal{G}^{\mathscr{F}}}^{sh}$ and $(\mathcal{R})^{\mathscr{F}}=(\mathcal{R}^{\mathscr{F}})$. 
    Hence $\Ope(\mathcal{G},\mathcal{R})^{\mathscr{F}}\cong\Ope_{sh}(\mathcal{G}^{\mathscr{F}},\mathcal{R}^{\mathscr{F}})$.
\end{theorem}

\noindent\textbf{Admissible order.} Let $\spe{S}$ be a $\mathscr{L}$-species. In order to define the notion of Gröbner bases, we need to introduce an order on the trees generating $\free{\spe{S}}^{sh}[V,l]$.
Instead of giving the broader notion of admissible order defined in \cite{DK10}, we only give a small variation of the {\em path-lexicographic ordering}.

First for every order $l$ on $V$, fix a basis of $\spe{S}[l]$ and an order on this basis such that for every $x$ in the chosen basis of $\spe{S}[l]$ and $\sigma:l\to l'$,  $\sigma\cdot x$ is also in the chosen basis of $\spe{S}[l']$ and for $y$ an other element of the basis greater than $x$ we have 
$\sigma\cdot x<\sigma\cdot y$.
That is to say, the order does not depend on the labels (but it can depend on their relative order). Given an ordered basis, we also have an order on the words on elements of the basis given by the lexicographic order.
Let now be $t\in\free{\spe{S}}^{sh}[l_1\dots l_n]$ such that every internal node is labelled by an element of the chosen bases. For all $i\in[n]$, there is a unique path from $l_i$ to the root of $t$. Denote by $a_i$ the word composed, from left
to right, of the labels of the nodes of this path, from the root to the leaf. We associate to $t$ the sequence $(a_1,\dots, a_i, w)$, where $w$ is the
word obtained by reading the leaves of $t$ from left to right. 

For two trees $t,t'\in\free{\spe{S}}^{sh}$, with associated sequences $(a_1,\dots, a_n, w)$ and $(b_1\dots, b_n, w')$ we then compare $t$ and $t'$ by lexicographically comparing $a_1$ with $b_1$ then $a_2$ with $b_2$, etc, and reverse lexicographically comparing $w$ with $w'$ if $a_i=b_i$ for all $i$.

\begin{exam}
    The $\mathscr{L}$-species $\spe{S}$ of Example~\ref{ex_free_shuffle} have natural ordered bases equal to the set $\set{1,\dots,n}$ with the natural order.
    The sequences attached to the given trees are then respectively $(xy,xy,123)$, $(xy,x,xy,132)$ and $(x,xy,xy,123)$ and we have
    \begin{equation}
        \begin{tikzpicture}[Centering,xscale=0.6,yscale=0.4]
            \node[NodeFree](x)at(0,0){$x$};
            \node[NodeFree](y)at(-1,2){$y$};
            \node[NodeFree, draw=white](1)at(-2,4){$1$};
            \node[NodeFree, draw=white](2)at(0,4){$2$};
            \node[NodeFree, draw=white](3)at(1,2){$3$};
            \draw[EdgeFree](x)--(y);
            \draw[EdgeFree](x)--(3);
            \draw[EdgeFree](1)--(y);
            \draw[EdgeFree](2)--(y);
        \end{tikzpicture}
        \enspace > \enspace
        \begin{tikzpicture}[Centering,xscale=0.6,yscale=0.4]
            \node[NodeFree](x)at(0,0){$x'$};
            \node[NodeFree](y)at(-1,2){$y'$};
            \node[NodeFree, draw=white](1)at(-2,4){$1$};
            \node[NodeFree, draw=white](3)at(0,4){$3$};
            \node[NodeFree, draw=white](2)at(1,2){$2$};
            \draw[EdgeFree](x)--(y);
            \draw[EdgeFree](x)--(2);
            \draw[EdgeFree](1)--(y);
            \draw[EdgeFree](3)--(y);
        \end{tikzpicture}
    \end{equation}
    if $x>x'$ or $x=x'$ and $y>y'$ or $x=x'$ and $y=y'$. 
\end{exam}

Now remark that trees in $\free{\spe{S}}^{sh}[V,l]$ with basis elements as internal node labels make a basis $\free{\spe{S}}^{sh}[V,l]$. Hence any $x\in\free{\spe{S}}^{sh}$ can be written as a sum of such elements and we define $\text{lt}(x)$ the {\em leading term} of $x$ as the maximal element in 
this sum.

\noindent\textbf{Divisibility and S-polynomials.} Let $\spe{S}$ be a $\mathscr{L}$-species. A tree $t$ of $\free{\spe{S}}^{sh}$ is {\em divisible} by another tree $t'$ of $\free{\spe{S}}^{sh}$ if $t'$
is a sub-tree of $t$. Here a sub-tree must also conserve the order of the leaves.
A tree $u$ of $\free{\spe{S}}^{sh}$ is a {\em small common multiple} of two tree $t$ and $t'$ if it is divisible by both $t$ and $t'$ and its number of vertices is 
less than the total number of vertices of $t$ and $t'$.

\begin{exam}
    Let $\spe{S}$ be an $\mathscr{L}$-species, let $l=l_1,\dots l_n$ be an order and let $\alpha(\beta(l_1,l_3),\gamma(\beta(l_2,l_6),l_4,l_5))$ be an element of $\free{\spe{S}}^{sh}$.
    This tree has among its divisors $\alpha(\beta(l_1,l_3),l_2)$ and $\gamma(\beta(l_1,l_4),l_2,l_3)$ but not $\gamma(\beta(l_1,l_3),l_2,l_4)$.
\end{exam}

If $t$ is divisible by $t'$, then there exists trees $\alpha$ and $\beta_1,\dots,\beta_k$ such that $t=\alpha(\dots, t'(\beta_1,\dots,\beta_k), \dots)$. We denote
by $m_{t,t'}$ the operation on any tree with same number of leaves than $t'$ which associate to a tree $u$ the tree $\alpha(\dots, u(\beta_1,\dots,\beta_k),\dots)$.
Let now $V$ be a finite set, $l$ an order on $V$ and $x,y\in\free{\spe{S}}^{sh}[V,l]$. Assume $\text{lt}(x)$ and $\text{lt}(y)$ have a small 
common multiple $u$. Then we have $m_{u,\text{lt}(x)}(\text{lt}(x)) = u = m_{u,\text{lt}(y)}(\text{lt}(y))$. We call {\em S-polynomial of $x$ and $y$ 
(corresponding to $u$)} the element
\begin{equation}
    s_u(x,y)= m_{u,\text{lt}(x)}(x) -\frac{c_x}{c_y}m_{u,\text{lt}(y)}(y),
\end{equation}
where $c_x$ and $c_y$ are the respective coefficient of the leading terms of $x$ and $y$.

\noindent\textbf{Gröbner bases and Koszulity.} We can finally give the definition of a Gröbner bases and a Koszul operad.

\begin{defi}[{\cite[Definition 13]{DK10}}]
    Let $\mathcal{G}$ be a $\mathscr{L}$-species and $\mathcal{R}$ be a $\mathscr{L}$-sub-species of $\free{\mathcal{G}}^{sh}$. Let $\mathcal{B}$ 
    be basis of $\mathcal{R}$.
    We say that $\mathcal{B}$ is a {\em Gröbner bases} of $\mathcal{R}$ if for every $x\in(\mathcal{R})$, the leading term of $x$ is divisible by 
    the leading term of one element in $\mathcal{B}$.
\end{defi}

\begin{defi}[{\cite[Corollary 3]{DK10}}]
    Let $\mathcal{G}$ be a $\mathscr{L}$-species and $\mathcal{R}$ be a quadratic $\mathscr{L}$-sub-species of $\free{\mathcal{G}}^{sh}$. We say that
    $\Ope_{sh}(\mathcal{G},\mathcal{R})$ is {\em Koszul} if $\mathcal{R}$ admits a Gröbner bases.

    Let $\mathcal{G}$ be a set species and $\mathcal{R}$ be a quadratic sub-species of $\free{\mathcal{G}}^{sh}$. We say that
    $\Ope(\mathcal{G},\mathcal{R})$ is {\em Koszul} if $\Ope_{sh}(\mathcal{G}^{\mathscr{F}},\mathcal{R}^{\mathscr{F}})$ is Koszul.
\end{defi}

When $\Operad$ is a Koszul symmetric operad, it admits a Koszul dual $\Operad^!$. In this case the Hilbert series of $\Operad$ and $\Operad^!$ are related by
the identity:
\begin{equation} \label{hdual}
    \hilbert{\Operad}(-\hilbert{\Operad^!}(-t)) = t.
\end{equation}

Let us finish by a characterisation of Gröbner bases.

\begin{prop}[{\cite[Theorem 1]{DK10}}]\label{prop_spol}
    Let $\mathcal{G}$ be a $\mathscr{L}$-species and $\mathcal{R}$ be a $\mathscr{L}$-sub-species of $\free{\mathcal{G}}^{sh}$. Let $\mathcal{B}$ 
    be basis of $\mathcal{R}$. Then $\mathcal{B}$ is a Gröbner bases if and only if for all pair of elements in $\mathcal{B}$, their S-polynomials
    are congruent to zero modulo $\mathcal{B}$ (i.e. they are in $(\mathcal{R})$).
\end{prop}

